% !TEX program = lualatex
% ---------------------------------------------------------------------------------------------
%  Readback note 43: The energy of amplitude sums: Schur's test on thin classes, block bounds, incoherent and degenerate pairs
%
%  Build:  python3 tools/rb.py note <slug>
%
%  One unit of a Lean artifact, rewritten as ordinary mathematics. The prose is hand-written. Lean names,
%  line ranges and counts come from data.tex; quoted Lean statements are copied from the artifact by the build.
%
%  Rules of the text: standard terminology; every symbol is explained where it first appears and listed again
%  in the notation at the end; no letter has two meanings; a proof is complete or says what it takes for granted;
%  Lean names appear only in the LEAN lines and in the concordance.
%
%  The signature rule (the build enforces it): a function, relation or predicate that the artifact defines is
%  written only through its notation macro, which the build makes from the spec (`notation`) with one slot for
%  each explicit argument of the Lean definition. No argument is fixed and dropped, also not "for brevity" and
%  not for one section. Every argument and every hypothesis of a stated declaration stands inside its
%  statement, or in a Setting that the statement cites by number; only implicit and instance arguments (the
%  ambient objects) may stand in the conventions.
% ---------------------------------------------------------------------------------------------
\documentclass[11pt,a4paper]{article}
\ifdefined\pdfvariable\pdfvariable suppressoptionalinfo \numexpr 1+2+4+8\relax\fi
\usepackage[a4paper,hmargin=35mm,top=22mm,bottom=27mm,footskip=12mm,nohead]{geometry}

\usepackage{fontspec}
\directlua{
  luaotfload.add_fallback("gfallback", {"STIXTwoMath-Regular.otf:mode=node;", "DejaVuSans.ttf:mode=node;"})
  luaotfload.add_fallback("gmonofallback", {"DejaVuSansMono.ttf:mode=node;", "STIXTwoMath-Regular.otf:mode=node;"})
}
\setmainfont{NewCM10}[Extension=.otf,UprightFont=*-Book,ItalicFont=*-BookItalic,
  BoldFont=*-Bold,BoldItalicFont=*-BoldItalic,SmallCapsFont=*-Book,SmallCapsFeatures={Letters=SmallCaps},
  RawFeature={fallback=gfallback}]
\setsansfont{NewCMSans10}[Extension=.otf,UprightFont=*-Book,ItalicFont=*-BookOblique,
  BoldFont=*-Bold,BoldItalicFont=*-BoldOblique,RawFeature={fallback=gfallback}]
\setmonofont{JetBrainsMono}[Extension=.otf,UprightFont=*-Regular,ItalicFont=*-Italic,
  BoldFont=*-SemiBold,BoldItalicFont=*-SemiBoldItalic,Scale=MatchLowercase,
  RawFeature={-calt;-liga;fallback=gmonofallback}]
\usepackage{amsmath}
\usepackage{amsthm}
\usepackage[math-style=TeX]{unicode-math}
\setmathfont{NewCMMath-Book.otf}
\emergencystretch=2.5em\relax

\usepackage{xcolor}
\usepackage{tikz}
\usetikzlibrary{calc,arrows.meta,positioning,backgrounds}
\usepackage[protrusion=true,expansion=true]{microtype}
\usepackage{enumitem}
\usepackage{needspace}
\usepackage{booktabs}
\usepackage{array}
\usepackage{longtable}
\usepackage{fancyhdr}
\usepackage{lastpage}
\usepackage{multicol}
\usepackage[hidelinks,pdfusetitle]{hyperref}

\definecolor{ink}{HTML}{0F172A}
\definecolor{inksoft}{HTML}{475569}
\definecolor{inkfaint}{HTML}{94A3B8}
\definecolor{rulegrey}{HTML}{CBD5E1}
\definecolor{arith}{HTML}{C2410C}
\definecolor{est}{HTML}{7C3AED}
\definecolor{ana}{HTML}{0891B2}
\definecolor{ok}{HTML}{047857}
\definecolor{leancol}{HTML}{3730A3}

\title{The energy of amplitude sums: Schur's test on thin classes, block bounds, incoherent and degenerate pairs}
\author{grh-compact}
\date{}

% ------------------------------------------------------------------------------------- macros
\newcommand{\leanfont}{\ttfamily\fontsize{8.6}{10.5}\selectfont\color{leancol}}
\newcommand{\lean}[1]{\csname LN@#1\endcsname}
\newcommand{\fillin}[1]{{\color{arith}\itshape[#1]}}  % a place of the template that is not written yet; the build reports every one that is left
\newcommand{\leanquote}[1]{\csname LQ@#1\endcsname}
% generated by rb.py note; do not edit
\newcommand{\ConcDecls}{31}
\newcommand{\Pin}{72fead10}
\newcommand{\PinText}{that of commit {\leanfont 72fead10}}
\newcommand{\ArtifactName}{grh-compact}
\expandafter\newcommand\csname LN@sumlatticeDiscinvle\endcsname{\href{https://github.com/lsernau/grh-lean/blob/72fead102f663d5c7a407964a13dc0cf9d977207/grh-compact/src/New/Filling/BlockBound.lean\#L21-L75}{{\leanfont Filling.\allowbreak{}sum\_latticeDisc\_inv\_le}}}
\expandafter\newcommand\csname LN@normsharedvectorlt\endcsname{\href{https://github.com/lsernau/grh-lean/blob/72fead102f663d5c7a407964a13dc0cf9d977207/grh-compact/src/New/Filling/BlockBound.lean\#L93-L125}{{\leanfont Filling.\allowbreak{}HasLevel.\allowbreak{}norm\_shared\_vector\_lt}}}
\expandafter\newcommand\csname LN@tsuminvPowWeightleofcodimle\endcsname{\href{https://github.com/lsernau/grh-lean/blob/72fead102f663d5c7a407964a13dc0cf9d977207/grh-compact/src/New/Filling/BlockBound.lean\#L130-L215}{{\leanfont Filling.\allowbreak{}tsum\_invPowWeight\_le\_of\_codim\_le}}}
\expandafter\newcommand\csname LN@submoduleIndexgIdelesmul\endcsname{\href{https://github.com/lsernau/grh-lean/blob/72fead102f663d5c7a407964a13dc0cf9d977207/grh-compact/src/New/Filling/BlockBound.lean\#L222-L226}{{\leanfont Filling.\allowbreak{}submoduleIndex\_gIdele\_smul}}}
\expandafter\newcommand\csname LN@sumpairWeightle\endcsname{\href{https://github.com/lsernau/grh-lean/blob/72fead102f663d5c7a407964a13dc0cf9d977207/grh-compact/src/New/Filling/BlockBound.lean\#L232-L297}{{\leanfont Filling.\allowbreak{}sum\_pairWeight\_le}}}
\expandafter\newcommand\csname LN@enormamplitudeSumleclassSizetop\endcsname{\href{https://github.com/lsernau/grh-lean/blob/72fead102f663d5c7a407964a13dc0cf9d977207/grh-compact/src/New/Filling/BlockBound.lean\#L304-L363}{{\leanfont Filling.\allowbreak{}enorm\_amplitudeSum\_le\_classSize\_top}}}
\expandafter\newcommand\csname LN@enormamplitudeSumleclassSize\endcsname{\href{https://github.com/lsernau/grh-lean/blob/72fead102f663d5c7a407964a13dc0cf9d977207/grh-compact/src/New/Filling/BlockBound.lean\#L368-L394}{{\leanfont Filling.\allowbreak{}enorm\_amplitudeSum\_le\_classSize}}}
\expandafter\newcommand\csname LN@amplitudeSumeqzeroofnotcharsAgree\endcsname{\href{https://github.com/lsernau/grh-lean/blob/72fead102f663d5c7a407964a13dc0cf9d977207/grh-compact/src/New/Filling/CrudeEnergy.lean\#L24-L41}{{\leanfont Filling.\allowbreak{}amplitudeSum\_eq\_zero\_of\_not\_charsAgree}}}
\expandafter\newcommand\csname LN@hasCommonFlagVectorofle\endcsname{\href{https://github.com/lsernau/grh-lean/blob/72fead102f663d5c7a407964a13dc0cf9d977207/grh-compact/src/New/Filling/CrudeEnergy.lean\#L46-L61}{{\leanfont Filling.\allowbreak{}hasCommonFlagVector\_of\_le}}}
\expandafter\newcommand\csname LN@existshasLevel\endcsname{\href{https://github.com/lsernau/grh-lean/blob/72fead102f663d5c7a407964a13dc0cf9d977207/grh-compact/src/New/Filling/CrudeEnergy.lean\#L67-L82}{{\leanfont Filling.\allowbreak{}CharsAgree.\allowbreak{}exists\_hasLevel}}}
\expandafter\newcommand\csname LN@energyleoflevelpairs\endcsname{\href{https://github.com/lsernau/grh-lean/blob/72fead102f663d5c7a407964a13dc0cf9d977207/grh-compact/src/New/Filling/CrudeEnergy.lean\#L108-L160}{{\leanfont Filling.\allowbreak{}energy\_le\_of\_level\_pairs}}}
\expandafter\newcommand\csname LN@energyleoflevellt\endcsname{\href{https://github.com/lsernau/grh-lean/blob/72fead102f663d5c7a407964a13dc0cf9d977207/grh-compact/src/New/Filling/CrudeEnergy.lean\#L166-L172}{{\leanfont Filling.\allowbreak{}energy\_le\_of\_level\_lt}}}
\expandafter\newcommand\csname LN@energyleoflinkEnergyle\endcsname{\href{https://github.com/lsernau/grh-lean/blob/72fead102f663d5c7a407964a13dc0cf9d977207/grh-compact/src/New/Filling/CrudeEnergy.lean\#L180-L191}{{\leanfont Filling.\allowbreak{}energy\_le\_of\_linkEnergy\_le}}}
\expandafter\newcommand\csname LQ@energyleoflinkEnergyle\endcsname{\leanquoteline{0}{lemma Filling.energy\_le\_of\_linkEnergy\_le R (h₂ : (∃ Λ, Λ ∉ R) →}\leanquoteline{2}{S.tailScale \textasciicircum{} 12 / X \textasciicircum{} 𝔰.k * S.mainEnergy ≤ C)}\leanquoteline{2}{(h : linkEnergy (T ∩ \{γ | S.termDisc γ ≠ 0\}) Ω R ≤ .ofReal (w * S.unipCovolume \textasciicircum{} 2)) :}\leanquoteline{2}{energy (S.amplitudeSum (σ := .full) T) Ω ≤ (1 + (S.b \textasciicircum{} 4)⁻¹) \textasciicircum{} 2 * w + (S.b \textasciicircum{} 12)⁻¹ * C}}
\expandafter\newcommand\csname LN@energy\endcsname{\href{https://github.com/lsernau/grh-lean/blob/72fead102f663d5c7a407964a13dc0cf9d977207/grh-compact/src/New/Filling/SchurTest.lean\#L20-L20}{{\leanfont Filling.\allowbreak{}energy}}}
\expandafter\newcommand\csname LQ@energy\endcsname{\leanquoteline{0}{noncomputable def energy := ∑ c, ∑ c', ∑ s with s ∈ Ω, ‖A c s c'‖ \textasciicircum{} 2 / S.unipCovolume \textasciicircum{} 2}}
\expandafter\newcommand\csname LN@linkEnergy\endcsname{\href{https://github.com/lsernau/grh-lean/blob/72fead102f663d5c7a407964a13dc0cf9d977207/grh-compact/src/New/Filling/SchurTest.lean\#L24-L25}{{\leanfont Filling.\allowbreak{}linkEnergy}}}
\expandafter\newcommand\csname LN@amplitudeSumsplit\endcsname{\href{https://github.com/lsernau/grh-lean/blob/72fead102f663d5c7a407964a13dc0cf9d977207/grh-compact/src/New/Filling/SchurTest.lean\#L30-L32}{{\leanfont Filling.\allowbreak{}amplitudeSum\_split}}}
\expandafter\newcommand\csname LN@energyaddle\endcsname{\href{https://github.com/lsernau/grh-lean/blob/72fead102f663d5c7a407964a13dc0cf9d977207/grh-compact/src/New/Filling/SchurTest.lean\#L35-L38}{{\leanfont Filling.\allowbreak{}energy\_add\_le}}}
\expandafter\newcommand\csname LN@energysubcutle\endcsname{\href{https://github.com/lsernau/grh-lean/blob/72fead102f663d5c7a407964a13dc0cf9d977207/grh-compact/src/New/Filling/SchurTest.lean\#L43-L58}{{\leanfont Filling.\allowbreak{}energy\_sub\_cut\_le}}}
\expandafter\newcommand\csname LN@classSize\endcsname{\href{https://github.com/lsernau/grh-lean/blob/72fead102f663d5c7a407964a13dc0cf9d977207/grh-compact/src/New/Filling/SchurTest.lean\#L65-L67}{{\leanfont Filling.\allowbreak{}classSize}}}
\expandafter\newcommand\csname LN@leclassSizeandcellMassle\endcsname{\href{https://github.com/lsernau/grh-lean/blob/72fead102f663d5c7a407964a13dc0cf9d977207/grh-compact/src/New/Filling/SchurTest.lean\#L72-L99}{{\leanfont Filling.\allowbreak{}IsNear.\allowbreak{}le\_classSize\_and\_cellMass\_le}}}
\expandafter\newcommand\csname LN@boxCutoff\endcsname{\href{https://github.com/lsernau/grh-lean/blob/72fead102f663d5c7a407964a13dc0cf9d977207/grh-compact/src/New/Filling/SchurTest.lean\#L103-L103}{{\leanfont Filling.\allowbreak{}boxCutoff}}}
\expandafter\newcommand\csname LN@schurWeight\endcsname{\href{https://github.com/lsernau/grh-lean/blob/72fead102f663d5c7a407964a13dc0cf9d977207/grh-compact/src/New/Filling/SchurTest.lean\#L107-L108}{{\leanfont Filling.\allowbreak{}schurWeight}}}
\expandafter\newcommand\csname LN@schurIntegral\endcsname{\href{https://github.com/lsernau/grh-lean/blob/72fead102f663d5c7a407964a13dc0cf9d977207/grh-compact/src/New/Filling/SchurTest.lean\#L114-L115}{{\leanfont Filling.\allowbreak{}schurIntegral}}}
\expandafter\newcommand\csname LN@inCutofbesselKernelnezero\endcsname{\href{https://github.com/lsernau/grh-lean/blob/72fead102f663d5c7a407964a13dc0cf9d977207/grh-compact/src/New/Filling/SchurTest.lean\#L121-L123}{{\leanfont Filling.\allowbreak{}inCut\_of\_besselKernel\_ne\_zero}}}
\expandafter\newcommand\csname LN@measurableboxCutoff\endcsname{\href{https://github.com/lsernau/grh-lean/blob/72fead102f663d5c7a407964a13dc0cf9d977207/grh-compact/src/New/Filling/SchurTest.lean\#L125-L126}{{\leanfont Filling.\allowbreak{}measurable\_boxCutoff}}}
\expandafter\newcommand\csname LN@measurableschurWeight\endcsname{\href{https://github.com/lsernau/grh-lean/blob/72fead102f663d5c7a407964a13dc0cf9d977207/grh-compact/src/New/Filling/SchurTest.lean\#L128-L129}{{\leanfont Filling.\allowbreak{}measurable\_schurWeight}}}
\expandafter\newcommand\csname LN@amplitudeAbssqle\endcsname{\href{https://github.com/lsernau/grh-lean/blob/72fead102f663d5c7a407964a13dc0cf9d977207/grh-compact/src/New/Filling/SchurTest.lean\#L131-L140}{{\leanfont Filling.\allowbreak{}amplitudeAbs\_sq\_le}}}
\expandafter\newcommand\csname LN@lintegralboxCutoffleofaverage\endcsname{\href{https://github.com/lsernau/grh-lean/blob/72fead102f663d5c7a407964a13dc0cf9d977207/grh-compact/src/New/Filling/SchurTest.lean\#L145-L169}{{\leanfont Filling.\allowbreak{}lintegral\_boxCutoff\_le\_of\_average}}}
\expandafter\newcommand\csname LN@termsofclassspec\endcsname{\href{https://github.com/lsernau/grh-lean/blob/72fead102f663d5c7a407964a13dc0cf9d977207/grh-compact/src/New/Filling/ThinClasses.lean\#L25-L39}{{\leanfont Filling.\allowbreak{}terms\_of\_class\_spec}}}
\expandafter\newcommand\csname LN@tsumschurIntegralle\endcsname{\href{https://github.com/lsernau/grh-lean/blob/72fead102f663d5c7a407964a13dc0cf9d977207/grh-compact/src/New/Filling/ThinClasses.lean\#L48-L123}{{\leanfont Filling.\allowbreak{}tsum\_schurIntegral\_le}}}
\expandafter\newcommand\csname LQ@tsumschurIntegralle\endcsname{\leanquoteline{0}{lemma Filling.tsum\_schurIntegral\_le : ∑' δ : doubleCoset γ (S.classUnip₀ c) (S.classUnip₀ c'),}\leanquoteline{2}{S.schurIntegral c c' δ ≤ .ofReal (S.unipCovolume * (card S.BesselCoord * S.tailScale \textasciicircum{} 2 *}\leanquoteline{2}{(S.b \textasciicircum{} 4 * S.g)⁻¹ \textasciicircum{} (min 𝔰.ϱ j * 𝔰.k * numInfPlaces E) * (S.xRowsIndex j (S.termPosition c c'}\leanquoteline{2}{γ) / √(S.minorNorm j γ)) \textasciicircum{} (𝔰.k + 1)))}}
\expandafter\newcommand\csname LN@tsumamplitudeAbsclassle\endcsname{\href{https://github.com/lsernau/grh-lean/blob/72fead102f663d5c7a407964a13dc0cf9d977207/grh-compact/src/New/Filling/ThinClasses.lean\#L131-L149}{{\leanfont Filling.\allowbreak{}tsum\_amplitudeAbs\_class\_le}}}
\newcommand{\ConcordanceRows}{\crow{sumlatticeDiscinvle}{\href{https://github.com/lsernau/grh-lean/blob/72fead102f663d5c7a407964a13dc0cf9d977207/grh-compact/src/New/Filling/BlockBound.lean\#L21-L75}{{\leanfont Filling.\allowbreak{}sum\_latticeDisc\_inv\_le}}}{BlockBound.lean}{21--75}
\crow{normsharedvectorlt}{\href{https://github.com/lsernau/grh-lean/blob/72fead102f663d5c7a407964a13dc0cf9d977207/grh-compact/src/New/Filling/BlockBound.lean\#L93-L125}{{\leanfont Filling.\allowbreak{}HasLevel.\allowbreak{}norm\_shared\_vector\_lt}}}{BlockBound.lean}{93--125}
\crow{tsuminvPowWeightleofcodimle}{\href{https://github.com/lsernau/grh-lean/blob/72fead102f663d5c7a407964a13dc0cf9d977207/grh-compact/src/New/Filling/BlockBound.lean\#L130-L215}{{\leanfont Filling.\allowbreak{}tsum\_invPowWeight\_le\_of\_codim\_le}}}{BlockBound.lean}{130--215}
\crow{submoduleIndexgIdelesmul}{\href{https://github.com/lsernau/grh-lean/blob/72fead102f663d5c7a407964a13dc0cf9d977207/grh-compact/src/New/Filling/BlockBound.lean\#L222-L226}{{\leanfont Filling.\allowbreak{}submoduleIndex\_gIdele\_smul}}}{BlockBound.lean}{222--226}
\crow{sumpairWeightle}{\href{https://github.com/lsernau/grh-lean/blob/72fead102f663d5c7a407964a13dc0cf9d977207/grh-compact/src/New/Filling/BlockBound.lean\#L232-L297}{{\leanfont Filling.\allowbreak{}sum\_pairWeight\_le}}}{BlockBound.lean}{232--297}
\crow{enormamplitudeSumleclassSizetop}{\href{https://github.com/lsernau/grh-lean/blob/72fead102f663d5c7a407964a13dc0cf9d977207/grh-compact/src/New/Filling/BlockBound.lean\#L304-L363}{{\leanfont Filling.\allowbreak{}enorm\_amplitudeSum\_le\_classSize\_top}}}{BlockBound.lean}{304--363}
\crow{enormamplitudeSumleclassSize}{\href{https://github.com/lsernau/grh-lean/blob/72fead102f663d5c7a407964a13dc0cf9d977207/grh-compact/src/New/Filling/BlockBound.lean\#L368-L394}{{\leanfont Filling.\allowbreak{}enorm\_amplitudeSum\_le\_classSize}}}{BlockBound.lean}{368--394}
\crow{amplitudeSumeqzeroofnotcharsAgree}{\href{https://github.com/lsernau/grh-lean/blob/72fead102f663d5c7a407964a13dc0cf9d977207/grh-compact/src/New/Filling/CrudeEnergy.lean\#L24-L41}{{\leanfont Filling.\allowbreak{}amplitudeSum\_eq\_zero\_of\_not\_charsAgree}}}{CrudeEnergy.lean}{24--41}
\crow{hasCommonFlagVectorofle}{\href{https://github.com/lsernau/grh-lean/blob/72fead102f663d5c7a407964a13dc0cf9d977207/grh-compact/src/New/Filling/CrudeEnergy.lean\#L46-L61}{{\leanfont Filling.\allowbreak{}hasCommonFlagVector\_of\_le}}}{CrudeEnergy.lean}{46--61}
\crow{existshasLevel}{\href{https://github.com/lsernau/grh-lean/blob/72fead102f663d5c7a407964a13dc0cf9d977207/grh-compact/src/New/Filling/CrudeEnergy.lean\#L67-L82}{{\leanfont Filling.\allowbreak{}CharsAgree.\allowbreak{}exists\_hasLevel}}}{CrudeEnergy.lean}{67--82}
\crow{energyleoflevelpairs}{\href{https://github.com/lsernau/grh-lean/blob/72fead102f663d5c7a407964a13dc0cf9d977207/grh-compact/src/New/Filling/CrudeEnergy.lean\#L108-L160}{{\leanfont Filling.\allowbreak{}energy\_le\_of\_level\_pairs}}}{CrudeEnergy.lean}{108--160}
\crow{energyleoflevellt}{\href{https://github.com/lsernau/grh-lean/blob/72fead102f663d5c7a407964a13dc0cf9d977207/grh-compact/src/New/Filling/CrudeEnergy.lean\#L166-L172}{{\leanfont Filling.\allowbreak{}energy\_le\_of\_level\_lt}}}{CrudeEnergy.lean}{166--172}
\crow{energyleoflinkEnergyle}{\href{https://github.com/lsernau/grh-lean/blob/72fead102f663d5c7a407964a13dc0cf9d977207/grh-compact/src/New/Filling/CrudeEnergy.lean\#L180-L191}{{\leanfont Filling.\allowbreak{}energy\_le\_of\_linkEnergy\_le}}}{CrudeEnergy.lean}{180--191}
\crow{energy}{\href{https://github.com/lsernau/grh-lean/blob/72fead102f663d5c7a407964a13dc0cf9d977207/grh-compact/src/New/Filling/SchurTest.lean\#L20-L20}{{\leanfont Filling.\allowbreak{}energy}}}{SchurTest.lean}{20--20}
\crow{linkEnergy}{\href{https://github.com/lsernau/grh-lean/blob/72fead102f663d5c7a407964a13dc0cf9d977207/grh-compact/src/New/Filling/SchurTest.lean\#L24-L25}{{\leanfont Filling.\allowbreak{}linkEnergy}}}{SchurTest.lean}{24--25}
\crow{amplitudeSumsplit}{\href{https://github.com/lsernau/grh-lean/blob/72fead102f663d5c7a407964a13dc0cf9d977207/grh-compact/src/New/Filling/SchurTest.lean\#L30-L32}{{\leanfont Filling.\allowbreak{}amplitudeSum\_split}}}{SchurTest.lean}{30--32}
\crow{energyaddle}{\href{https://github.com/lsernau/grh-lean/blob/72fead102f663d5c7a407964a13dc0cf9d977207/grh-compact/src/New/Filling/SchurTest.lean\#L35-L38}{{\leanfont Filling.\allowbreak{}energy\_add\_le}}}{SchurTest.lean}{35--38}
\crow{energysubcutle}{\href{https://github.com/lsernau/grh-lean/blob/72fead102f663d5c7a407964a13dc0cf9d977207/grh-compact/src/New/Filling/SchurTest.lean\#L43-L58}{{\leanfont Filling.\allowbreak{}energy\_sub\_cut\_le}}}{SchurTest.lean}{43--58}
\crow{classSize}{\href{https://github.com/lsernau/grh-lean/blob/72fead102f663d5c7a407964a13dc0cf9d977207/grh-compact/src/New/Filling/SchurTest.lean\#L65-L67}{{\leanfont Filling.\allowbreak{}classSize}}}{SchurTest.lean}{65--67}
\crow{leclassSizeandcellMassle}{\href{https://github.com/lsernau/grh-lean/blob/72fead102f663d5c7a407964a13dc0cf9d977207/grh-compact/src/New/Filling/SchurTest.lean\#L72-L99}{{\leanfont Filling.\allowbreak{}IsNear.\allowbreak{}le\_classSize\_and\_cellMass\_le}}}{SchurTest.lean}{72--99}
\crow{boxCutoff}{\href{https://github.com/lsernau/grh-lean/blob/72fead102f663d5c7a407964a13dc0cf9d977207/grh-compact/src/New/Filling/SchurTest.lean\#L103-L103}{{\leanfont Filling.\allowbreak{}boxCutoff}}}{SchurTest.lean}{103--103}
\crow{schurWeight}{\href{https://github.com/lsernau/grh-lean/blob/72fead102f663d5c7a407964a13dc0cf9d977207/grh-compact/src/New/Filling/SchurTest.lean\#L107-L108}{{\leanfont Filling.\allowbreak{}schurWeight}}}{SchurTest.lean}{107--108}
\crow{schurIntegral}{\href{https://github.com/lsernau/grh-lean/blob/72fead102f663d5c7a407964a13dc0cf9d977207/grh-compact/src/New/Filling/SchurTest.lean\#L114-L115}{{\leanfont Filling.\allowbreak{}schurIntegral}}}{SchurTest.lean}{114--115}
\crow{inCutofbesselKernelnezero}{\href{https://github.com/lsernau/grh-lean/blob/72fead102f663d5c7a407964a13dc0cf9d977207/grh-compact/src/New/Filling/SchurTest.lean\#L121-L123}{{\leanfont Filling.\allowbreak{}inCut\_of\_besselKernel\_ne\_zero}}}{SchurTest.lean}{121--123}
\crow{measurableboxCutoff}{\href{https://github.com/lsernau/grh-lean/blob/72fead102f663d5c7a407964a13dc0cf9d977207/grh-compact/src/New/Filling/SchurTest.lean\#L125-L126}{{\leanfont Filling.\allowbreak{}measurable\_boxCutoff}}}{SchurTest.lean}{125--126}
\crow{measurableschurWeight}{\href{https://github.com/lsernau/grh-lean/blob/72fead102f663d5c7a407964a13dc0cf9d977207/grh-compact/src/New/Filling/SchurTest.lean\#L128-L129}{{\leanfont Filling.\allowbreak{}measurable\_schurWeight}}}{SchurTest.lean}{128--129}
\crow{amplitudeAbssqle}{\href{https://github.com/lsernau/grh-lean/blob/72fead102f663d5c7a407964a13dc0cf9d977207/grh-compact/src/New/Filling/SchurTest.lean\#L131-L140}{{\leanfont Filling.\allowbreak{}amplitudeAbs\_sq\_le}}}{SchurTest.lean}{131--140}
\crow{lintegralboxCutoffleofaverage}{\href{https://github.com/lsernau/grh-lean/blob/72fead102f663d5c7a407964a13dc0cf9d977207/grh-compact/src/New/Filling/SchurTest.lean\#L145-L169}{{\leanfont Filling.\allowbreak{}lintegral\_boxCutoff\_le\_of\_average}}}{SchurTest.lean}{145--169}
\crow{termsofclassspec}{\href{https://github.com/lsernau/grh-lean/blob/72fead102f663d5c7a407964a13dc0cf9d977207/grh-compact/src/New/Filling/ThinClasses.lean\#L25-L39}{{\leanfont Filling.\allowbreak{}terms\_of\_class\_spec}}}{ThinClasses.lean}{25--39}
\crow{tsumschurIntegralle}{\href{https://github.com/lsernau/grh-lean/blob/72fead102f663d5c7a407964a13dc0cf9d977207/grh-compact/src/New/Filling/ThinClasses.lean\#L48-L123}{{\leanfont Filling.\allowbreak{}tsum\_schurIntegral\_le}}}{ThinClasses.lean}{48--123}
\crow{tsumamplitudeAbsclassle}{\href{https://github.com/lsernau/grh-lean/blob/72fead102f663d5c7a407964a13dc0cf9d977207/grh-compact/src/New/Filling/ThinClasses.lean\#L131-L149}{{\leanfont Filling.\allowbreak{}tsum\_amplitudeAbs\_class\_le}}}{ThinClasses.lean}{131--149}}
% the notations of the defined functions: one slot for each explicit argument of the Lean declaration
\newcommand{\csize}[3]{\operatorname{cs}_{#1}(#2;#3)}
\newcommand{\bcut}[2]{\beta^{0}_{#1,#2}}
\newcommand{\schw}[4]{\mathrm{sw}_{#1}(#2;#3,#4)}
\newcommand{\schi}[4]{\mathrm{SI}_{#1}(#2,#3;#4)}
\newcommand{\UnitMembers}{497}
\newcommand{\UnitInNote}{31}
\newcommand{\UnitProofs}{179}
\newcommand{\UnitSteps}{1371}
\newcommand{\UnitAllSteps}{1479}

% generated by notes/filling-setup/shared/mkspec.py from shared/notation-148ab.json; do not edit.
% \input after data.tex: a macro that the build made from the spec is kept, the others are defined here (same TeX).
\providecommand{\fs}{\mathfrak{s}}                      % the shape (rho, kappa, k_0)
\providecommand{\sT}{\mathsf{T}}                        % a finite set of finite places
\providecommand{\cL}{\mathcal{L}}                       % a level
\providecommand{\AEf}{\mathbb{A}_{E,f}}
\providecommand{\AdE}{\mathbb{A}_{E}}
\providecommand{\hO}{\hat{O}_E}
\providecommand{\GL}{\mathrm{GL}}
\providecommand{\full}{\mathrm{full}}
\providecommand{\cut}{\mathrm{cut}}
\providecommand{\rh}[1]{\varrho(#1)}                    % Shape.ϱ
\providecommand{\kk}[1]{k(#1)}                          % Shape.k = #κ
\providecommand{\nn}[1]{n(#1)}                          % Shape.n = ϱ + 1
\providecommand{\rr}[1]{r(#1)}                          % Shape.r = 2ϱ + k + 1
\providecommand{\Ix}[1]{\operatorname{Ind}(#1)}                          % Shape.BesselIndex
\providecommand{\Mx}[1]{\operatorname{Mid}(#1)}                          % Shape.MiddleIndex
\providecommand{\ninf}[1]{f(#1)}                        % numInfPlaces E
\providecommand{\ET}[1]{E_{#1}}                         % LocProd T
\providecommand{\bS}[1]{b(#1)}                          % S.b
\providecommand{\NS}[1]{N(#1)}                          % S.N
\providecommand{\qS}[2]{q_{#2}(#1)}                     % S.levelPrime i
\providecommand{\piS}[1]{\bar\pi(#1)}                   % S.piBar
\providecommand{\vpS}[1]{\varpi(#1)}                    % S.ϖ
\providecommand{\ChS}[2]{\mathrm{Ch}_{#2}(#1)}          % S.chosenPlaces i
\providecommand{\KcolS}[1]{\mathsf{K}^{\mathrm{col}}(#1)}        % S.columnLevel
\providecommand{\BsysS}[2]{\mathfrak{B}_{#2}(#1)}       % S.besselSys v
\providecommand{\mass}[2]{\operatorname{mass}(#1,#2)}   % mass E r
% Lib: forms, groups and frames (letters of the notes of unit 132)
\providecommand{\fr}[1]{\mathfrak{f}_{#1}}              % the frame vector of index a (e_a for a < rho, z for a = rho)
\providecommand{\Trg}[2]{T^{\mathrm{rat}}(#1;#2)}       % Shape.middleGram N: the form diag(1,...,1,-N)
\providecommand{\JT}[1]{J(#1)}                          % Shape.Form.splitGram
\providecommand{\GT}[1]{G(#1)}                          % Shape.Form.G
\providecommand{\Gone}[1]{G_{1}(#1)}                    % Shape.G₁
\providecommand{\UT}[1]{\mathcal{U}(#1)}                % besselU
\providecommand{\UoT}[1]{\mathcal{U}^{0}(#1)}           % besselU0
\providecommand{\HT}[1]{\mathcal{H}(#1)}                % besselH
\providecommand{\ST}[1]{\mathcal{S}(#1)}                % besselS
\providecommand{\Hl}[2]{\mathcal{H}_{#1}(#2)}           % GT l
\providecommand{\Ql}[2]{\mathcal{Q}_{#1}(#2)}           % unitriGroup l
\providecommand{\sd}[1]{\operatorname{sd}(#1)}          % simpleRootSum
\providecommand{\ZT}[2]{Z_{#1}(#2)}                     % Shape.Form.frameGram
\providecommand{\Zr}[1]{Z(#1)}                          % realFrameGram
\providecommand{\lmin}[2]{\mathrm{D}_{#1}(#2)}          % leadingMinor Z l
\providecommand{\weyl}[2]{w_{#1}(#2)}                   % weylElem l p
\providecommand{\gix}[1]{[#1]}                          % genIndex (the generalized index)
\providecommand{\den}[1]{\operatorname{den}(#1)}        % denomIndex
\providecommand{\Kf}[2]{\mathsf{K}_{f}(#1,#2)}                   % Kf E n
\providecommand{\Gf}[2]{G_{f}(#1,#2)}                   % Gf E n
\providecommand{\volO}[1]{\operatorname{vol}(\hat O_{#1})} % volÔ E
\providecommand{\rdE}[1]{\operatorname{rd}_{#1}}        % rootDiscr E
\providecommand{\dE}[1]{d_{#1}}                         % discr E
\providecommand{\CC}{\mathbb{C}}
\providecommand{\RR}{\mathbb{R}}
\providecommand{\NN}{\mathbb{N}}
\providecommand{\vol}{\operatorname{vol}}
% ----- the definitions of chunks a and b of unit 148 (and, for this note, its own additions)
\providecommand{\Tlatg}[2]{T^{\mathrm{lat}}(#1;#2)}                  % latMiddleGram (𝔰, N)
\providecommand{\mbc}[2]{\Lambda^{\circ}(#1;#2)}                     % middleBaseChange (𝔰, p)
\providecommand{\blb}[2]{\Lambda(#1;#2)}                             % besselLatticeBasis (𝔰, p)
\providecommand{\fdisc}[2]{\mathsf{g}(#1;#2)}                        % frameDisc (N, Z)
\providecommand{\fpair}[1]{[\mathfrak{f}\mid #1\mathfrak{f}]}        % framePair (γ)
\providecommand{\ifpair}[1]{[\mathfrak{e}\mid #1\mathfrak{e}]}       % isotropicFramePair (γ)
\providecommand{\lmat}[1]{\mathrm{L}(#1)}                            % linkMatrix (γ)
\providecommand{\link}[1]{\operatorname{link}(#1)}                   % link (γ)
\providecommand{\damp}[2]{\operatorname{dp}_{#1}(#2)}                % dampedPower (k, r)
\providecommand{\hw}[3]{\Lambda^{\mathbb{C}}_{#1}(#2;#3)}            % headWeight (K, i, p)
\providecommand{\Sieve}[8]{\operatorname{Sieve}(#1,#2,#3,#4,#5,#6,#7,#8)} % SieveSystem (u, ι, n, J, ψ, En, Λ, β)
\providecommand{\fhw}[3]{\Lambda_{#1}(#2;#3)}                        % finHeadWeight (K, i, p)
\providecommand{\occ}[1]{c_{G}(#1)}                                  % Shape.openCellConst (𝔰)
\providecommand{\bandc}[1]{c_{\mathrm{band}}(#1)}                    % Shape.bandConst (𝔰)
\providecommand{\bsize}[1]{\operatorname{bs}(#1)}                    % besselSize (y)
\providecommand{\InCut}[2]{#1\in\operatorname{Cut}(#2)}              % InCut (y, b)
\providecommand{\frbox}[1]{\operatorname{FBox}(#1)}                  % frameBox (ρ)
\providecommand{\gdr}[1]{\mathsf{g}^{\mathbb{R}}(#1)}                % gramDiscReal (Z)
\providecommand{\cellm}[4]{\operatorname{cm}(#1;#2,#3,#4)}           % cellMass (y, deg, ρ, d)
\providecommand{\bcorr}[3]{\mathrm{J}(#1,#2;#3)}                     % besselCorrelation (μ, W, γ)
\providecommand{\SchurLe}[2]{\operatorname{Sch}(#1)\le #2}           % SchurLe (W, D)
\providecommand{\HasSB}[4]{\operatorname{SchB}(#1,#2;#3,#4)}         % HasSchurBound (W, rep, B, D)
\providecommand{\ncorr}[3]{\mathrm{J}^{\circ}(#1;#2,#3)}             % normCorrelation (y, W, p)
\providecommand{\BandL}[4]{\operatorname{BL}(#1,#2,#3,#4)}           % HasBandLimits (b, W, deg, J)
\providecommand{\CutPair}[4]{\operatorname{CP}(#1,#2,#3,#4)}         % IsCutPair (b, f, rep, W)
\providecommand{\BSys}[3]{\operatorname{BSys}(#1,#2,#3)}             % BesselSystem (𝔰, b, f)
\providecommand{\Fill}[2]{\operatorname{Fill}(#1,#2)}                % Filling (E, 𝔰)
\providecommand{\gS}[1]{g(#1)}                                       % Filling.g (S)
\providecommand{\Trat}[2]{T^{\mathrm{rat}}_{#1}(#2)}                 % Filling.ratMiddle (S, R)
\providecommand{\Tlat}[2]{T^{\mathrm{lat}}_{#1}(#2)}                 % Filling.latMiddle (S, R)
\providecommand{\GS}[2]{G_{#1}(#2)}                                  % Filling.G (S, R)
\providecommand{\GFin}[1]{G^{\mathrm{lat}}_{#1}}                     % Filling.GFin (S)
\providecommand{\GRes}[1]{\mathsf{G}(#1)}                            % Filling.GRes (S)
\providecommand{\Bc}[1]{\mathrm{Bc}(#1)}                             % Filling.BesselCoord (S)
\providecommand{\Cl}[1]{\mathrm{Cl}(#1)}                             % Filling.Classes (S)
\providecommand{\Zf}[2]{Z_{#1}(#2)}                                  % Filling.frameProducts (S, γ)
\providecommand{\KS}[1]{\mathsf{K}(#1)}                              % Filling.latStabilizer (S)
\providecommand{\redS}[1]{\mathrm{red}_{#1}}                         % Filling.reduction (S)
\providecommand{\KgS}[1]{\mathsf{K}_g(#1)}                           % Filling.principalLevel (S)
\providecommand{\redim}[2]{\overline{#2}^{\,#1}}                     % Filling.reducedImage (S, A)
\providecommand{\kapS}[2]{\kappa_{#2}(#1)}                           % Filling.kappa (S, v)
\providecommand{\bscale}[3]{\mathrm{sc}_{#1,#2}(#3)}                 % Filling.basisScale (S, v, i)
\providecommand{\sdc}[1]{\psi_{#1}}                                  % Filling.selfDualChar (S)
\providecommand{\fscale}[3]{\mathrm{fs}_{#1}(#2,#3)}                 % Filling.frameScaleSq (S, v, a)
\providecommand{\fsprod}[1]{\mathrm{fs}(#1)}                         % Filling.frameScaleProd (S)
\providecommand{\ucov}[1]{\mathrm{cov}(#1)}                          % Filling.unipCovolume (S)
\providecommand{\siegC}[1]{c_{*}(#1)}                                % Filling.siegelConst (S)
\providecommand{\tdisc}[2]{\mathsf{g}_{#1}(#2)}                      % Filling.termDisc (S, γ)
\providecommand{\bsGL}[2]{\mathrm{Sc}_{#1,#2}}                       % Filling.basisScaleGL (S, v)
\providecommand{\latB}[1]{\Lambda(#1)}                               % Filling.latBasis (S)
\providecommand{\tolb}[1]{\mathrm{lb}_{#1}}                          % Filling.toLatBasis (S)
\providecommand{\radl}[1]{\mathrm{r}^{\mathbb{A}}_{#1}}              % Filling.ratToAdelic (S)
\providecommand{\ainf}[1]{\mathrm{a}^{\infty}_{#1}}                  % Filling.adelicToInf (S)
\providecommand{\afin}[1]{\mathrm{a}^{f}_{#1}}                       % Filling.adelicToFin (S)
\providecommand{\rinf}[1]{\mathrm{r}^{\infty}_{#1}}                  % Filling.ratToInf (S)
\providecommand{\rfin}[1]{\mathrm{r}^{f}_{#1}}                       % Filling.ratToFin (S)
\providecommand{\incH}[1]{\mathrm{inc}_{#1}}                         % Filling.inclHFin (S)
\providecommand{\hc}[2]{h_{#1}(#2)}                                  % Filling.h (S, c)
\providecommand{\pos}[4]{\mathrm{pos}_{#1}(#2,#3;#4)}                % Filling.termPosition (S, c, c', γ)
\providecommand{\terms}[4]{\mathrm{Tm}_{#1}(#2,#3,#4)}               % Filling.termsOver (S, c, s, c')
\providecommand{\clU}[2]{\mathcal{U}_{#1,#2}}                        % Filling.classUnip (S, c)
\providecommand{\clUo}[2]{\mathcal{U}^{0}_{#1,#2}}                   % Filling.classUnip₀ (S, c)
\providecommand{\bxs}[1]{\operatorname{bx}(#1)}                      % Filling.boxSize (S)
\providecommand{\IsBP}[3]{#3\in\operatorname{BP}(#1,#2)}             % Filling.IsBoxPartition (S, c, ω)
\providecommand{\besfn}[4]{W^{#1}_{#2}(#3)_{#4}}                     % Filling.besselFn (σ, S, y, i)
\providecommand{\besker}[4]{\mathcal{W}^{#1}_{#2}(#3;#4)}            % Filling.besselKernel (σ, S, γ, y)
\providecommand{\cdim}[1]{\mathrm{cd}(#1)}                           % Filling.columnDim (S)
\providecommand{\bdeg}[1]{\deg(#1)}                                  % Filling.besselDegree (S)
\providecommand{\mainE}[1]{\mathsf{d}(#1)}                           % Filling.mainEnergy (S)
\providecommand{\uord}[1]{\mathrm{uo}(#1)}                           % Filling.unipOrder (S)
\providecommand{\packC}[1]{a_{*}(#1)}                                % Filling.packetConst (S)
\providecommand{\besE}[1]{\mathsf{d}_{1}(#1)}                        % Filling.besselEnergy (S)
\providecommand{\tailS}[1]{\mathrm{tl}(#1)}                          % Filling.tailScale (S)
\providecommand{\codimF}[1]{\mathrm{cf}(#1)}                         % Filling.codimFactor (S)
\providecommand{\IsAdm}[1]{#1\text{ is admissible}}                  % Filling.IsAdmissible (S)
\providecommand{\cutoff}[2]{\beta_{#1,#2}}                           % Filling.classCutoff (S, c)
\providecommand{\amp}[5]{a^{#1}_{#2}(#3,#4;#5)}                      % Filling.amplitude (σ, S, c, c', γ)
\providecommand{\ampabs}[5]{|a|^{#1}_{#2}(#3,#4;#5)}                 % Filling.amplitudeAbs (σ, S, c, c', γ)
\providecommand{\ampsum}[6]{A^{#1}_{#2}(#3;#4,#5,#6)}                % Filling.amplitudeSum (σ, S, T, c, s, c')
\providecommand{\gid}[1]{\hat{g}(#1)}                                % Filling.gIdele (S)
\providecommand{\uchar}[2]{\psi^{\mathcal{U}}_{#1}(#2)}              % Filling.unipotentChar (S, u)
\providecommand{\Agr}[2]{#2\in\operatorname{Agr}(#1)}                % Filling.Agrees (S, s)
\providecommand{\IsStab}[2]{#2\in\operatorname{Stab}(#1)}            % Filling.IsStable (S, T)
\providecommand{\tlink}[2]{\mathrm{Tm}^{\mathrm{lk}}_{#1}(#2)}       % Filling.termsOfLink (S, Λ)
\providecommand{\clB}[2]{\mathrm{B}_{#1}(#2)}                        % Filling.classBasis (S, c)
\providecommand{\ioia}[2]{\operatorname{adj}_{#1}^{#2}}              % Filling.idOrInvAdjoint (S, b)
\providecommand{\rag}[2]{\mathrm{red}_{#1,#2}}                       % Filling.reduceAboveG (S, x)
\providecommand{\vfix}[3]{\mathrm{Fix}_{#1,#2}(#3)}                  % Filling.vectorFixer (S, w, d)
\providecommand{\FixN}[2]{\operatorname{fix}(#1)\ge #2}              % Filling.FixesNormAtLeast (s, x)
\providecommand{\fixset}[3]{\mathrm{Fix}_{#1}(#2,#3)}                % Filling.fixingSet (S, Sg, d)
\providecommand{\zfp}[1]{\mathrm{L}_{z}(#1)}                         % Filling.zFormPair (T)
\providecommand{\clink}[2]{\operatorname{link}^{\circ}(#1,#2)}       % commonLink (γ, γ')
\providecommand{\HCFV}[4]{\operatorname{CF}_{#4}(#1;#2,#3)}          % HasCommonFlagVector (γ, c, c', a)
\providecommand{\IsNF}[2]{#2\in\operatorname{NF}_{#1}}               % IsNormalForm (j, γ)
\providecommand{\ratG}[1]{J_{#1}}                                    % Filling.ratGram (S)
\providecommand{\lform}[1]{Q^{\pm}_{#1}}                             % Filling.linkForm (S)
\providecommand{\clC}[2]{\mathrm{C}_{#1}(#2)}                        % Filling.classCoord (S, c)
\providecommand{\clCX}[2]{\mathrm{C}^{\circ}_{#1}(#2)}               % Filling.classCoordX₀ (S, c)
\providecommand{\cw}[2]{\operatorname{cw}_{#2}(#1)}                  % Filling.countWeight (S, i)
\providecommand{\HasLev}[3]{#3\in\operatorname{Lv}_{#2}(#1)}         % Filling.HasLevel (S, j, γ)
\providecommand{\Cong}[5]{#4\sim_{#1;#2,#3}#5}                       % Filling.CongruentTerms (S, c, c', γ, γ')
\providecommand{\tdnorm}[2]{\|\mathsf{g}\|_{#1}(#2)}                 % Filling.termDiscNorm (S, γ)
\providecommand{\bxp}[1]{\operatorname{BxP}(#1)}                     % Filling.boxedPoints (S)
\providecommand{\IsNr}[2]{#2\in\operatorname{Near}(#1)}              % Filling.IsNear (S, γ)
\providecommand{\ChA}[2]{#2\in\operatorname{ChA}(#1)}                % Filling.CharsAgree (S, γ)
\providecommand{\fadm}[1]{(#1)_{\mathbb{A}_{f}}}                     % Filling.finAdeleMat (M)
\providecommand{\slb}[2]{\mathrm{m}_{#1}(#2)}                        % Filling.sublatticeBasis (H, V)
\providecommand{\sld}[2]{\operatorname{disc}_{#1}(#2)}               % Filling.sublatticeDisc (u, V)
\providecommand{\lmet}[1]{Q^{\mathrm{link}}_{#1}}                    % Filling.linkMetric (S)
\providecommand{\ldsq}[3]{\mathrm{ld}_{#1}(#2,#3)}                   % Filling.linkDiscSq (S, v, e)
\providecommand{\HIL}[3]{#3\in\operatorname{IL}(#1;#2)}              % Filling.HasIntegralLifts (c, γ, m)
\providecommand{\IdxB}[1]{\mathrm{Ib}(#1)}                           % Filling.IndexBelow (j)
\providecommand{\xrows}[1]{#1^{\mathrm{x}}}                          % Filling.xRows (Z)
\providecommand{\padl}[2]{\operatorname{pad}_{#1}(#2)}               % Filling.paddedLeading (j, Z)
\providecommand{\fmid}[1]{(#1\mathfrak{f})^{\mathrm{mid}}}           % Filling.frameMiddle (γ)
\providecommand{\latd}[2]{\operatorname{disc}_{#1,#2}}               % Filling.latticeDisc (S, c)
\providecommand{\fvi}[5]{\mathrm{o}_{#1,#2}(#3;#4,#5)}               % Filling.frameValueIndex (S, c, k, a, V)
\providecommand{\xri}[3]{\mathrm{i}_{#1,#2}(#3)}                     % Filling.xRowsIndex (S, j, k)
\providecommand{\sati}[2]{\mathfrak{d}_{#1}(#2)}                     % Filling.saturationIndex (S, k)
\providecommand{\pinv}[5]{\mathrm{pi}_{#1}(#2,#3;#4,#5)}             % Filling.pairInvariant (S, c, c', γ, γ')
\providecommand{\telem}[3]{\mathrm{el}_{#1}(#2;#3)}                  % Filling.termElem (S, c, p)
\providecommand{\mnorm}[3]{\mathrm{mn}_{#1}(#2;#3)}                  % Filling.minorNorm (S, j, γ₁)
\providecommand{\cgauge}[5]{\mathrm{cg}_{#1}(#2;#3,#4;#5)}           % Filling.classGauge (S, j, c, c', γ)
\providecommand{\IsNRep}[4]{#4\in\operatorname{NRep}_{#1}(#2,#3)}    % Filling.IsNearRep (S, c, c', γ)
\providecommand{\ipw}[4]{\mathrm{ipw}_{#1}(#2;#3,#4)}                % Filling.invPowWeight (S, c, n, p)
\providecommand{\termw}[7]{\mathrm{w}_{#1}(#2,#3;#4;#5,#6;#7)}       % Filling.termWeight (S, j, i, c, p, p', V)
\providecommand{\pairw}[6]{\mathrm{pw}_{#1}(#2;#3,#4;#5,#6)}         % Filling.pairWeight (S, j, c, c', γ, γ')
\providecommand{\IsLab}[5]{#5\in\operatorname{Lab}_{#1}(#2;#3,#4)}   % Filling.IsLabelledTerm (S, j, c, t, p)
\providecommand{\energy}[3]{\mathrm{En}_{#1}(#2;#3)}                 % Filling.energy (S, A, Ω)
\providecommand{\linkE}[5]{\mathrm{LE}^{#1}_{#2}(#3;#4,#5)}          % Filling.linkEnergy (σ, S, T, Ω', R)
\providecommand{\csize}[3]{\operatorname{cs}_{#1}(#2;#3)}            % Filling.classSize (S, j, x)
\providecommand{\bcut}[2]{\beta^{0}_{#1,#2}}                         % Filling.boxCutoff (S, c)
\providecommand{\schw}[4]{\mathrm{sw}_{#1}(#2;#3,#4)}                % Filling.schurWeight (S, c, γ, y)
\providecommand{\schi}[4]{\mathrm{SI}_{#1}(#2,#3;#4)}                % Filling.schurIntegral (S, c, c', γ')
\providecommand{\charT}[2]{\psi_{#1,#2}}                             % Filling.charOnPlaces (S, T)
\providecommand{\fgain}[3]{\beta_{#1,#2}(#3)}                        % Filling.frequencyGain (S, T, ξ)
\providecommand{\ratpt}[3]{\mathrm{pt}_{#1}(#2;#3)}                  % Filling.ratPoint (T, l, t)
\providecommand{\Fits}[6]{\operatorname{Fits}(#1;#2,#3,#4,#5,#6)}    % Filling.Fits (S, V, l, kw, Y, η)
\providecommand{\stdK}[1]{\mathsf{K}^{\mathrm{std}}_{g}(#1)}         % Filling.stdPrincipalLevel (S)
\providecommand{\tailv}[3]{\operatorname{tv}(#1,#2,#3)}              % tailVector (π, t, x)


\linespread{1.08}                       % 11 pt on 14.7 pt
\setlength{\parindent}{0pt}
\setlength{\parskip}{5.5pt plus 1.5pt minus 0.5pt}
\setlength{\abovedisplayskip}{7.5pt plus 2pt minus 3pt}
\setlength{\belowdisplayskip}{7.5pt plus 2pt minus 3pt}
\setlength{\abovedisplayshortskip}{3pt plus 2pt}
\setlength{\belowdisplayshortskip}{6pt plus 2pt minus 2pt}
\emergencystretch=1.5em
\clubpenalty=10000 \widowpenalty=10000 \displaywidowpenalty=10000
\relpenalty=7000 \binoppenalty=9000
\raggedbottom
\pagestyle{fancy}\fancyhf{}\renewcommand{\headrulewidth}{0pt}
\fancyfoot[C]{\sffamily\fontsize{7.6}{9}\selectfont\color{inkfaint}\thepage\,/\,\pageref*{LastPage}}
\setlist{topsep=2pt,partopsep=0pt,parsep=0pt,itemsep=4pt,leftmargin=2.1em,labelsep=0.55em}

\newcommand{\Q}{\mathbb{Q}}
\newcommand{\Z}{\mathbb{Z}}
\newcommand{\F}{\mathbb{F}}
\newcommand{\len}{\operatorname{len}}
\newcommand{\cok}{\operatorname{cok}}
\newcommand{\End}{\operatorname{End}}
\newcommand{\Hom}{\operatorname{Hom}}
\newcommand{\Gal}{\operatorname{Gal}}
\newcommand{\diag}{\operatorname{diag}}
\newcommand{\kw}[1]{\emph{#1}}
\newcommand{\m}[1]{\mbox{$#1$}}        % a short formula that must not be broken across lines
\newcommand{\head}[1]{\par\addvspace{17pt plus 4pt minus 2pt}\Needspace{6\baselineskip}%
  \noindent{\fontsize{13.2}{16}\selectfont\bfseries\color{ink}#1}\par\nopagebreak\vspace{-1pt}}
\newcommand{\sect}[2]{\refstepcounter{section}\label{#2}\par\addvspace{19pt plus 4pt minus 2pt}\Needspace{7\baselineskip}%
  \noindent\llap{{\fontsize{13.2}{16}\selectfont\bfseries\color{inkfaint}\thesection}\hspace{9pt}}%
  {\fontsize{13.2}{16}\selectfont\bfseries\color{ink}#1}\par\nopagebreak\vspace{-1pt}}
\newcommand{\subhead}[1]{\par\addvspace{9pt plus 2pt}\Needspace{4\baselineskip}\noindent\textbf{#1}\enspace\ignorespaces}
% where a statement lives in the Lean archive
\newcommand{\leantag}{{\sffamily\bfseries\fontsize{6.3}{8}\selectfont\color{leancol}\addfontfeatures{LetterSpace=7}LEAN}}
\newcommand{\sep}{\unskip\nobreak\hspace{5pt}{\color{inkfaint}\textperiodcentered}\hspace{5pt plus 2pt}\ignorespaces}
\newcommand{\leanline}[1]{\par\nopagebreak\vspace{-2.5pt}{\fontsize{8.9}{12.4}\selectfont\color{inksoft}\raggedright
  \noindent\llap{\leantag\hspace{9pt}}#1\par}\vspace{1pt}}
% a displayed statement in a tinted box: #1 colour, #2 text
\newcommand{\statement}[2]{\par\addvspace{7pt}\begingroup\setlength{\fboxsep}{8.5pt}\setlength{\fboxrule}{0.6pt}%
  \noindent\fcolorbox{#1!30!white}{#1!3!white}{\parbox{\dimexpr\linewidth-2\fboxsep-2\fboxrule\relax}{#2}}%
  \endgroup\par\addvspace{7pt}}
\newcommand{\nt}[2]{\par\noindent\hangindent=1.1em\hangafter=1 #1\enspace{\color{inksoft}#2}\par\vspace{1pt}}
\newcommand{\legdot}[1]{\tikz[baseline=-0.62ex]\fill[#1] (0,0) circle (2.5pt);\hspace{3pt}}
\newcommand{\crow}[4]{\csname CS@#1\endcsname & #2 & \textcolor{inksoft}{#3} & \textcolor{inksoft}{#4}\\}
\newcommand{\leanquoteline}[2]{\par\noindent\hangindent=\dimexpr#1\fontdimen2\font+2.2em\relax\hangafter=1
  \hspace*{#1\fontdimen2\font}#2}
\newcommand{\quoted}[1]{\par\vspace{3pt}{\ttfamily\fontsize{7.5}{10.4}\selectfont\color{leancol}\raggedright
  \setlength{\parskip}{0pt}\leanquote{#1}\par}\vspace{4pt}}
\newcommand{\bx}[3]{{\sffamily\bfseries\fontsize{7.3}{9.2}\selectfont\color{#1!82!black}#2}\\[1.5pt]#3}
\newcommand{\cs}[2]{\expandafter\newcommand\csname CS@#1\endcsname{#2}}

% theorems
\newtheoremstyle{plainx}{7pt plus 2pt}{7pt plus 2pt}{\itshape}{}{\bfseries}{.}{0.5em}{}
\newtheoremstyle{uprightx}{7pt plus 2pt}{7pt plus 2pt}{}{}{\bfseries}{.}{0.5em}{}
\theoremstyle{plainx}
\newtheorem{theorem}{Theorem}[section]
\newtheorem{lemma}[theorem]{Lemma}
\newtheorem{proposition}[theorem]{Proposition}
\newtheorem{corollary}[theorem]{Corollary}
\theoremstyle{uprightx}
\newtheorem{definition}[theorem]{Definition}
\newtheorem{setting}[theorem]{Setting}          % hypotheses shared by several statements; each of them cites it by number
\newtheorem{remark}[theorem]{Remark}
\newtheorem{example}[theorem]{Example}
\renewcommand{\qedsymbol}{{\setlength{\fboxsep}{0pt}\setlength{\fboxrule}{0.5pt}\fbox{\rule{0pt}{0.46em}\hspace{0.46em}}}}
% what holds throughout the note: the ambient objects only (implicit and instance arguments), never a hypothesis
\newenvironment{conventions}{\par\addvspace{5pt}\noindent\textbf{Conventions.}\enspace\ignorespaces}{\par\addvspace{3pt}}
\makeatletter
\renewenvironment{proof}[1][Proof]{\par\pushQED{\qed}\normalfont\topsep0pt\partopsep0pt
  \trivlist\item[\hskip\labelsep\itshape #1.]\ignorespaces}{\popQED\endtrivlist\@endpefalse}
\makeatother

% inline formulas may break at relations and operations: with the long full forms of this unit a boxed formula
% often does not fit a line
\renewcommand{\m}[1]{$#1$}
% ----- shorthands of this note; each prints its full form with every argument (no argument is dropped)
\newcommand{\sk}{\kk{\fs}}
\newcommand{\sn}{\nn{\fs}}
\newcommand{\sro}{\rh{\fs}}
\newcommand{\fE}{\ninf{E}}
\newcommand{\bb}{\bS{S}}
\newcommand{\Nn}{\NS{S}}
\newcommand{\gq}{\gS{S}}
\newcommand{\tl}{\tailS{S}}
\newcommand{\nuS}{\ucov{S}}
\newcommand{\dm}{\mainE{S}}
\newcommand{\dd}{\besE{S}}
\newcommand{\ri}[1]{\rinf{S}(#1)}
\newcommand{\rfi}[1]{\rfin{S}(#1)}
\newcommand{\Ginf}{G_{\infty}}
\newcommand{\Tone}{\Trg{\fs}{1}}
\newcommand{\TE}{\Trat{S}{E}}
\newcommand{\Tl}{\Tlat{S}{\AEf}}
\newcommand{\UE}{\UT{\TE}}
\newcommand{\UoE}{\UoT{\TE}}
\newcommand{\Uone}{\UT{\Tone}}
\newcommand{\Uoone}{\UoT{\Tone}}
\newcommand{\Hone}{\HT{\Tone}}
\newcommand{\Sbar}{\redim{S}{\ST{\Tl}}}
\newcommand{\Wt}[3]{W^{#1}_{#2,#3}}
\newcommand{\It}[2]{I_{#2}(#1)}
\newcommand{\degt}[2]{\deg_{#2}(#1)}
\newcommand{\cj}[1]{\overline{#1}}
\newcommand{\Zl}[1]{\ZT{\Tl}{#1}}
\newcommand{\Mat}{\operatorname{Mat}}
\newcommand{\adj}{\operatorname{adj}}
\newcommand{\Nm}{\mathrm{N}_{E/\Q}}
\newcommand{\inp}[1]{\textup{(#1)}}
\newcommand{\ip}{\textbf{Inputs.}\enspace}
\newcommand{\ReP}{\operatorname{Re}}

% where each declaration of the concordance is treated
\cs{energy}{Definition~\ref{d:energy}}
\cs{linkEnergy}{Definition~\ref{d:energy}}
\cs{amplitudeSumsplit}{Lemma~\ref{l:split}}
\cs{energyaddle}{Lemma~\ref{l:add}}
\cs{energysubcutle}{Theorem C}
\cs{classSize}{Definition~\ref{d:cs}}
\cs{leclassSizeandcellMassle}{Lemma~\ref{l:cell}}
\cs{boxCutoff}{Definition~\ref{d:schur}}
\cs{schurWeight}{Definition~\ref{d:schur}}
\cs{schurIntegral}{Definition~\ref{d:schur}}
\cs{inCutofbesselKernelnezero}{Lemma~\ref{l:incut}}
\cs{measurableboxCutoff}{Lemma~\ref{l:cs}}
\cs{measurableschurWeight}{Lemma~\ref{l:cs}}
\cs{amplitudeAbssqle}{Lemma~\ref{l:cs}}
\cs{lintegralboxCutoffleofaverage}{Proposition~\ref{p:avg}}
\cs{termsofclassspec}{Lemma~\ref{l:class}}
\cs{tsumschurIntegralle}{Theorem A}
\cs{tsumamplitudeAbsclassle}{Theorem A}
\cs{sumlatticeDiscinvle}{Proposition~\ref{p:latdisc}}
\cs{normsharedvectorlt}{Lemma~\ref{l:short}}
\cs{tsuminvPowWeightleofcodimle}{Proposition~\ref{p:count}}
\cs{submoduleIndexgIdelesmul}{Lemma~\ref{l:gindex}}
\cs{sumpairWeightle}{Proposition~\ref{p:pairs}}
\cs{enormamplitudeSumleclassSizetop}{Theorem B}
\cs{enormamplitudeSumleclassSize}{Theorem B}
\cs{amplitudeSumeqzeroofnotcharsAgree}{Proposition~\ref{p:incoh}}
\cs{hasCommonFlagVectorofle}{Lemma~\ref{l:flag}}
\cs{existshasLevel}{Proposition~\ref{p:level}}
\cs{energyleoflevelpairs}{Theorem D}
\cs{energyleoflevellt}{Theorem D}
\cs{energyleoflinkEnergyle}{Theorem E}

\begin{document}
\color{ink}

{\raggedright
{\sffamily\bfseries\fontsize{7.2}{9}\selectfont\color{leancol}\addfontfeatures{LetterSpace=8}READBACK NOTE 43\par}
\vspace{5pt}
{\fontsize{21.5}{25.5}\selectfont\bfseries The energy of amplitude sums: Schur's test on thin classes, block bounds, incoherent and degenerate pairs\par}
\vspace{5pt}
{\fontsize{12.4}{16}\selectfont\itshape\color{inksoft} One unit of the Lean artifact \ArtifactName, read back into
mathematics\par}}
\vspace{8pt}
{\color{rulegrey}\hrule height 0.6pt}
\vspace{7pt}

\textbf{Summary.} In the proof of Theorem~B of the write-up (enough forms), the second moment of translates of
Poincar\'e series unfolds into sums of \emph{amplitudes} over rational elements of a unitary group, collected in
blocks \m{\ampsum{\sigma}{S}{\Xi}{c}{s}{c'}}. This note bounds the \emph{energy} of such blocks, the sum of their squared norms: a Schur
test bounds the total amplitude of one double coset of the arithmetic unipotent group (Theorem~A); summing over the
cosets gives bounds for blocks (Theorem~B); passing from the Bessel functions to their cuts costs little
(Theorem~C); pairs of terms with different links or with degenerate frames are negligible, by lattice counts
(Theorem~D); and the result is the crude estimate (Theorem~E) that the proof of the filling theorem uses.

\statement{leancol}{\fontsize{9.6}{13.2}\selectfont\textbf{What this note is.} The artifact \ArtifactName{} is a
body of Lean code whose final theorems state the Riemann hypothesis and its analogue for Dirichlet $L$-functions.
That is a claim: its authors report that Lean's kernel accepts the code, a replay of that check by the owner of
this write-up is recorded elsewhere, and nothing was compiled for this note; no mathematician has read the proof
through. This note treats one part of one unit (the folder of the filling theorem): \UnitInNote{} of its
\UnitMembers{} declarations, those of four files on the energy of amplitude sums. Every statement was compared,
argument by argument, with the Lean text. The proofs are complete, with two exceptions that are said where they
occur: results of the other notes of the unit, of the unit on the real place and of the background library are
quoted as \emph{Inputs}, with their owners, and the comparisons of exponents that the formal text leaves to a
decision procedure are carried out by hand with cruder intermediate bounds (Section~\ref{s:notes}). The Lean text read is \PinText;
Appendix~A is the concordance of \ConcDecls{} declarations.}

\sect{The statements}{s:intro}
The objects are defined in Section~\ref{s:setting}; most of them belong to the set-up of the unit, which is
the subject of Note~S (\emph{The filling data: Bessel subgroups, links of frames, Siegel's lemma and strong
approximation}) and of Note~B (\emph{Box partitions of arithmetic unipotent groups: frame counts, the covolume,
amplitude sums}). In a few words: $S$ is the data of the filling theorem over a CM field~$E$; \m{\GS{S}{E}} is the
unitary group of a hermitian space with \m{\sro} hyperbolic planes; \m{\Cl{S}} is a finite set of classes
and \m{\GRes{S}} a finite group; a rational element \m{\gamma\in\GS{S}{E}} lies over \m{(c,s,c')} if it carries the
lattice of the class $c'$ to that of $c$ and reduces to~$s$; its \emph{amplitude} \m{\amp{\sigma}{S}{c}{c'}{\gamma}} is
an integral of a product of two Bessel functions, full (\m{\sigma=\full}) or cut (\m{\sigma=\cut}), against
cutoffs for the arithmetic unipotent groups \m{\clU{S}{c}}, \m{\clU{S}{c'}}; and
\m{\ampsum{\sigma}{S}{\Xi}{c}{s}{c'}} is the sum of the amplitudes of the elements of a set $\Xi$ that lie over
\m{(c,s,c')}. The two standing hypotheses are collected in Setting~\ref{set:data}, those of the energy estimates
in Setting~\ref{set:en}.

\statement{arith}{\phantomsection\label{t:A}\textbf{Theorem A (Schur's test on a thin class).}\enspace
In Setting~\ref{set:data}, let \m{j\in\{0,\dots,\sro+1\}}, let \m{c,c'\in\Cl{S}}, \m{s\in\GRes{S}}, and let
\m{\gamma\in\GS{S}{E}} be such that \m{\gamma\in\terms{S}{c}{s}{c'}} and \m{\HasLev{S}{j}{\gamma}}. Let
\m{\mathcal D=\clUo{S}{c}\,\gamma\,\clUo{S}{c'}} be its double coset.

\textup{(a)} With \m{j_0=\min(\sro,j)},
\[\sum_{\delta\in\mathcal D}\schi{S}{c}{c'}{\delta}\ \le\ \nuS\cdot\#\Bc{S}\cdot\tl^{2}\cdot
\bigl(\bb^{4}\,\gq\bigr)^{-j_0\,\sk\,\fE}\cdot
\Bigl(\frac{\xri{S}{j}{\pos{S}{c}{c'}{\gamma}}}{\sqrt{\mnorm{S}{j}{\gamma}}}\Bigr)^{\sk+1}.\]

\textup{(b)} If moreover \m{\tdnorm{S}{\gamma}\le\Nn^{\fE}\,\mnorm{S}{j}{\gamma}}, then
\[\sum_{\delta\in\mathcal D}\ampabs{\cut}{S}{c}{c'}{\delta}\ \le\ \nuS\cdot\csize{S}{j}{\cgauge{S}{j}{c}{c'}{\gamma}}.\]}

\statement{arith}{\phantomsection\label{t:B}\textbf{Theorem B (block bounds).}\enspace
In Setting~\ref{set:data}, let \m{\Xi\subseteq\GS{S}{E}} be stable, \m{\IsStab{S}{\Xi}}, and let \m{c,c'\in\Cl{S}}
and \m{s\in\GRes{S}}.

\textup{(a)} If \m{\tdisc{S}{\gamma}\ne0} for all \m{\gamma\in\Xi}, then
\[\bigl\|\ampsum{\cut}{S}{\Xi}{c}{s}{c'}\bigr\|\ \le\
\sum_{\substack{\gamma\in\Xi\cap\terms{S}{c}{s}{c'}\\ \IsNRep{S}{c}{c'}{\gamma}}}
\nuS\cdot\csize{S}{\sro+1}{\tdnorm{S}{\gamma}}.\]

\textup{(b)} If \m{j\in\{0,\dots,\sro+1\}} and \m{\HasLev{S}{j}{\gamma}} for all \m{\gamma\in\Xi}, then
\m{\ampsum{\cut}{S}{\Xi}{c}{s}{c'}=0} unless \m{\Agr{S}{s}}, and if \m{\Agr{S}{s}} then
\[\bigl\|\ampsum{\cut}{S}{\Xi}{c}{s}{c'}\bigr\|\ \le\
\sum_{\substack{\gamma\in\Xi\cap\terms{S}{c}{s}{c'}\\ \IsNRep{S}{c}{c'}{\gamma}}}
\nuS\cdot\csize{S}{j}{\cgauge{S}{j}{c}{c'}{\gamma}}.\]}

\statement{arith}{\phantomsection\label{t:C}\textbf{Theorem C (full against cut).}\enspace
In Setting~\ref{set:data}, for every set \m{\Xi\subseteq\GS{S}{E}} and every set \m{\Omega\subseteq\GRes{S}},
\[\energy{S}{\ampsum{\full}{S}{\Xi}{\cdot}{\cdot}{\cdot}-\ampsum{\cut}{S}{\Xi}{\cdot}{\cdot}{\cdot}}{\Omega}
\ \le\ \bb^{-30}\cdot\dd.\]}

\statement{arith}{\phantomsection\label{t:D}\textbf{Theorem D (incoherent and degenerate pairs are
negligible).}\enspace In Setting~\ref{set:en}, with its $C$, $\Omega$, $\Xi$ and $X$:

\textup{(a)} Let \m{w\ge0} be a real number, \m{j\in\{0,\dots,\sro+1\}}, and $R$ a set of $E$-subspaces of
\m{(E^{\Mx{\fs}})^2}. Suppose that \m{\HasLev{S}{j}{\gamma}} for all \m{\gamma\in\Xi}; that
\m{\tl^{12}X^{-\sk}\,\dm\le C} in case \m{\sro<j} and some subspace does not belong to~$R$; and that
\m{\linkE{\cut}{S}{\Xi}{\Omega}{R}\le w\,\nuS^{2}}. Then
\[\energy{S}{\ampsum{\cut}{S}{\Xi}{\cdot}{\cdot}{\cdot}}{\Omega}\ \le\ w+\frac{C}{\bb^{30}\cdot4^{\,\sro+1-j}}.\]

\textup{(b)} Let \m{l\in\{0,\dots,\sro+1\}} be such that every \m{\gamma\in\Xi} with \m{\ChA{S}{\gamma}} satisfies
\m{\HasLev{S}{j}{\gamma}} for some \m{j<l}. Then
\[\energy{S}{\ampsum{\cut}{S}{\Xi}{\cdot}{\cdot}{\cdot}}{\Omega}\ \le\ \frac{C}{\bb^{30}\cdot4^{\,\sro+1-l}}.\]}

\statement{arith}{\phantomsection\label{t:E}\textbf{Theorem E (the crude estimate).}\enspace
In Setting~\ref{set:en}, with its $C$, $\Omega$, $\Xi$ and $X$, let \m{w\ge0} be a real number and $R$ a set of
$E$-subspaces of \m{(E^{\Mx{\fs}})^2}. Suppose that \m{\tl^{12}X^{-\sk}\,\dm\le C} in case some subspace does not
belong to~$R$, and that
\[\linkE{\cut}{S}{\Xi\cap\{\gamma:\tdisc{S}{\gamma}\ne0\}}{\Omega}{R}\ \le\ w\,\nuS^{2}.\]
Then
\[\energy{S}{\ampsum{\full}{S}{\Xi}{\cdot}{\cdot}{\cdot}}{\Omega}\ \le\ \bigl(1+\bb^{-4}\bigr)^{2}\,w+\bb^{-12}\,C.\]}

Theorem~E is what the other notes of the unit use: the note on the filling theorem applies it twice, in its
estimate of the main term and, with $R$ empty and \m{w=0}, for the terms outside the Bessel subgroup. Theorems A to D are its proof: E follows from C, D\,(a) at the top level \m{j=\sro+1} and
D\,(b); D\,(a) follows from B\,(b), from the inequality of pair invariants of Note~P (Input~\inp{P2}) and from the
count of pairs (Proposition~\ref{p:pairs}); B follows from A and from the open-cell bound of the real place
(Input~\inp{R3}).

\sect{The setting}{s:setting}
\begin{conventions}\phantomsection\label{c:conventions}
The natural numbers are \m{0,1,2,\dots}; for natural numbers \m{j_0\le j} the difference \m{j-j_0} is a natural
number. $E$ is a CM field with its CM type: a number field that is a CM field, with complex conjugation
\m{x\mapsto\cj x}, together with the choice of one embedding at each infinite place~$\tau$; \m{x_\tau\in\CC} is
the image of \m{x\in E} under it, and \m{\fE} is the number of infinite places, so that \m{[E:\Q]=2\,\fE}.
\m{\Nm} is the norm, \m{\rdE{E}} the root discriminant and \m{\dE{E}} the discriminant of~$E$; \m{\AEf} is the
ring of finite adeles, \m{\hO} its subring of integral adeles, and \m{\psi_{E,f}} the standard additive character
of~\m{\AEf}. For \m{a\in\AEf}, \m{\gix{a}} is the \emph{generalized index}: for a unit~$a$, the inverse of the factor
by which multiplication by $a$ scales a Haar measure of \m{\AEf}, and $0$ for a non-unit; so
\m{\gix{a}=[\hO:a\hO]} for a unit \m{a\in\hO}, and \m{\gix{x}=|\Nm(x)|=\prod_\tau|x_\tau|^2} for \m{x\in E}. For subgroups $M$, $M'$
of a group, \m{[M':M]} is the index of \m{M\cap M'} in~$M'$, read as $0$ when it is infinite. A matrix over
\m{\AEf} is \emph{integral} if its entries lie in~\m{\hO}; for a matrix $M$ with rows indexed by a finite set~$I_1$,
\m{\hO\langle M\rangle\subseteq\AEf^{I_1}} is the \m{\hO}-module spanned by its columns. The field $E$ is \emph{stably unramified} if it
contains a primitive root of unity of order~$5040$. A \emph{shape} is a triple \m{\fs=(\varrho,\kappa,k_0)} of a natural
number, a finite set and an element of it; we write \m{\sro}, \m{\sk=\#\kappa}, \m{\sn=\sro+1} and
\m{\rr{\fs}=2\sro+\sk+1}. For a commutative ring with involution, $M^*$ is the conjugate transpose of a
matrix~$M$. $S$ is filling data over $E$ of shape $\fs$ (Setting~\ref{set:S}); it is \emph{admissible} if it
satisfies the standing hypotheses of Note~S, of which Setting~\ref{set:data} lists those used here. A \emph{side} $\sigma$ is one of the two
words \m{\full}, \m{\cut}. $V$ is a normed additive group. Sums of non-negative terms over arbitrary index sets and
integrals of functions with values in \m{[0,\infty]} are taken in \m{[0,\infty]}; the norm of a vector is also
read there.
\end{conventions}

\begin{setting}[filling data; from Note~S]\label{set:S}
$E$ is a CM field with its CM type, $\fs$ a shape, and $S$ is filling data over $E$ of shape~$\fs$,
\m{S\in\Fill{E}{\fs}}. Of its components and conditions this note uses:
\begin{enumerate}[label=\textup{(\roman*)},leftmargin=2.3em,itemsep=2pt]
\item a real number \m{\bb\ge1000}; a prime number \m{\Nn\ge5} and a finite idele \m{\piS{S}\in\hO} with
\m{\piS{S}\,\cj{\piS{S}}=\Nn}; two prime numbers \m{\qS{S}{0}}, \m{\qS{S}{1}} with
\m{\bb^{10}\le\qS{S}{i}\le\bb^{11}}, and \m{\gq=\qS{S}{0}\,\qS{S}{1}}, so that \m{\bb^{19}\le\gq\le\bb^{23}};
\item for \m{i=0,1} a set \m{\ChS{S}{i}} of \m{\fE} finite places of~$E$, the \emph{chosen places} above
\m{\qS{S}{i}}, whose residue fields have \m{\qS{S}{i}} elements;
\item a generator \m{\vpS{S}} of the different of~$E$ with \m{\cj{\vpS{S}}=-\vpS{S}}; an open compact subgroup
\m{\KcolS{S}} of the unitary group \m{\Gf{E}{\kappa}} of the standard hermitian form on \m{\AEf^\kappa} (the
\emph{column level}), all of whose class automorphism groups are trivial; and for every infinite place $\tau$ a
Bessel system \m{\BsysS{S}{\tau}} for \m{(\fs,\bb,\fE)}, described below;
\item the inequalities \m{10^5\le\sro}, \m{1000\,(2\sro+2)^9\le\sk}, \m{\sk\le\fE\le\bb/1000}, \m{\sk\le\bb},
\m{\sn\le\bb} and \m{(200\,\fE\,\sk\,\rdE{E})^{1000}\le\bb}.
\end{enumerate}
\end{setting}

\subhead{Forms and groups.} These notions belong to the background library and to Note~S. The index set
\m{\Ix{\fs}} consists of \m{e_0,\dots,e_{\sro-1}}, of the elements of~$\kappa$, of~$z$, and of
\m{e^*_0,\dots,e^*_{\sro-1}}; \m{\Mx{\fs}=\kappa\sqcup\{z\}} is the set of the \emph{middle} indices. The \emph{frame} is
\m{\fr{a}=e_a} for \m{a<\sro} and \m{\fr{\sro}=z}. For a commutative ring $A$ with involution, an
\emph{$\fs$-form} over $A$ is a matrix $T$ indexed by \m{\Mx{\fs}}; its \emph{split Gram matrix} \m{\JT{T}},
indexed by \m{\Ix{\fs}}, pairs $e_a$ with $e^*_a$ to $1$ and is $T$ on the middle; \m{\GT{T}} is the unitary group
of \m{\JT{T}} over~$A$, acting on the column vectors \m{t\in A^{\Ix{\fs}}}, and \m{h_T(t,t')=t^*\JT{T}\,t'}. For \m{\gamma\in\GT{T}} we write
\m{\gamma\mathfrak f} for the matrix with the columns \m{\gamma\fr{0},\dots,\gamma\fr{\sro}}, \m{\mathfrak f} for
this matrix at \m{\gamma=1}, and \m{\gamma\mathfrak f\,v=\sum_a v_a\,\gamma\fr{a}} for \m{v\in A^{\sn}}. The
\emph{frame Gram matrix} is \m{\ZT{T}{\gamma}=(h_T(\fr{a},\gamma\fr{a'}))_{a,a'\le\sro}}. For a square matrix $Z$
of size \m{\sn} and \m{l\le\sn}, \m{\lmin{l}{Z}} is its leading principal minor of order~$l$, and for \m{N'\in A}
\[\fdisc{N'}{Z}=\det Z\cdot\cj{\det Z}-N'^{\,2}\,\lmin{\sro}{Z}\,\cj{\lmin{\sro}{Z}}.\]
In \m{\GT{T}}: \m{\UT{T}} is the \emph{Bessel unipotent subgroup}, the elements that are upper triangular with
identity diagonal blocks for the order \m{e_0,\dots,e_{\sro-1}}, middle, \m{e^*_{\sro-1},\dots,e^*_0}; \m{\UoT{T}} is
its subgroup of the elements that fix every frame vector; \m{\HT{T}} is the subgroup of the elements that fix all
$e_a$, all $e^*_a$ and~$z$; and \m{\ST{T}=\UT{T}\,\HT{T}} is the \emph{Bessel subgroup}. For \m{u\in\ST{T}},
\m{\sd{u}=\sum_{a<\sro}u_{\fr{a},\fr{a+1}}} is the sum of its entries above the diagonal of the frame.

Over \m{\CC} we use the form \m{\Tone}, diagonal with entries $1$ on $\kappa$ and $-1$ at~$z$, and
\m{\Gone{\fs}=\GT{\Tone}}, with the Haar measure that the library fixes on it; \m{\Zr{y}=\ZT{\Tone}{y}} for
\m{y\in\Gone{\fs}}, and \m{\gdr{Z}=\ReP\fdisc{1}{Z}=|\det Z|^2-|\lmin{\sro}{Z}|^2}. For \m{y\in\Gone{\fs}},
\m{\bsize{y}=\inf_u\|u\,y\|+1}, the infimum over \m{u\in\ST{\Tone}} of operator norms; for a real number~$b_0$,
\m{\InCut{y}{b_0}} means \m{\bsize{y}\le b_0^{\,6\sn+2}}; and for a real number $\rho$, \m{\frbox{\rho}} is the set
of the \m{y\in\Gone{\fs}} with \m{\|y^{-1}\fr{a}\|\le\rho} for all \m{a\le\sro} (Euclidean norms). The \emph{cell mass} of
\m{y\in\Gone{\fs}} for real numbers $\deg_0$, $\rho$, $d_0$ is
\[\cellm{y}{\deg_0}{\rho}{d_0}=\Bigl(\frac{\rho^{\,9\sn}}{\gdr{\Zr{y}}}\Bigr)^{10^5\sn^4}\cdot
\frac{d_0^{\,2}}{\deg_0\cdot\occ{\fs}\cdot\bigl(\gdr{\Zr{y}}/2\bigr)^{\sk-1}},\]
with a positive constant \m{\occ{\fs}} of the shape. We write \m{\Ginf=\prod_\tau\Gone{\fs}}, with the product
measure; it depends on $E$ and~$\fs$.

\subhead{The groups of the data.} In Setting~\ref{set:S}, for a commutative ring $A$, \m{\Trat{S}{A}} is the
$\fs$-form over $A$ that is diagonal with $1$ on $\kappa$ and \m{-\Nn} at~$z$, and \m{\GS{S}{A}=\GT{\Trat{S}{A}}}
when $A$ has an involution. For \m{\gamma\in\GS{S}{E}}:
\m{\Zf{S}{\gamma}=\ZT{\TE}{\gamma}}, \m{\tdisc{S}{\gamma}=\fdisc{\Nn}{\Zf{S}{\gamma}}\in E}, and
\m{\tdnorm{S}{\gamma}=\prod_\tau|\tdisc{S}{\gamma}_\tau|}. \m{\Tlat{S}{A}} is the Gram matrix of the same hermitian
space in the \emph{lattice basis}, in which the basis vector of index $k_0$ is replaced by an isotropic vector;
\m{\GFin{S}=\GT{\Tl}}, \m{\KS{S}\subseteq\GFin{S}} is the subgroup of the integral matrices with integral inverse,
\m{\GRes{S}=\prod_{i}\prod_{w\in\ChS{S}{i}}\GL_{\Ix{\fs}}(\F_w)} with the residue fields \m{\F_w}, and
\m{\redS{S}\colon\KS{S}\to\GRes{S}} is the reduction of the entries; its kernel is \m{\KgS{S}}, and
\m{\Sbar} is the image of \m{\KS{S}\cap\ST{\Tl}}, a subgroup of \m{\GRes{S}}. There are homomorphisms
\m{\rfin{S}\colon\GS{S}{E}\to\GFin{S}} (the diagonal embedding, written in the lattice basis) and
\m{\rinf{S}\colon\GS{S}{E}\to\Ginf}, where \m{\ri{\gamma}_\tau} is the matrix of~$\gamma$ at~$\tau$ conjugated by
the diagonal matrix of the \emph{scales} \m{\bscale{S}{\tau}{\cdot}}:
\m{\bscale{S}{\tau}{e_a}=\kapS{S}{\tau}^{\sro-a}/\sqrt{\Nn}}, \m{\bscale{S}{\tau}{i}=1} on $\kappa$,
\m{\bscale{S}{\tau}{z}=1/\sqrt{\Nn}}, \m{\bscale{S}{\tau}{e^*_a}=\bscale{S}{\tau}{e_a}^{-1}}. Here
\m{\kapS{S}{\tau}} is a non-zero real number, the real part of \m{\gq\,\vpS{S}_\tau/(2\pi i)}, and
\m{\fscale{S}{\tau}{a}=\kapS{S}{\tau}^{2(\sro-a)}/\Nn} for \m{a\le\sro}.

\subhead{Classes and terms.} \m{\Cl{S}} is the finite set of the double cosets of \m{\Gf{E}{\kappa}} by the
rational points of the unitary group on the left and \m{\KcolS{S}} on the right; \m{\incH{S}} embeds
\m{\Gf{E}{\kappa}} in \m{\HT{\Tl}\subseteq\GFin{S}}, and \m{\hc{S}{c}} is the image of a fixed representative of the class~$c$. For
\m{c,c'\in\Cl{S}} and \m{\gamma\in\GS{S}{E}},
\[\pos{S}{c}{c'}{\gamma}=\hc{S}{c}^{-1}\,\rfi{\gamma}\,\hc{S}{c'}\in\GFin{S},\]
and for \m{s\in\GRes{S}}, \m{\terms{S}{c}{s}{c'}} is the set of the $\gamma$ with \m{\pos{S}{c}{c'}{\gamma}\in\KS{S}}
and \m{\redS{S}(\pos{S}{c}{c'}{\gamma})=s}: the \emph{terms over} \m{(c,s,c')}. The arithmetic unipotent groups are
\m{\clU{S}{c}=\{u\in\UE:\pos{S}{c}{c}{u}\in\KgS{S}\}} and \m{\clUo{S}{c}=\clU{S}{c}\cap\UoE}. The matrix
\m{\clC{S}{c}} over \m{\AEf} takes rational coordinates to the coordinates in the lattice basis of the class~$c$,
so that the matrix of \m{\pos{S}{c}{c'}{\gamma}} is \m{\clC{S}{c}\,\fadm{\gamma}\,\clC{S}{c'}^{-1}}, where
\m{\fadm{M}} is a matrix $M$ over $E$ read over~\m{\AEf}; \m{\clCX{S}{c}} is its part on~$\kappa$, a matrix with
rows in \m{\Mx{\fs}} and columns in~$\kappa$.

\subhead{The real places.} The Bessel system \m{\BsysS{S}{\tau}} consists of a non-empty finite set
\m{\It{S}{\tau}}, a number \m{\degt{S}{\tau}>0} and two continuous functions
\m{\Wt{\full}{S}{\tau},\Wt{\cut}{S}{\tau}\colon\Gone{\fs}\to\CC^{\It{S}{\tau}}}, the \emph{full} and the \emph{cut
Bessel function}, with further properties of which this note uses three. \inp{W1}~For each side~$\sigma$ there is
a character $\chi$ of \m{\Uone} with values of absolute value~$1$ such that \m{\Wt{\sigma}{S}{\tau}(u\,y)=\chi(u)\,\Wt{\sigma}{S}{\tau}(y)} for \m{u\in\Uone}.
\inp{W2}~\m{\int\|\Wt{\cut}{S}{\tau}\|^2\le\#\It{S}{\tau}}, the integral over the quotient of \m{\Gone{\fs}} by
\m{\Uone}. \inp{W3}~If \m{\Wt{\cut}{S}{\tau}(y)\ne0} then \m{\InCut{y}{\bb}}. We put \m{\Bc{S}=\prod_\tau\It{S}{\tau}}
and \m{\bdeg{S}=\prod_\tau\degt{S}{\tau}}. For a side~$\sigma$, \m{y\in\Ginf} and \m{\gamma\in\GS{S}{E}}, the
vector \m{\besker{\sigma}{S}{\gamma}{y}\in\CC^{\Bc{S}\times\Bc{S}}} is the tensor product over $\tau$ of the
vectors \m{\Wt{\sigma}{S}{\tau}(y_\tau)} and of the complex conjugates of the vectors
\m{\Wt{\sigma}{S}{\tau}((\ri{\gamma}^{-1}y)_\tau)}; its norm is the product of their norms.

The \emph{box size} is \m{\bxs{S}=\bb^{2\sn^2}}, and \m{\bxp{S}} is the set of the \emph{boxed points}
\m{y\in\Ginf} with \m{y_\tau\in\frbox{\bxs{S}}} and \m{\InCut{y_\tau}{\bb}} for all~$\tau$. An element
\m{\gamma\in\GS{S}{E}} is \emph{near}, \m{\IsNr{S}{\gamma}}, if there is \m{y\in\bxp{S}} with
\m{\ri{\gamma}^{-1}y\in\bxp{S}}; and \m{\IsNRep{S}{c}{c'}{\gamma}} means that $\gamma$ is near and is the
representative, fixed once and for all, of its double coset \m{\clUo{S}{c}\,\gamma\,\clUo{S}{c'}}. For a class $c$,
\m{\cutoff{S}{c}\colon\Ginf\to[0,\infty]} is the \emph{cutoff} of Note~B, a product
\m{\cutoff{S}{c}=\theta_{S,c}\cdot\omega_{S,c}} of two functions described in Inputs \inp{S9}
and~\inp{B1}. The \emph{amplitude} and the \emph{absolute amplitude} of $\gamma$ over \m{c,c'} are
\begin{gather*}
\amp{\sigma}{S}{c}{c'}{\gamma}=\int_{\Ginf}\cutoff{S}{c}(y)\,\cutoff{S}{c'}(\ri{\gamma}^{-1}y)\;\besker{\sigma}{S}{\gamma}{y}\,dy,\\
\ampabs{\sigma}{S}{c}{c'}{\gamma}=\int_{\Ginf}\cutoff{S}{c}(y)\,\cutoff{S}{c'}(\ri{\gamma}^{-1}y)\;\bigl\|\besker{\sigma}{S}{\gamma}{y}\bigr\|\,dy,
\end{gather*}
and for a set \m{\Xi\subseteq\GS{S}{E}} the \emph{amplitude sum} is
\[\ampsum{\sigma}{S}{\Xi}{c}{s}{c'}=\sum_{\gamma\in\Xi\cap\terms{S}{c}{s}{c'}}\amp{\sigma}{S}{c}{c'}{\gamma}\ \in\CC^{\Bc{S}\times\Bc{S}}.\]
The formal text defines the integral and the sum by conventions that give $0$ when they do not converge; under
the hypotheses of Setting~\ref{set:data} they converge (Input~\inp{B4}).

\subhead{Levels, links, characters.} Let \m{\gamma\in\GS{S}{E}}. For \m{a\le\sro} and \m{c_1,c_2\in E},
\m{\HCFV{\gamma}{c_1}{c_2}{a}} means: there are \m{v,v'\in E^{\sn}} with \m{\mathfrak f\,v=\gamma\mathfrak f\,v'},
\m{v_a=c_1}, \m{v'_a=c_2} and \m{v_{a'}=v'_{a'}=0} for \m{a'>a}; the flags of the two frames share a vector with these
leading coefficients. For a natural number~$j$, \m{\IsNF{j}{\gamma}} (\emph{normal form of level~$j$}) means: there
is \m{c_1\in E} with \m{c_1\cj{c_1}=1} such that \m{\HCFV{\gamma}{c_1}{1}{a}} for all \m{a\le\sro} with \m{j\le a}.
For \m{j\le\sro+1}, $\gamma$ \emph{has level~$j$}, \m{\HasLev{S}{j}{\gamma}}, if \m{\IsNF{j}{\gamma}},
\m{\lmin{j}{\Zf{S}{\gamma}}\ne0}, and \m{\sro<j} holds exactly if \m{\tdisc{S}{\gamma}\ne0}. The first two
notions are defined in the same way for an element of \m{\GT{T}}, for any $\fs$-form $T$ over a commutative ring
with involution; they are used below also for \m{\Gone{\fs}}.

The \emph{link} \m{\link{\gamma}} is the $E$-subspace of \m{(E^{\Mx{\fs}})^2} of the pairs formed by the middle
coordinates of $t$ and of \m{\gamma^{-1}t}, for the vectors \m{t\in E^{\Ix{\fs}}} orthogonal to all $e_a$ and all
\m{\gamma e_a}, \m{a<\sro}. For \m{\gamma,\gamma'}, the \emph{common link} \m{\clink{\gamma}{\gamma'}} is the
subspace of \m{E^\kappa} of the vectors $t_0$ such that some pair in \m{\link{\gamma}\cap\link{\gamma'}}, both of whose
members have $z$-coordinate~$0$, has first member~$t_0$ on~$\kappa$.

For \m{u\in\GFin{S}} put \m{\uchar{S}{u}=\psi_{E,f}\bigl(-\sd{u}/(\gq\,\vpS{S})\bigr)}; for
\m{u\in\UE} this gives \m{\uchar{S}{\rfi{u}}=\psi_{E,f}(-\sd{u}/(\gq\,\vpS{S}))}. \m{\ChA{S}{\gamma}} (\emph{the
characters agree at~$\gamma$}) means that \m{\uchar{S}{\rfi{u}}=\uchar{S}{\rfi{u'}}} for all \m{u,u'\in\UE} with
\m{u\gamma=\gamma u'}. An element \m{s\in\GRes{S}} \emph{agrees}, \m{\Agr{S}{s}}, if \m{\uchar{S}{u}=\uchar{S}{u'}}
for all \m{u,u'\in\KS{S}\cap\UT{\Tl}} with \m{\redS{S}(u)\,s=s\,\redS{S}(u')}. A set \m{\Xi\subseteq\GS{S}{E}} is
\emph{stable}, \m{\IsStab{S}{\Xi}}, if \m{u^{-1}\gamma u'\in\Xi} is equivalent to \m{\gamma\in\Xi} for all
\m{u,u'\in\UE}. For \m{s\in\GRes{S}} and a natural number~$X$, \m{\FixN{s}{X}} means that there is a finite set
of chosen places and of conjugates of chosen places, with product of residue cardinalities at least~$X$, at each
of which $s$, or its inverse adjoint at a conjugate, fixes a non-zero vector of the span of the frame
(Note~S).

\subhead{Constants.} With \m{\cdim{S}=\#\Cl{S}\cdot\#\Bc{S}}:
\begin{gather*}
\tl=\bb^{\,\sk\fE/4},\qquad \codimF{S}=(\bb/\gq)^{\sk\fE}\cdot2\,\tl^2,\\
\dd=\#\Sbar\cdot\cdim{S},\qquad
\dm=\frac{\cdim{S}^2}{\bdeg{S}\cdot\mass{E}{\rr{\fs}}},\\
\cw{S}{i}=\frac{\#\Sbar\cdot\#\Cl{S}}{\Nn^{\sk\fE}}\,(2\,\tl^2)^{i}\quad(i\in\NN),
\end{gather*}
where \m{\mass{E}{\rr{\fs}}>0} is the mass of the genus of the standard lattice of rank \m{\rr{\fs}}; and
\m{\nuS>0} is the \emph{covolume} of Note~B, given there by a formula in $\gq$, $\Nn$, \m{\dE{E}} and the
\m{\kapS{S}{\tau}}. The names in the formal text are: tail scale, codimension factor, Bessel energy, main energy,
count weight. Three more constants occur in one inequality of admissibility:
\begin{gather*}
\fsprod{S}=\prod_\tau\prod_{a\le\sro}\fscale{S}{\tau}{a}^{\,\sk+2(\sro-a)},\\
\siegC{S}=\Bigl(\frac{\pi^{\sk}}{(\sk-1)!\,\Nn^{\sk}}\Bigr)^{\fE}|\dE{E}|^{1/4}\,\volO{E}^{\sk},\\
\packC{S}=\#\Sbar\cdot\gq^{\,\fE\sro(2\sk+2\sro+1)}\cdot\#\Cl{S}\cdot\siegC{S},
\end{gather*}
with \m{\volO{E}=(2/\rdE{E})^{\fE}}, and \m{|\dE{E}|=\rdE{E}^{2\fE}}.

\begin{setting}[standing hypotheses]\label{set:data}
In Setting~\ref{set:S}, the field $E$ is stably unramified and the filling data $S$ is admissible. Of the
conditions of admissibility (Note~S) this note uses, for all infinite places~$\tau$:
\begin{enumerate}[label=\textup{(A\arabic*)},leftmargin=2.8em,itemsep=2pt]
\item \m{\KcolS{S}} is the preimage of \m{\KgS{S}} under \m{\incH{S}};
\item \m{1\le|\kapS{S}{\tau}|\le\bb^{23}}, and \m{\Nn\le\bb^{60\sn}};
\item \m{\tl^{44}\,\dd\le\dm}, \m{\ \tl^{18}\,\codimF{S}\,\dm\le\dd\ } and \m{\ \codimF{S}\le1};
\item \m{\#\Bc{S}\le\bb^{2\sk\fE}};
\item \m{(\Nn^{\sk}/\bb)^{\fE}\le[\Kf{E}{\kappa}:\incH{S}^{-1}(\KS{S})]}, where \m{\Kf{E}{\kappa}} is the group of
the integral elements of \m{\Gf{E}{\kappa}} with integral inverse;
\item \m{\packC{S}\le(1+\bb^{-9})\cdot\dm\cdot\bdeg{S}\cdot\occ{\fs}^{\fE}\cdot\volO{E}^{-\sn}\cdot
\bigl(\nuS\,\fsprod{S}/\#\Bc{S}\bigr)^{2}}.
\end{enumerate}
\end{setting}

\subhead{Lattices and weights (from Note~B).} Let \m{\mathsf k\in\GFin{S}}. Write \m{\Zl{\mathsf k}} for its frame
Gram matrix, \m{\xrows{Z}} for the matrix of the first \m{\sro} rows of a matrix~$Z$ with \m{\sn} rows, and, for
\m{j\le\sro+1}, \m{\IdxB{j}} for the set of the \m{a<\sro} with \m{a<j}. Let
\m{\mathcal R(\mathsf k)=\hO\langle(\xrows{\Zl{\mathsf k}})^*\rangle} be the lattice spanned by the conjugates of
the first \m{\sro} rows of \m{\Zl{\mathsf k}}, and \m{\mathcal R_j(\mathsf k)} the sublattice spanned by the conjugates of the
rows of index in~\m{\IdxB{j}}. Then
\[\xri{S}{j}{\mathsf k}=[\mathcal R(\mathsf k):\mathcal R_j(\mathsf k)],\qquad
\sati{S}{\mathsf k}=\bigl[\hO^{\Ix{\fs}}\cap\operatorname{span}_{\AEf}(e_a,\mathsf ke_a):\hO\langle e_a,\mathsf ke_a\rangle\bigr],\]
the second with the vectors \m{e_a}, \m{\mathsf ke_a}, \m{a<\sro} (the \emph{saturation index}). For
\m{\gamma\in\GS{S}{E}}, classes \m{c,c'} and \m{j\le\sro+1}, with \m{j_1=j-\min(\sro,j)\in\{0,1\}} and
\m{\mathsf k=\pos{S}{c}{c'}{\gamma}},
\[\mnorm{S}{j}{\gamma}=\frac{|\Nm(\lmin{j}{\Zf{S}{\gamma}})|}{\Nn^{2j_1\fE}},\qquad
\cgauge{S}{j}{c}{c'}{\gamma}=\frac{\max\bigl(\tdnorm{S}{\gamma},\ \Nn^{\fE}\,\mnorm{S}{j}{\gamma}\bigr)}
{\xri{S}{j}{\mathsf k}\cdot\xri{S}{j}{\mathsf k^{-1}}}\]
(the \emph{minor norm} and the \emph{class gauge}). For a class $c$ and an $E$-subspace \m{U\subseteq E^\kappa},
\m{\latd{S}{c}(U)} is the discriminant of the lattice of the points $t_0$ of \m{U\otimes_E\AEf} with
\m{\clCX{S}{c}\,t_0} integral: it is \m{\sqrt{\gix{\det(M_U^*M_U)}}} for a matrix
\m{M_U=\slb{\clCX{S}{c}}{U}} over \m{\AEf}, fixed by Note~S, whose columns form an \m{\hO}-basis of this lattice
(more precisely, $M_U$ is a basis matrix of $U$ over~$E$ times an invertible adelic matrix, and
\m{\clCX{S}{c}\,M_U} is integral with an integral left inverse). For a unit \m{a\in\AEf} the \emph{frame value
index} is
\[\fvi{S}{c}{\mathsf k}{a}{U}=\bigl[\Lambda_U(\mathsf k):a\,\mathcal R(\mathsf k)\bigr],\qquad
\Lambda_U(\mathsf k)=\hO\bigl\langle(\fmid{\mathsf k})^*\ \Tl\ \clCX{S}{c}\,M_U\bigr\rangle,\]
where \m{\fmid{\mathsf k}} is the matrix of the middle rows of \m{\mathsf k\mathfrak f}; so \m{\Lambda_U(\mathsf k)} is
the lattice of the vectors \m{(h(\mathsf k\fr{a},t_0))_{a}} for $t_0$ in the lattice of $U$ and~$c$. The \emph{pair
invariant} of \m{\gamma,\gamma'} is, with \m{U=\clink{\gamma}{\gamma'}} and \m{\mathsf k=\pos{S}{c}{c'}{\gamma}},
\[\pinv{S}{c}{c'}{\gamma}{\gamma'}=\latd{S}{c}(U)\cdot\fvi{S}{c}{\mathsf k}{1}{U}\cdot\sati{S}{\mathsf k}.\]
For a pair \m{p=(c',\gamma)\in\Cl{S}\times\GS{S}{E}} put \m{\telem{S}{c}{p}=\pos{S}{c}{c'}{\gamma}}, and for a function
\m{\eta\colon\GFin{S}\to\NN},
\[\ipw{S}{c}{\eta}{p}=\eta(\telem{S}{c}{p})^{-\sk}\in[0,\infty],\]
read as $0$ when \m{\eta(\telem{S}{c}{p})=0}. For \m{\mathsf t\in\Sbar}, \m{\IsLab{S}{j}{c}{\mathsf t}{p}} (\emph{$p$ is a
labelled term}) means that \m{\IsNRep{S}{c}{c'}{\gamma}}, \m{\HasLev{S}{j}{\gamma}}, and
\m{\gamma\in\terms{S}{c}{s}{c'}} for some $s$ whose component in \m{\Sbar} is~\m{\mathsf t}, for a fixed bijection of
\m{\GRes{S}} with \m{(\GRes{S}/\Sbar)\times\Sbar}. Two elements are \emph{congruent terms},
\m{\Cong{S}{c}{c'}{\gamma}{\gamma'}}, if both lie in \m{\terms{S}{c}{s}{c'}} for one and the same~$s$. Finally, for
\m{i\in\NN}, pairs \m{p,p'} and a subspace~$U$, the \emph{term weight} and the \emph{pair weight} are
\begin{gather*}
\termw{S}{j}{i}{c}{p}{p'}{U}=\cw{S}{i}^{-1}\cdot\latd{S}{c}(U)^{-(\sk+1)}\cdot\ipw{S}{c}{\eta_{j,U}}{p}\cdot
\ipw{S}{c}{\sati{S}{\cdot}}{p'},\\
\eta_{j,U}(\mathsf k)=\gix{\gid{S}}^{\#\IdxB{j}}\cdot\fvi{S}{c}{\mathsf k}{1}{U},\qquad
\gix{\gid{S}}=\gq^{\,2\fE},
\end{gather*}
where \m{\gid{S}} is $\gq$ as a finite idele, and
\[\pairw{S}{j}{c}{c'}{\gamma}{\gamma'}=\termw{S}{j}{\sk-\dim_E U}{c}{(c',\gamma)}{(c',\gamma')}{U},\qquad
U=\clink{\gamma}{\gamma'},\]
if \m{\Cong{S}{c}{c'}{\gamma}{\gamma'}}, \m{\IsNRep{S}{c}{c'}{\gamma}}, \m{\HasLev{S}{j}{\gamma}},
\m{\IsNRep{S}{c}{c'}{\gamma'}} and \m{\HasLev{S}{j}{\gamma'}}; otherwise the pair weight is~$0$. Finally, for a
square matrix $Z$ of size~\m{\sn}, \m{\padl{j}{Z}} is the matrix whose rows of index \m{a<j} are those of $Z$ and
whose other rows are those of the identity matrix; \m{\det\padl{j}{Z}=\lmin{j}{Z}}.

\subhead{The notions of this note.} The remaining definitions are those of the four files treated here.

\begin{definition}[energy, link energy]\label{d:energy}
In Setting~\ref{set:S}: \textup{(a)} for a normed additive group $V$, a family
\m{\mathcal A\colon\Cl{S}\times\GRes{S}\times\Cl{S}\to V} and a set \m{\Omega\subseteq\GRes{S}}, the \emph{energy}
of $\mathcal A$ over $\Omega$ is
\[\energy{S}{\mathcal A}{\Omega}=\sum_{c\in\Cl{S}}\ \sum_{c'\in\Cl{S}}\ \sum_{s\in\Omega}\frac{\|\mathcal A(c,s,c')\|^2}{\nuS^{2}}.\]
\textup{(b)} For a side $\sigma$, a set \m{\Xi\subseteq\GS{S}{E}}, a set \m{\Omega\subseteq\GRes{S}} and a set $R$
of $E$-subspaces of \m{(E^{\Mx{\fs}})^2}, the \emph{link energy} is
\[\linkE{\sigma}{S}{\Xi}{\Omega}{R}=\sum_{c\in\Cl{S}}\ \sum_{c'\in\Cl{S}}\ \sum_{s\in\Omega}\ \sum_{\Lambda\in R}
\bigl\|\ampsum{\sigma}{S}{\Xi\cap\{\gamma:\link{\gamma}=\Lambda\}}{c}{s}{c'}\bigr\|^2\ \in[0,\infty].\]
\end{definition}
\leanline{\lean{energy}\sep\lean{linkEnergy}}

The group \m{\GRes{S}} is finite, so the energy is a finite sum of real numbers. For the family of the amplitude
sums of a set~$\Xi$ we write \m{\ampsum{\sigma}{S}{\Xi}{\cdot}{\cdot}{\cdot}}; here
\m{V=\CC^{\Bc{S}\times\Bc{S}}} with its Euclidean norm.

\begin{definition}[class size]\label{d:cs}
In Setting~\ref{set:S}, for \m{j\in\{0,\dots,\sro+1\}} and a real number \m{x>0},
\[\csize{S}{j}{x}=\tl^{2}\cdot\sqrt{\frac{\mathsf D_j}{\cw{S}{0}\cdot\gq^{\,2\min(\sro,j)\,\sk\,\fE}\cdot x^{\sk+1}}},\]
where \m{\mathsf D_j=\dm} if \m{\sro<j}, and \m{\mathsf D_j=\tl^{13}\,\tl^{-32j}\,\dd} if \m{j\le\sro}.
\end{definition}
\leanline{\lean{classSize}}

\begin{definition}[boxed cutoff, Schur weight, Schur integral]\label{d:schur}
In Setting~\ref{set:S}, let \m{c,c'\in\Cl{S}}, \m{\gamma,\gamma'\in\GS{S}{E}} and \m{y\in\Ginf}. \textup{(a)} The
\emph{boxed cutoff} \m{\bcut{S}{c}\colon\Ginf\to[0,\infty]} is equal to \m{\cutoff{S}{c}} on the set \m{\bxp{S}}
of the boxed points and to $0$ outside it. \textup{(b)} With the frame volume $\Delta$ on \m{\Gone{\fs}}
(Input~\inp{R1}), the \emph{Schur weight} is
\[\schw{S}{c}{\gamma}{y}=\bcut{S}{c}\bigl(\ri{\gamma}^{-1}y\bigr)\cdot
\prod_\tau\Bigl(\frac{\Delta(y_\tau)}{\Delta((\ri{\gamma}^{-1}y)_\tau)}\Bigr)^{\sk+\sn}\in[0,\infty].\]
\textup{(c)} The \emph{Schur integral} is
\[\schi{S}{c}{c'}{\gamma'}=\int_{\Ginf}\bcut{S}{c}(y)\cdot\prod_\tau\bigl\|\Wt{\cut}{S}{\tau}(y_\tau)\bigr\|^{2}\cdot
\schw{S}{c'}{\gamma'}{y}\,dy\ \in[0,\infty].\]
\end{definition}
\leanline{\lean{boxCutoff}\sep\lean{schurWeight}\sep\lean{schurIntegral}}

\begin{setting}[the hypotheses of the energy estimates]\label{set:en}
The hypotheses of Setting~\ref{set:data} hold: $E$ is stably unramified and $S$ is admissible. Further, $C$ is a
real number with \m{0\le C}; \m{\Omega\subseteq\GRes{S}} is a set; \m{\Xi\subseteq\GS{S}{E}} is a stable set,
\m{\IsStab{S}{\Xi}}; $X$ is a natural number with \m{X\ne0} such that \m{\FixN{s}{X}} for every \m{s\in\Omega};
\m{\dd\le C}; and \m{\tl^{20}\cdot\dd\le C} in case there are \m{\gamma\in\Xi}, classes \m{c,c'\in\Cl{S}} and
\m{s\in\Omega} with \m{\gamma\in\terms{S}{c}{s}{c'}} and \m{\IsNF{0}{\gamma}}.
\end{setting}

\sect{Inputs}{s:inputs}
What follows is taken for granted. Each statement was compared with its Lean text; the owners are the two other
notes on the set-up of this unit (Note~S, Note~B), the note on the pairing of links of the same unit (Note~P),
the unit on the real place (unit~132 of the reading, Chapter~VI of the write-up), and the background library.
For the last, the catalogue of the write-up classes all the statements used here as elementary, with one
exception that is said in~\inp{B8}. Throughout, $S$ is as in Setting~\ref{set:S}.

\statement{ana}{\ip\emph{From Note~S.}
\inp{S1} \m{\tl\ge2}, \m{\tl^{4}=\bb^{\sk\fE}}, and \m{\bb^{e}\le\tl} for every natural number
\m{e\le10^5\,\sn^9\,\fE}; \m{\nuS>0}.
\inp{S2} \m{\|\amp{\sigma}{S}{c}{c'}{\gamma}\|\le\ampabs{\sigma}{S}{c}{c'}{\gamma}}.
\inp{S3} \m{\rinf{S}} maps \m{\UE} into the elements $u$ of \m{\Ginf} with all \m{u_\tau\in\Uone}, and \m{\UoE} into those
with all \m{u_\tau\in\Uoone}. For such $u$ with all \m{u_\tau\in\Uoone}: \m{u\,y\in\bxp{S}} exactly if \m{y\in\bxp{S}}. If
\m{\IsNr{S}{\gamma}} and \m{u,u'\in\UoE}, then \m{\IsNr{S}{u\gamma u'}}.
\inp{S4} For \m{\gamma\in\GS{S}{E}}, \m{a,a'\le\sro} and \m{j\le\sn}:
\m{\Zr{\ri{\gamma}_\tau}_{a,a'}=\kapS{S}{\tau}^{\sro-a}\kapS{S}{\tau}^{\sro-a'}\Nn^{-1}\,(\Zf{S}{\gamma}_{a,a'})_\tau};
\ \m{\lmin{j}{\Zr{\ri{\gamma}_\tau}}=\prod_{a<j}\fscale{S}{\tau}{a}\cdot\lmin{j}{\Zf{S}{\gamma}}_\tau};
the number \m{\tdisc{S}{\gamma}_\tau} is real and non-negative; and
\m{\gdr{\Zr{\ri{\gamma}_\tau}}=\bigl(\prod_{a\le\sro}\fscale{S}{\tau}{a}\bigr)^{2}\,|\tdisc{S}{\gamma}_\tau|}.
\inp{S5} If $S$ is admissible and \m{\IsNr{S}{\gamma}}, then \m{|(\Zf{S}{\gamma}_{a,a'})_\tau|\le\bb^{5\sn^2}} and
\m{|\lmin{l}{\Zf{S}{\gamma}}_\tau|^{2}\le\bb^{11\sn^3}} for all \m{a,a'\le\sro}, \m{l\le\sn} and all~$\tau$.
\inp{S6} If \m{\gamma\in\terms{S}{c}{s}{c'}}, then the entries of \m{\Zf{S}{\gamma}} are algebraic integers,
\m{\pos{S}{c}{c'}{\gamma}} is integral, \m{\gamma^{-1}\in\terms{S}{c'}{s^{-1}}{c}}, and
\m{\gamma\delta\in\terms{S}{c}{s\,s'}{c''}} for \m{\delta\in\terms{S}{c'}{s'}{c''}}. For all $\gamma$:
\m{\Zl{\pos{S}{c}{c'}{\gamma}}=\fadm{\Zf{S}{\gamma}}}, and the matrix \m{\pos{S}{c}{c'}{\gamma}\,\mathfrak f} is
\m{\clC{S}{c}\,\fadm{\gamma\mathfrak f}}.
\inp{S7} \m{\KgS{S}} is the set of the \m{\mathsf k\in\KS{S}} such that all entries of \m{\mathsf k-1} lie in
\m{\gq\,\hO}.
\inp{S8} With the representative \m{x_c\in\Gf{E}{\kappa}} of the class~$c$ and the matrix
\m{\mbc{\fs}{\piS{S}}} of Note~S: \m{\clCX{S}{c}=\mbc{\fs}{\piS{S}}\,x_c^{-1}}; for a matrix $M$ over \m{\AEf} with
rows in $\kappa$, \m{\mbc{\fs}{\piS{S}}\,M} is integral exactly if $M$ is integral and its row $k_0$ has entries
in \m{\piS{S}\hO}; and for \m{\mathsf k\in\Gf{E}{\kappa}}, \m{\incH{S}(\mathsf k)\in\KS{S}} exactly if
\m{\mathsf k\in\Kf{E}{\kappa}} and the row $k_0$ of \m{\mathsf k-1} has entries in \m{\piS{S}\hO}. \inp{S9} For every subgroup $\Gamma$ of \m{\clU{S}{c}} there is a measurable \m{\theta\colon\Ginf\to[0,\infty]}
with \m{\sum_{u\in\Gamma}\theta(\ri{u}\,y)=1} for all~$y$; the image of $\Gamma$ in \m{\Ginf} is discrete, and for
every measurable function $F_0$ on \m{\Ginf} with values in a finite-dimensional complex vector space that is invariant under left translation by the image of~$\Gamma$, the integral
\m{\int\theta\,F_0} is the integral of $F_0$ over the quotient of \m{\Ginf} by this image. \m{\theta_{S,c}} is the function
fixed in this way for \m{\Gamma=\clUo{S}{c}}.
\inp{S10} For \m{u,u'\in\UE}: \m{\link{u^{-1}\gamma u'}=\link{\gamma}},
\m{\tdisc{S}{u^{-1}\gamma u'}=\tdisc{S}{\gamma}}, \m{\lmin{l}{\Zf{S}{u^{-1}\gamma u'}}=\lmin{l}{\Zf{S}{\gamma}}},
and \m{\HasLev{S}{j}{u^{-1}\gamma u'}} exactly if \m{\HasLev{S}{j}{\gamma}}. If \m{u'\in\UoE} then
\m{(\gamma u')\mathfrak f=\gamma\mathfrak f}.
\inp{S11} A pair \m{(t_1,t_2)} of vectors of \m{E^{\Mx{\fs}}} lies in \m{\link{\gamma}} exactly if there is a
vector~$t$ with middle coordinates $t_2$ and no $e^*$-coordinates such that \m{\gamma t} has middle coordinates
$t_1$ and no $e^*$-coordinates. The link is Lagrangian for the form \m{\Trat{S}{E}\oplus(-\Trat{S}{E})}: a pair lies in
\m{\link{\gamma}} exactly if it is orthogonal to \m{\link{\gamma}}. If \m{t_1\in E^{\Mx{\fs}}} has $z$-coordinate
$0$ and \m{t_1^*\,\Trat{S}{E}\,t_1=0}, then \m{t_1=0}.
\inp{S12} If \m{\IsNF{j}{\gamma}} then \m{\IsNF{j}{\gamma^{-1}}} and \m{\IsNF{j}{\ri{\gamma}_\tau}} for every~$\tau$;
for an $\fs$-form $T$, \m{u\in\UT{T}} and \m{\gamma_1\in\GT{T}} with \m{\IsNF{j}{\gamma_1}} also
\m{\IsNF{j}{u\gamma_1}}; and for every \m{v_0\in E^{\sn}} there are
\m{v,v'\in E^{\sn}} with \m{\mathfrak f\,v=\gamma\mathfrak f\,v'} and \m{v'_a=(v_0)_a} for all \m{a\ge j}. If
\m{\tdisc{S}{\gamma}\ne0} then \m{\HasLev{S}{\sro+1}{\gamma}}.
\inp{S13} If \m{\mathfrak f\,v=\gamma\mathfrak f\,v'} then \m{\cj{v_{\sro}}\,v_{\sro}=\cj{v'_{\sro}}\,v'_{\sro}} and
\m{\Zf{S}{\gamma}\,v'} is the vector with entry \m{-\Nn\,v_{\sro}} at \m{\sro} and $0$ elsewhere. If
\m{\Zf{S}{\gamma}\,v'} is the vector with entry \m{D_0\in E} at \m{\sro} and $0$ elsewhere, then
\m{t=\gamma\mathfrak f\,v'} has no $e^*$-coordinates, $z$-coordinate \m{-D_0/\Nn}, and for every $\tau$
\[\Nn\sum_{i\in\kappa}|(t_i)_\tau|^{2}=|(D_0)_\tau|^{2}-\bigl(\Nn\,|(v'_{\sro})_\tau|\bigr)^{2};\]
and if this number is \m{\le0} for some $\tau$, then \m{t=\mathfrak f\,v} for some \m{v\in E^{\sn}}.
\inp{S14} (Eichler elements.) Let \m{a\le\sro}, \m{v_0\in E^{\sn}} with \m{(v_0)_{a'}=0} for \m{a'\ge a},
\m{r_0\in E}, and $t$ a vector such that \m{h(t,\fr{a'})=0} for \m{a'<a}, with \m{h=h_{\TE}}. Then the map
\m{t'\mapsto t'+t_0\,h(t,t')-(t+\tfrac12h(t,t)\,t_0)\,h(t_0,t')}, with \m{t_0=r_0\,\mathfrak f\,v_0}, is an element
$u$ of \m{\UE} with \m{\sd{u}=r_0\sum_{1\le a'\le\sro}h(t,\fr{a'})\,(v_0)_{a'-1}}. Every $E$-linear form on
\m{E^{\Ix{\fs}}} is \m{h(t,\cdot)} for some~$t$.
\inp{S15} (Lattice discriminants.) Let $E$ be stably unramified. For \m{x\in\Gf{E}{\kappa}} and an $E$-subspace
\m{U\subseteq E^\kappa} let \m{\sld{x^{-1}}{U}} be the discriminant of the lattice \m{(U\otimes\AEf)\cap x\hO^\kappa}. Then for
every \m{m\in\NN}
\[\sum_{U:\ \dim_EU+m=\sk}\sld{x^{-1}}{U}^{-(\sk+1)}\ \le\ \bigl(200\,\fE\,\sk\,\rdE{E}\bigr)^{\fE(\sk+1)m}.\]
Further, let $H$ be a matrix over \m{\AEf} with columns in a finite set $I_1$, having a left inverse,
\m{U\subseteq E^{I_1}} an $E$-subspace of dimension~$d_1$, $x_0$ a matrix over $E$ with $d_1$ columns in~$U$,
\m{w_0\in\GL_{d_1}(\AEf)}, and \m{D_1=\gix{\det\fadm{x_0^*x_0}}\cdot\gix{\det w_0}^{2}\ne0}. If
\m{H\,\fadm{x_0}\,w_0} is integral and all columns of \m{Y_0\,H\,\fadm{x_0}\,w_0} lie in an \m{\hO}-module $K_0$, for a
matrix $Y_0$ over \m{\AEf}, then \m{(\sld{H}{U}\cdot i_0)^{2}\le D_1}, \m{\sld{H}{U}>0} and \m{i_0\ge1}, where
\m{i_0=[\hO\langle Y_0HM_U\rangle:K_0]} and \m{M_U=\slb{H}{U}} is the basis matrix of the lattice of the points of
\m{U\otimes\AEf} that $H$ makes integral.}

\statement{ana}{\ip\emph{From Note~B.} In Setting~\ref{set:data} unless said otherwise.
\inp{B1} \m{\cutoff{S}{c}=\theta_{S,c}\cdot\omega_{S,c}}, where \m{\omega_{S,c}\colon\Ginf\to[0,\infty]} is a
\emph{box partition}, \m{\IsBP{S}{c}{\omega_{S,c}}}: it is measurable; \m{\omega_{S,c}\le1}; if
\m{\omega_{S,c}(y)\ne0} and \m{\InCut{y_\tau}{\bb}} then \m{y_\tau\in\frbox{\bxs{S}}}; \m{\omega_{S,c}(u\,y)=\omega_{S,c}(y)} whenever
every $u_\tau$ lies in the subgroup generated by \m{\Uoone} and \m{\Hone}; and
\m{\sum_{\gamma\in\clU{S}{c}}\alpha(\gamma)\,\omega_{S,c}(\ri{\gamma}y)=1} for every $y$ and every
\m{\alpha\colon\clU{S}{c}\to[0,\infty]} with \m{\sum_{u\in\clUo{S}{c}}\alpha(u\gamma)=1} for all~$\gamma$. Moreover
\m{\cutoff{S}{c}} is measurable and \m{\sum_{u\in\clU{S}{c}}\cutoff{S}{c}(\ri{u}y)=1} for all~$y$.
\inp{B2} If \m{\InCut{y_\tau}{\bb}} for all $\tau$, there is \m{u\in\clU{S}{c}} with \m{\ri{u}\,y\in\bxp{S}}.
\inp{B3} (Frame count.) Let \m{c\in\Cl{S}}, \m{s\in\GRes{S}}, \m{y,y'\in\Ginf}, \m{C_1\in[0,\infty]}, and let
$\mathcal P$ be a set with, for each \m{p\in\mathcal P}, a class $c_p$, elements \m{\gamma_p,\delta_p\in\GS{S}{E}}
and a weight \m{\varphi_p\in[0,\infty]}, such that \m{p\mapsto(c_p,\delta_p)} is injective. Suppose that for every $p$
with \m{\varphi_p\ne0}: \m{\gamma_p\in\terms{S}{c}{s}{c_p}}, \m{\delta_p^{-1}\gamma_p\in\clU{S}{c_p}},
\m{\ri{\gamma_p}^{-1}y\in\bxp{S}} and \m{\varphi_p\le\cutoff{S}{c_p}(\ri{\delta_p}^{-1}y')\cdot C_1}. Let $\mathcal V$ be a
real subspace of \m{\prod_\tau\prod_{a\le\sro}\CC^{\Ix{\fs}}} that contains the differences of the families
\m{\bigl((\ri{\gamma_p}_\tau^{-1}y_\tau)^{-1}\fr{a}\bigr)_{\tau,a}} for any two $p$ with \m{\varphi_p\ne0}. Then
\m{\sum_{p}\varphi_p\le\bb^{\,6\sn^2\dim_{\RR}\mathcal V}\cdot C_1}.
\inp{B4} For a side $\sigma$ and a set \m{\Xi\subseteq\GS{S}{E}}:
\m{\sum_{\gamma\in\Xi\cap\terms{S}{c}{s}{c'}}\ampabs{\sigma}{S}{c}{c'}{\gamma}<\infty}, so the amplitudes are absolutely
summable; and for fixed $c$ and~$s$,
\begin{align*}
\sum_{c'}\ \sum_{\gamma\in\Xi\cap\terms{S}{c}{s}{c'}}\int\cutoff{S}{c}(y)\,\cutoff{S}{c'}(\ri{\gamma}^{-1}y)\,
&\bigl\|\besker{\sigma}{S}{\gamma}{y}-\besker{\cut}{S}{\gamma}{y}\bigr\|\,dy\\
&\le\ \tfrac12\,\nuS\,\bb^{-\sk^3\fE}.
\end{align*}
\inp{B5} (Change of weight.) Let $\sigma$ be a side, $\Xi$ stable, and
\m{\varphi\colon\GS{S}{E}\times\Ginf\to[0,\infty)} measurable in the second variable with
\m{\sum_{(u,u')\in\clU{S}{c}\times\clU{S}{c'}}\varphi(u\gamma u'^{-1},\ri{u}y)=1} for all
\m{\gamma\in\Xi\cap\terms{S}{c}{s}{c'}} and all~$y$. Then the functions
\m{y\mapsto\varphi(\gamma,y)\,\besker{\sigma}{S}{\gamma}{y}} are integrable,
\m{\ampsum{\sigma}{S}{\Xi}{c}{s}{c'}=\sum_\gamma\int\varphi(\gamma,y)\,\besker{\sigma}{S}{\gamma}{y}\,dy}, and
\m{\sum_\gamma\int\varphi(\gamma,y)\,\|\besker{\sigma}{S}{\gamma}{y}-\besker{\cut}{S}{\gamma}{y}\|\,dy\le\tfrac12\nuS\bb^{-\sk^3\fE}};
the sums are over \m{\gamma\in\Xi\cap\terms{S}{c}{s}{c'}}.
\inp{B6} If $\Xi$ is stable and $s$ does not agree, then \m{\ampsum{\sigma}{S}{\Xi}{c}{s}{c'}=0}.
\inp{B7} In Setting~\ref{set:S}: for \m{u,u'\in\UE}, \m{\besker{\sigma}{S}{u\gamma u'^{-1}}{\ri{u}y}=
\uchar{S}{\rfi{u}}^{-1}\,\uchar{S}{\rfi{u'}}\cdot\besker{\sigma}{S}{\gamma}{y}}; and for \m{u\in\clU{S}{c}} the product
over $\tau$ of the Bessel functions is invariant under left translation by~\m{\ri{u}}.
\inp{B8} Let $\eta^0$ be the product over $\tau$ of the Haar measures of \m{\Uoone} fixed by the library; they are
invariant under left and right translations. There is a measure $\mu^0$ on \m{\Ginf}, product of measures on the
factors, such that \m{(y,u)\mapsto u\,y} carries \m{\mu^0\times\eta^0} to the Haar measure of~\m{\Ginf} (the
catalogue of the write-up refers for this to Weil's integration formula for a quotient; the second reader of the
reading project could not confirm the cited place and derived the statement from Tonelli's theorem). And if, for
each~$\tau$, \m{\phi_\tau\colon\Gone{\fs}\to[0,\infty]} is measurable and invariant under left translation by \m{\Uone}, then
\m{\int\cutoff{S}{c}(y)\prod_\tau \phi_\tau(y_\tau)\,dy=\nuS\cdot\prod_\tau\int \phi_\tau}, the integrals on the right
over the quotient of \m{\Gone{\fs}} by~\m{\Uone}, as in~\inp{W2}.
\inp{B9} (Colourings.) Let $\mathcal P$ be a set with weights \m{\varphi_p\in[0,\infty]}. \textup{(i)} If a map
from $\mathcal P$ to some set takes at most \m{N_1\in[0,\infty]} values and every non-empty subset on which it is
constant has total weight at most~$C_1$, then \m{\sum_p\varphi_p\le N_1\,C_1}. \textup{(ii)} Let $I_1$ be a finite
set, \m{K_0\subseteq K_1} subgroups of an abelian group with \m{[K_1:K_0]\ne0}, and \m{F_1(p)\in K_1^{I_1}} for
\m{p\in\mathcal P}; if every non-empty \m{\mathcal B\subseteq\mathcal P} with
\m{F_1(p)-F_1(p')\in K_0^{I_1}} for all \m{p,p'\in\mathcal B} has total weight at most~$C_1$, then
\m{\sum_p\varphi_p\le[K_1:K_0]^{\#I_1}\,C_1}. \textup{(iii)} A set of vectors in \m{E^{I_1}} whose entries are
algebraic integers of absolute value at most \m{r_1\ge0} at every infinite place has at most
\m{(1+2r_1)^{2\,\#I_1\fE}} elements. \textup{(iv)} For a unit \m{d_0\in\hO} of \m{\AEf}:
\m{[\hO^{I_1}:d_0\hO^{I_1}]=\gix{d_0}^{\#I_1}\ne0}.
\inp{B10} (Unipotent elements with prescribed first rows.) Let $E$ be stably unramified, $c$ a class, $I_1$ a
finite set, \m{p_1\colon I_1\to\{0,\dots,\sro-1\}}, $x_1$ the matrix over $E$ with the columns \m{e_{p_1(i)}},
$U_1$ a matrix over $E$ with rows in \m{\Ix{\fs}} and columns in~$I_1$ whose $e^*_a$-rows vanish for all
\m{a\le\max p_1}, and $\Psi$ a square matrix over $E$ indexed by $I_1$ with \m{\Psi^*=-\Psi}, such that
\m{\clC{S}{c}\,\fadm{U_1}} and \m{\fadm{\Psi}} are integral. Then there are \m{\gamma_1\in\clU{S}{c}} and
\m{\vartheta_1\in E} with, for \m{J=\JT{\TE}},
\[\gamma_1=1+\gq\cdot\bigl(x_1\,U_1^*J-(U_1+x_1(\vartheta_1\,U_1^*JU_1-\Psi))\,x_1^*J\bigr),\]
and \m{\gamma_1\in\clUo{S}{c}} if the columns of $U_1$ are orthogonal to all frame vectors.
\inp{B11} (Shared vectors.) In Setting~\ref{set:S}: if \m{\HasLev{S}{j}{\gamma}} then \m{\HasLev{S}{j}{\gamma^{-1}}}
and \m{\mnorm{S}{j}{\gamma}>0}; the matrix \m{Y=\padl{j}{\Zf{S}{\gamma}}} is invertible, and for every \m{a\le\sro}
with \m{j\le a} there is \m{v\in E^{\sn}} with \m{\mathfrak f\,v=\gamma\mathfrak f\,(Y^{-1})_{\cdot,a}}, the column
$a$ of~$Y^{-1}$.
\inp{B12} In Setting~\ref{set:S}: if \m{\gamma\mathfrak f} and \m{\gamma'\mathfrak f} have the same middle rows
and the same $e^*$-rows, then for \m{y\in\Gone{\fs}}, every $\tau$ and \m{a'\le\sro},
\begin{align*}
(\ri{\gamma}_\tau^{-1}y)^{-1}\fr{a'}&-(\ri{\gamma'}_\tau^{-1}y)^{-1}\fr{a'}\\
&=\sum_{a<\sro}
\bigl((\gamma-\gamma')_{e_a,\fr{a'}}\bigr)_\tau\,\frac{\bscale{S}{\tau}{\fr{a'}}}{\bscale{S}{\tau}{e_a}}\cdot y^{-1}e_a.
\end{align*}
For matrices $t$, $t'$ over $E$ with rows in \m{\Ix{\fs}},
\[(\clC{S}{c}\fadm{t})^*\,\JT{\Tl}\,(\clC{S}{c}\fadm{t'})=\fadm{t^*\JT{\TE}\,t'},\]
and \m{\clCX{S}{c}^*\,\Tl\,\clCX{S}{c}=1}. The matrix \m{\clC{S}{c}} is the identity on the coordinates $e_a$ and
$e^*_a$, and on the middle coordinates it has the columns of \m{\clCX{S}{c}} on $\kappa$ and the column of~$z$; so it fixes the vectors $e_a$, and it maps \m{\gamma e_a}, read over \m{\AEf}, to \m{\pos{S}{c}{c'}{\gamma}\,e_a},
for all \m{a<\sro}.}

\statement{ana}{\ip\emph{From Note~P} (the note \emph{link-pairing} of this unit). Let $E$ be stably unramified,
\m{j\le\sro+1}, \m{j_1=j-\min(\sro,j)}, and \m{\gamma\in\terms{S}{c}{s}{c'}} with \m{\HasLev{S}{j}{\gamma}}.
\inp{P1} \m{\xri{S}{j}{\pos{S}{c}{c'}{\gamma}}\ge1}, \m{\sati{S}{\pos{S}{c}{c'}{\gamma}}\ge1} and
\m{\cgauge{S}{j}{c}{c'}{\gamma}>0}.
\inp{P2} Let moreover \m{\Agr{S}{s}}, let $X$ be a natural number with \m{\FixN{s}{X}}, and let
\m{\gamma'\in\terms{S}{c}{s}{c'}} with \m{\HasLev{S}{j}{\gamma'}}. Then
\m{\latd{S}{c}(\clink{\gamma}{\gamma'})>0} and
\m{\fvi{S}{c}{\pos{S}{c}{c'}{\gamma}}{1}{\clink{\gamma}{\gamma'}}\ge1}, and there is \m{i\in\NN} such that
\m{\sk\le\dim_E\clink{\gamma}{\gamma'}+i+j_1}, such that \m{\link{\gamma}=\link{\gamma'}} if \m{i=0}, and such that,
with \m{X'=X} if \m{\sro<j} and \m{i=0}, and \m{X'=1} otherwise,
\[\pinv{S}{c}{c'}{\gamma}{\gamma'}\cdot\pinv{S}{c}{c'}{\gamma'}{\gamma}\cdot X'^{\,2}\ \le\
\Bigl(\frac{2}{\gq}\Bigr)^{2i\fE}\cdot\cgauge{S}{j}{c}{c'}{\gamma}\cdot\cgauge{S}{j}{c}{c'}{\gamma'}.\]}

\statement{ana}{\ip\emph{From unit~132 (the real place).}
\inp{R1} The \emph{frame volume} \m{\Delta\colon\Gone{\fs}\to\RR} is continuous and positive, and
\m{\Delta(u\,y)=\Delta(y)} for \m{u\in\Uoone}. (It is the absolute value of the determinant of the Euclidean
products of the vectors \m{y^{-1}\fr{a}} with a basis of a maximal negative subspace.)
\inp{R2} (Volume kernel.) Let \m{y\in\Gone{\fs}} and $\theta_0$ a real number with \m{\bsize{y}\le\theta_0}; let
\m{j\le\sn}, $\Sigma$ a finite set with an injection into \m{\{0,\dots,j-1\}}, $R_0$ an invertible complex matrix
indexed by~$\Sigma$, $Z$ a complex square matrix of size \m{\sn} whose leading block $Z_j$ of size $j$ has
\m{\det Z_j\ne0}, and $V_0$ a complex matrix with rows in $\kappa$ and $j$ columns. Then there is a measurable
\m{\mathcal K\colon\Mat_{\kappa\times\Sigma}(\CC)\to[0,\infty]} with
\[\int\mathcal K\ \le\ \Bigl(\frac{(2\theta_0^2)^{j}}{|\det Z_j|}\Bigr)^{\sk+\sn}\cdot
\Bigl(\frac{\pi^{\#\Sigma}}{(\#\Sigma)!\,(\theta_0^{\#\Sigma}\,|\det R_0|)^{2}}\Bigr)^{\sk}\]
(Lebesgue measure) such that \m{(\Delta(y)/\Delta(\gamma^{-1}y))^{\sk+\sn}\le\mathcal K(A)} for all
\m{\gamma\in\Gone{\fs}} and \m{A\in\Mat_{\kappa\times\Sigma}(\CC)} with: \m{\Zr{\gamma}=Z}; \m{\IsNF{j}{\gamma}}; and
the matrix of the $\kappa$-coordinates of \m{\gamma\fr{0},\dots,\gamma\fr{j-1}} is \m{V_0+A\,R_0\,Z_{\Sigma}}, where
\m{Z_\Sigma} consists of the rows of $Z_j$ with index in the image of~$\Sigma$.
\inp{R3} (Open-cell bound.) Let $\tau$ be an infinite place, \m{\gamma,y\in\Gone{\fs}} and \m{\rho\ge\bb} a real
number with \m{\gdr{\Zr{\gamma}}>0}, \m{y\in\frbox{\rho}} and \m{\gamma^{-1}y\in\frbox{\rho}}. Then there are a compact
set \m{\mathcal C\subseteq\Gone{\fs}}, a finite measure $n_1$ on \m{\Gone{\fs}} with total mass at most
\m{\cellm{\gamma}{\degt{S}{\tau}}{\rho}{\#\It{S}{\tau}}}, a probability space \m{(Q_1,m_1)} and a measurable map
\m{t_1\colon Q_1\to\Gone{\fs}} such that: for $m_1$-almost every $q$, \m{t_1(q)} lies in the subgroup generated by
\m{\Uoone} and \m{\Hone}, and every \m{y'\in\frbox{\rho}} with \m{\bsize{\gamma\,t_1(q)\,y'}\le\rho} lies in
$\mathcal C$; and for every bounded measurable \m{\phi\colon\Gone{\fs}\to\CC},
\[\int_{Q_1}\Bigl\|\int_{\mathcal C}\phi(y')\;\Wt{\full}{S}{\tau}(\gamma\,t_1(q)\,y')\otimes
\cj{\Wt{\full}{S}{\tau}(t_1(q)\,y')}\,dy'\Bigr\|^{2}dm_1(q)\ \le\ \int|\phi|^{2}\,dn_1.\]
\inp{R4} If \m{\bsize{y}\le\theta_0}, there is \m{u\in\Uone} with \m{\|u\,y\|\le\theta_0}.
\inp{R5} (Congruence count.) Let $E$ be stably unramified; let \m{d_0\in\hO} be a unit of \m{\AEf}; $I_1$, $I_2$,
$I_3$ finite sets with \m{\#I_2=\#I_3+m_1}; $\mathcal P$ a set with integral matrices \m{Y_p} (rows in $I_1$,
columns in $I_2$) for \m{p\in\mathcal P}; $M_0$ an integral matrix with rows in $I_2$ and columns in $I_3$ that has an integral
left inverse; \m{K_0\supseteq d_0\hO^{I_1}} an \m{\hO}-module; \m{\eta_p\in\NN} with
\m{[\hO\langle Y_pM_0\rangle:K_0]} dividing \m{\eta_p}; \m{l_1\ge\#I_3}; and \m{C_1\in[0,\infty]} such that every
non-empty \m{\mathcal B\subseteq\mathcal P} with \m{\hO\langle Y_p-Y_{p'}\rangle\subseteq K_0} for all
\m{p,p'\in\mathcal B} has at most $C_1$ elements. Then
\m{\sum_{p\in\mathcal P}\eta_p^{-l_1}\le\gix{d_0}^{\#I_1(\#I_1+m_1)}\cdot C_1}.
\inp{R6} If $M$ is a matrix over $E$ with rows in $\kappa$ such that \m{\clCX{S}{c}\fadm{M}} is integral, and a column of $M$
has \m{\sum_{i\in\kappa}|(M_{i})_\tau|^{2}<1} at every~$\tau$ for its entries $M_i$, then this column is zero.
\inp{R7} If $S$ is admissible, \m{p=(c_1,\gamma_1)} and \m{p'=(c_2,\gamma_2)} satisfy
\m{\IsLab{S}{j}{c}{\mathsf t}{p}} and \m{\IsLab{S}{j}{c}{\mathsf t}{p'}}, and
\m{(u\gamma_1)\mathfrak f=(u'\gamma_2)\mathfrak f} for some \m{u,u'\in\clUo{S}{c}}, then \m{p=p'}.
\inp{R8} Let $E$ be stably unramified, \m{\gamma\in\GS{S}{E}} and \m{\gamma'\in\GS{S}{E}} with
\m{\HasLev{S}{j}{\gamma}}, \m{\HasLev{S}{j}{\gamma'}} and \m{\Zf{S}{\gamma'}=\Zf{S}{\gamma}}; let $\Sigma$ be a set of
indices \m{a<\sro} containing all \m{a<j}, $\xi$ the matrix of the rows of \m{\Zf{S}{\gamma}} with index
in~$\Sigma$, and \m{Y=\padl{j}{\Zf{S}{\gamma}}}. Suppose that, for each middle index, the conjugate of the
corresponding row of
\m{\fmid{\pos{S}{c}{c''}{\gamma'}}-\fmid{\pos{S}{c}{c'}{\gamma}}} lies in
\m{\gq\,\hO\langle\fadm{\xi}^*\rangle+\{v:\gq^{-1}\fadm{Y^{-1}}^*v\ \text{is integral}\}}. Then there is a matrix
$\ell$ over $E$ with rows in $\kappa$ and columns in $\Sigma$ such that \m{\clCX{S}{c}\fadm{\ell}} is integral
and \m{\gamma'\mathfrak f-\gamma\mathfrak f} has the rows \m{\gq\cdot\ell\,\xi} on~$\kappa$ and vanishes on the row
of $z$ and on the rows \m{e^*_a}.}

\statement{ana}{\ip\emph{From the background library.}
\inp{L1} (Cauchy--Schwarz.) For measurable \m{F_1,F_2\ge0} and a function \m{F_0\ge0} on a measure space with
\m{F_0^2\le F_1F_2} pointwise: \m{(\int F_0)^2\le\int F_1\cdot\int F_2}.
\inp{L2} (Class integrals.) Let $\Gamma_0$ be the group of the rational points in \m{\Gf{E}{\kappa}}. There is a
functional \m{F_0\mapsto\int_{\Gamma_0\backslash}F_0} on the real functions on \m{\Gf{E}{\kappa}} that are left
\m{\Gamma_0}-invariant and right invariant under an open subgroup, with these properties: it is additive, and
non-negative on non-negative functions; \m{\int_{\Gamma_0\backslash}F_0(\,\cdot\,x)=\int_{\Gamma_0\backslash}F_0} for
\m{x\in\Gf{E}{\kappa}}; and if $F_0$ is right invariant under \m{\KcolS{S}}, then
\m{\int_{\Gamma_0\backslash}F_0=\sum_{c\in\Cl{S}}F_0(x_c)}, by the triviality of the class automorphism groups.
\inp{L3} Let $\Gamma$ be a discrete subgroup of \m{\Ginf}, and $F_0$ a continuous function on \m{\Ginf} with values
in a finite-dimensional complex vector space, left invariant under~$\Gamma$. If \m{F_0(u\,y)=c_0\,F_0(y)} for all
$y$, for some \m{u\in\Ginf} and a complex number $c_0$ with \m{|c_0|=1}, \m{c_0\ne1}, then the integral of $F_0$
over the quotient by $\Gamma$ is~$0$.
\inp{L4} Let $M$, $M'$ be matrices over a field with the same rows, $J_0$, $J_0'$ sets of column indices of $M$,
$M'$, and $v_0$, $v_0'$ coefficient vectors. If every linear form that vanishes on the columns of $M$ in $J_0$ and
on those of $M'$ in $J_0'$ takes the same value at \m{Mv_0} and at \m{M'v_0'}, then \m{Mv=M'v'} for some $v$, $v'$
that agree with $v_0$, $v_0'$ outside $J_0$, $J_0'$.
\inp{L5} If \m{a_1,a_2\in\AEf} satisfy \m{\psi_{E,f}(t\,a_1)=\psi_{E,f}(t\,a_2)} for all \m{t\in E}, then
\m{a_1=a_2} (the character is non-degenerate and $E$ is dense in \m{\AEf}).
\inp{L6} A non-zero algebraic integer $x$ of $E$ has \m{\prod_\tau|x_\tau|\ge1}; in particular an algebraic integer
with \m{|x_\tau|<1} for all $\tau$ is~$0$.
\inp{L7} A finite commutative ring with $q_1$ elements has at most \m{q_1^{\,n_1^2}} invertible matrices of size~$n_1$. A
complex square matrix of size $n_1$ with entries of absolute value at most $r_1$ has operator norm at most
\m{n_1^2r_1}, and its determinant is at most the $n_1$-th power of its operator norm.
\inp{L8} For \m{x,x'\in\UT{T}}: \m{\ZT{T}{x'^{-1}\gamma x}=M(x')^*\,\ZT{T}{\gamma}\,M(x)}, where \m{M(x)} is the block
of $x$ on the frame, upper triangular with diagonal entries~$1$, and equal to $1$ for \m{x\in\UoT{T}}.
\inp{L9} The map \m{\alpha_1\mapsto u(\alpha_1)} that sends a complex matrix $\alpha_1$ with rows \m{a<\sro} and
columns in~$\kappa$ to the element of \m{\Uoone} with \m{u(\alpha_1)e_a=e_a},
\m{u(\alpha_1)x_i=x_i+\sum_a(\alpha_1)_{a,i}e_a} for the basis vectors $x_i$, \m{i\in\kappa}, and
\m{u(\alpha_1)e^*_a=e^*_a-\sum_i\cj{(\alpha_1)_{a,i}}\,x_i-\tfrac12\sum_{a'}(\alpha_1\alpha_1^*)_{a',a}e_{a'}} is continuous;
every element of \m{\Uoone} is \m{u(\alpha_1)} times an element that moves each $e^*_a$ by a combination of the
$e_{a'}$ and fixes the other basis vectors.
\inp{L10} An element $x$ of \m{\UT{T}} that fixes \m{\sro} vectors whose matrix of $e^*$-coordinates is the
identity is~$1$.
\inp{L11} For a linear map $P_0$ and submodules: if \m{D_2\subseteq P_0^{-1}(K_0)} and \m{K_1\subseteq P_0(M_2)+K_0},
then \m{[K_1:K_0]} divides \m{[M_2:D_2]}.
\inp{L12} (Products of mean-square bounds.) If for each $\tau$ a measurable kernel \m{k_\tau(q,y')}, integrable
in~$y'$, satisfies \m{\int\sum_i|\int\phi(y')k_{\tau,i}(q,y')\,dy'|^2\,dm_\tau(q)\le\int|\phi|^2dn_\tau} for all
bounded measurable~$\phi$, with $\sigma$-finite measures, then the product kernels satisfy the same bound for the
product measures.}

\sect{Energies, and the passage to the cut functions}{s:energy}
\begin{lemma}[splitting]\label{l:split}
In Setting~\ref{set:data}, let $\sigma$ be a side and \m{\Xi,\Xi'\subseteq\GS{S}{E}}. Then, for all
\m{c,c'\in\Cl{S}} and \m{s\in\GRes{S}},
\[\ampsum{\sigma}{S}{\Xi}{c}{s}{c'}=\ampsum{\sigma}{S}{\Xi\setminus\Xi'}{c}{s}{c'}+\ampsum{\sigma}{S}{\Xi\cap\Xi'}{c}{s}{c'}.\]
\end{lemma}
\begin{proof}
By Input~\inp{B4} the amplitudes of the elements of \m{\Xi\cap\terms{S}{c}{s}{c'}} are absolutely summable, and a
sum of an absolutely summable family splits over a partition of its index set.
\end{proof}
\leanline{\lean{amplitudeSumsplit}}

\begin{lemma}\label{l:add}
In Setting~\ref{set:S}, let $V$ be a normed additive group,
\m{\mathcal A,\mathcal A'\colon\Cl{S}\times\GRes{S}\times\Cl{S}\to V} two families, \m{\Omega\subseteq\GRes{S}} and
$\varepsilon$ a real number with \m{\varepsilon>0}. Then
\[\energy{S}{\mathcal A+\mathcal A'}{\Omega}\ \le\ (1+\varepsilon)\,\energy{S}{\mathcal A}{\Omega}+(1+\varepsilon^{-1})\,\energy{S}{\mathcal A'}{\Omega}.\]
\end{lemma}
\begin{proof}
Term by term: for \m{v_1,v_2\in V},
\m{\|v_1+v_2\|^2\le\|v_1\|^2+2\|v_1\|\|v_2\|+\|v_2\|^2} and
\m{2\|v_1\|\|v_2\|\le\varepsilon\|v_1\|^2+\varepsilon^{-1}\|v_2\|^2}; divide by \m{\nuS^2} and sum.
\end{proof}
\leanline{\lean{energyaddle}}

\begin{proof}[Proof of Theorem C]
Fix $c$, $s$, $c'$. By Input~\inp{B4}, for both sides, the amplitudes are absolutely summable and each amplitude is
an absolutely convergent integral, so the block of the difference at \m{(c,s,c')} is
\[\sum_{\gamma\in\Xi\cap\terms{S}{c}{s}{c'}}\int\cutoff{S}{c}(y)\,\cutoff{S}{c'}(\ri{\gamma}^{-1}y)\,
\bigl(\besker{\full}{S}{\gamma}{y}-\besker{\cut}{S}{\gamma}{y}\bigr)\,dy,\]
and its norm is at most the number \m{\epsilon(c,s,c')} obtained by putting the norm inside the integral. By the
second statement of~\inp{B4}, with \m{\sigma=\full}, \m{\sum_{c'}\epsilon(c,s,c')\le\tfrac12\nuS\,\bb^{-\sk^3\fE}}
for all $c$ and~$s$. Since \m{\sum_{c'}\epsilon(c,s,c')^2\le(\sum_{c'}\epsilon(c,s,c'))^2}, summing over the classes $c$ and
over \m{s\in\Omega\subseteq\GRes{S}} gives
\[\energy{S}{\ampsum{\full}{S}{\Xi}{\cdot}{\cdot}{\cdot}-\ampsum{\cut}{S}{\Xi}{\cdot}{\cdot}{\cdot}}{\Omega}\ \le\
\#\Cl{S}\cdot\#\GRes{S}\cdot\tfrac14\,\bb^{-2\sk^3\fE}.\]
The group \m{\GRes{S}} is a product, over the \m{2\fE} chosen places, of general linear groups of size
\m{\rr{\fs}} over fields with \m{\qS{S}{0}} or \m{\qS{S}{1}} elements, \m{\fE} of each kind (Setting~\ref{set:S});
by~\inp{L7}, \m{\#\GRes{S}\le(\qS{S}{0}\qS{S}{1})^{\rr{\fs}^2\fE}=\gq^{\,\rr{\fs}^2\fE}\le\bb^{23\,\rr{\fs}^2\fE}}. Since
\m{2\sro+1\le\sk}, we have \m{\rr{\fs}\le2\sk}, hence \m{23\,\rr{\fs}^2\fE+30\le92\,\sk^2\fE+30\le2\,\sk^3\fE}, the
last because \m{\sk\ge61}. As \m{\bb\ge1},
\m{\#\GRes{S}\cdot\bb^{-2\sk^3\fE}\le\bb^{-30}}. Finally \m{\#\Cl{S}\le\#\Sbar\cdot\#\Cl{S}\cdot\#\Bc{S}=\dd}.
\end{proof}
\leanline{\lean{energysubcutle}}

\sect{Class sizes and Schur weights}{s:schur}
\begin{lemma}[the open-cell mass fits under the class size]\label{l:cell}
In Setting~\ref{set:S}, let $S$ be admissible, and let \m{\gamma\in\GS{S}{E}} be near, \m{\IsNr{S}{\gamma}}, with
\m{\tdnorm{S}{\gamma}\ge1}. Put \m{x=\tdnorm{S}{\gamma}}. Then
\begin{gather*}
\bb^{-\sk^3\fE}\ \le\ \csize{S}{\sro+1}{x}\qquad\text{and}\\
\prod_\tau\cellm{\ri{\gamma}_\tau}{\degt{S}{\tau}}{\bxs{S}}{\#\It{S}{\tau}}\ \le\
\Bigl(\frac{\nuS\cdot\csize{S}{\sro+1}{x}}{2}\Bigr)^{2}.
\end{gather*}
\end{lemma}
\begin{proof}
Write \m{\varkappa=\prod_\tau|\kapS{S}{\tau}|}; by \inp{A2}, \m{\varkappa\ge1}. At the level \m{\sro+1},
Definition~\ref{d:cs} and the formula for \m{\cw{S}{0}} give
\begin{equation}\label{e:cstop}
\csize{S}{\sro+1}{x}^{2}=\frac{\tl^{4}\cdot\dm\cdot\Nn^{\sk\fE}}{\#\Sbar\cdot\#\Cl{S}\cdot\gq^{\,2\sro\sk\fE}\cdot x^{\sk+1}}.
\end{equation}
\emph{First inequality.} By \inp{A3}, \m{\dm\ge\tl^{44}\dd\ge\#\Sbar\cdot\#\Cl{S}}, and \m{\Nn\ge1}, \m{\tl\ge1},
\m{\gq\le\bb^{23}}; so the right side of \eqref{e:cstop} is at least \m{\bb^{-46\sro\sk\fE}x^{-(\sk+1)}}. By
\inp{S4} the number \m{\tdisc{S}{\gamma}_\tau} is real and non-negative; as it equals
\m{|\lmin{\sn}{\Zf{S}{\gamma}}_\tau|^2-\Nn^2|\lmin{\sro}{\Zf{S}{\gamma}}_\tau|^2}, it is at most
\m{|\lmin{\sn}{\Zf{S}{\gamma}}_\tau|^2\le\bb^{11\sn^3}} by~\inp{S5}. Hence \m{x\le\bb^{11\sn^3\fE}} and
\[\csize{S}{\sro+1}{x}^{2}\ \ge\ \bb^{-(46\sro\sk+11\sn^3(\sk+1))\fE}\ \ge\ \bb^{-2\sk^3\fE},\]
because \m{46\sro\le\sk^2} and \m{11\sn^3(\sk+1)\le22\sn^3\sk\le\sk^2}, both by \m{\sk\ge1000(2\sn)^9}.

\emph{Second inequality.} By \inp{S4}, and since \m{\prod_{a\le\sro}\fscale{S}{\tau}{a}=|\kapS{S}{\tau}|^{\sro\sn}\Nn^{-\sn}},
\[\prod_\tau\gdr{\Zr{\ri{\gamma}_\tau}}=\varkappa^{2\sro\sn}\,\Nn^{-2\sn\fE}\,x.\]
With \m{\bxs{S}^{\,9\sn\cdot10^5\sn^4}=\bb^{18\cdot10^5\sn^7}}, \m{\prod_\tau\#\It{S}{\tau}=\#\Bc{S}} and
\m{\prod_\tau\degt{S}{\tau}=\bdeg{S}}, the product of the cell masses is
\[L_0=\frac{\bb^{18\cdot10^5\sn^7\fE}\cdot2^{(\sk-1)\fE}\cdot(\#\Bc{S})^{2}}
{\bigl(\varkappa^{2\sro\sn}\Nn^{-2\sn\fE}x\bigr)^{10^5\sn^4+\sk-1}\cdot\bdeg{S}\cdot\occ{\fs}^{\fE}}.\]
On the other side, \inp{A6} and the definition of \m{\packC{S}} give
\[\frac{\nuS^{2}\,\dm}{\#\Sbar\cdot\#\Cl{S}}\ \ge\ \frac{\gq^{\,\fE\sro(2\sk+2\sro+1)}\cdot\siegC{S}\cdot\volO{E}^{\sn}\cdot(\#\Bc{S})^{2}}
{(1+\bb^{-9})\cdot\bdeg{S}\cdot\occ{\fs}^{\fE}\cdot\fsprod{S}^{2}},\]
so that, by \eqref{e:cstop}, the right side \m{R_0=(\nuS\,\csize{S}{\sro+1}{x}/2)^2} of the claim satisfies
\[R_0\ \ge\ \frac{\tl^{4}\cdot\gq^{\,\fE\sro(2\sro+1)}\cdot\Nn^{\sk\fE}\siegC{S}\cdot\volO{E}^{\sn}\cdot(\#\Bc{S})^{2}}
{4\,(1+\bb^{-9})\cdot x^{\sk+1}\cdot\bdeg{S}\cdot\occ{\fs}^{\fE}\cdot\fsprod{S}^{2}}.\]
Since \m{\sum_{a\le\sro}(\sk+2(\sro-a))=\sn(\sk+\sro)} and
\m{\sum_{a\le\sro}2(\sro-a)(\sk+2(\sro-a))=\sk\sro\sn+\tfrac23\sro\sn(2\sro+1)},
\[\fsprod{S}=\varkappa^{\,\sk\sro\sn+\frac23\sro\sn(2\sro+1)}\cdot\Nn^{-\fE\sn(\sk+\sro)}.\]
In the quotient \m{L_0/R_0} the factors \m{(\#\Bc{S})^2}, \m{\bdeg{S}} and \m{\occ{\fs}^{\fE}} cancel, and the
powers of $\varkappa$, $x$ and \m{\Nn} that remain are \m{\varkappa^{-\epsilon_1}}, \m{x^{-\epsilon_2}} and
\m{\Nn^{\epsilon_3\fE}} with
\begin{gather*}
\epsilon_1=2\sro\sn(10^5\sn^4-1)-\tfrac43\sro\sn(2\sro+1)\ge0,\\
\epsilon_2=10^5\sn^4-2\ge0,\qquad
\epsilon_3=2\sn(10^5\sn^4-1-\sro)\le2\cdot10^5\sn^5.
\end{gather*}
Using \m{\varkappa\ge1}, \m{x\ge1}, \m{\gq\ge1}, \m{\Nn\le\bb^{60\sn}} \inp{A2}, \m{\tl^4=\bb^{\sk\fE}} \inp{S1},
\m{4(1+\bb^{-9})\le8} and \m{\Nn^{\sk\fE}\siegC{S}=(\pi^{\sk}/(\sk-1)!)^{\fE}|\dE{E}|^{1/4}\volO{E}^{\sk}}, we get
\[\frac{L_0}{R_0}\ \le\ \frac{\bb^{(18\cdot10^5\sn^7+12\cdot10^6\sn^6)\fE}\cdot2^{(\sk-1)\fE+3}}
{\bb^{\sk\fE}\cdot(\pi^{\sk}/(\sk-1)!)^{\fE}\cdot|\dE{E}|^{1/4}\cdot\volO{E}^{\sk+\sn}}.\]
Now \m{|\dE{E}|^{1/4}\volO{E}^{\sk+\sn}=2^{\fE(\sk+\sn)}\,\rdE{E}^{-\fE(\sk+\sn-1/2)}} and
\m{\pi^{\sk}/(\sk-1)!\ge2^{\sk}\sk^{-\sk}}; the power of $2$ that results is \m{2^{3-\fE(\sk+\sn+1)}\le1}. Further
\m{\sk^{\sk\fE}\,\rdE{E}^{\fE(\sk+\sn-1/2)}\le(\sk\,\rdE{E})^{\fE(\sk+\sn)}\le\bb^{\fE(\sk+\sn)/1000}}, by
\m{\rdE{E}\ge1} and the last inequality of Setting~\ref{set:S}. Altogether
\[\frac{L_0}{R_0}\ \le\ \bb^{\,\fE\,(18\cdot10^5\sn^7+12\cdot10^6\sn^6+(\sk+\sn)/1000-\sk)}\ \le\ 1,\]
since \m{18\cdot10^5\sn^7+12\cdot10^6\sn^6\le3\cdot10^6\sn^7\le\sk/2} and \m{(\sk+\sn)/1000\le\sk/2}, again by
\m{\sk\ge1000(2\sn)^9}.
\end{proof}
\leanline{\lean{leclassSizeandcellMassle}}

\begin{lemma}\label{l:incut}
In Setting~\ref{set:S}, let \m{\gamma\in\GS{S}{E}} and \m{y\in\Ginf} with \m{\besker{\cut}{S}{\gamma}{y}\ne0}. Then,
for every infinite place~$\tau$, \m{\InCut{y_\tau}{\bb}} and \m{\InCut{(\ri{\gamma}^{-1}y)_\tau}{\bb}}.
\end{lemma}
\begin{proof}
The norm of the kernel is the product of the norms of the vectors \m{\Wt{\cut}{S}{\tau}(y_\tau)} and
\m{\Wt{\cut}{S}{\tau}((\ri{\gamma}^{-1}y)_\tau)}; so all of them are non-zero, and \inp{W3} applies.
\end{proof}
\leanline{\lean{inCutofbesselKernelnezero}}

\begin{lemma}[Cauchy--Schwarz with Schur weights]\label{l:cs}
In Setting~\ref{set:data}, for \m{c,c'\in\Cl{S}} and \m{\gamma\in\GS{S}{E}},
\[\ampabs{\cut}{S}{c}{c'}{\gamma}^{2}\ \le\ \schi{S}{c}{c'}{\gamma}\cdot\schi{S}{c'}{c}{\gamma^{-1}}.\]
\end{lemma}
\begin{proof}
The boxed cutoffs \m{\bcut{S}{c}} and the Schur weights \m{y\mapsto\schw{S}{c}{\gamma}{y}} are measurable: the cutoff is
measurable \inp{B1}, the set \m{\bxp{S}} is measurable, and $\Delta$ is continuous and positive \inp{R1}. In
the second Schur integral substitute \m{y\mapsto\ri{\gamma}^{-1}y}, which preserves the Haar measure; as
\m{\ri{\gamma^{-1}}^{-1}=\ri{\gamma}}, it becomes the integral of
\[F_2(y)=\bcut{S}{c'}(\ri{\gamma}^{-1}y)\prod_\tau\|\Wt{\cut}{S}{\tau}((\ri{\gamma}^{-1}y)_\tau)\|^{2}\cdot
\bcut{S}{c}(y)\prod_\tau\Bigl(\frac{\Delta((\ri{\gamma}^{-1}y)_\tau)}{\Delta(y_\tau)}\Bigr)^{\sk+\sn},\]
while the first is the integral of
\[F_1(y)=\bcut{S}{c}(y)\prod_\tau\|\Wt{\cut}{S}{\tau}(y_\tau)\|^{2}\cdot
\bcut{S}{c'}(\ri{\gamma}^{-1}y)\prod_\tau\Bigl(\frac{\Delta(y_\tau)}{\Delta((\ri{\gamma}^{-1}y)_\tau)}\Bigr)^{\sk+\sn}.\]
Let \m{F_0(y)=\cutoff{S}{c}(y)\,\cutoff{S}{c'}(\ri{\gamma}^{-1}y)\,\|\besker{\cut}{S}{\gamma}{y}\|} be the integrand of the
absolute amplitude. We claim \m{F_0^2\le F_1F_2} pointwise; then \inp{L1} gives the lemma. Let \m{F_0(y)\ne0}. Then
the kernel is non-zero, so by Lemma~\ref{l:incut} all \m{y_\tau} and all \m{(\ri{\gamma}^{-1}y)_\tau} are in the cut;
and \m{\omega_{S,c}(y)\ne0}, \m{\omega_{S,c'}(\ri{\gamma}^{-1}y)\ne0}, so by \inp{B1} both points have all their
components in \m{\frbox{\bxs{S}}}. Thus $y$ and \m{\ri{\gamma}^{-1}y} are boxed points, the boxed cutoffs coincide there
with the cutoffs, the frame volumes cancel, and \m{F_1(y)F_2(y)=F_0(y)^2}.
\end{proof}
\leanline{\lean{amplitudeAbssqle}\sep\lean{measurableboxCutoff}\sep\lean{measurableschurWeight}}

\begin{proposition}[averaging over unipotent translates]\label{p:avg}
In Setting~\ref{set:data}, let \m{c\in\Cl{S}}; let $\mu$ be an $s$-finite measure on a measurable space~$Q$, and
\m{\phi\colon Q\to\Ginf} a measurable map with \m{\phi(q)_\tau\in\Uoone} for all \m{q\in Q} and all~$\tau$; let
\m{F\colon\Ginf\to[0,\infty]} be measurable with \m{F(\ri{u}\,y)=F(y)} for all \m{u\in\clUo{S}{c}} and
\m{y\in\Ginf}; let \m{Q_0\subseteq Q} with \m{\mu(Q_0)\ge1}; and let \m{C_0\in[0,\infty]} be such that
\[\int_{Q_0}F\bigl(u\,\phi(q)\,y\bigr)\,d\mu(q)\ \le\ C_0\]
for every \m{u\in\Ginf} with all \m{u_\tau\in\Uoone} and every \m{y\in\bxp{S}}. Then
\[\int_{\Ginf}\bcut{S}{c}(y)\cdot\prod_\tau\bigl\|\Wt{\cut}{S}{\tau}(y_\tau)\bigr\|^{2}\cdot F(y)\,dy\ \le\
C_0\cdot\nuS\cdot\#\Bc{S}.\]
\end{proposition}
\begin{proof}
Write \m{W(y)=\prod_\tau\|\Wt{\cut}{S}{\tau}(y_\tau)\|^2} and \m{\Uoone^{\infty}} for the group of the \m{u\in\Ginf} with all
\m{u_\tau\in\Uoone}, with its Haar measure $\eta^0$ of~\inp{B8}. By \inp{W1} the function $W$ is invariant under
\m{\Uoone^\infty}, even under all $u$ with \m{u_\tau\in\Uone}; so by \inp{B8} and \inp{W2}
\begin{equation}\label{e:avg1}
\nuS\cdot\#\Bc{S}\ \ge\ \int\cutoff{S}{c}(y)\,W(y)\,dy=\int\!\!\int\omega_{S,c}(y)\,W(y)\;\theta_{S,c}(u\,y)\,d\eta^0(u)\,d\mu^0(y),
\end{equation}
where we used the measure $\mu^0$ of~\inp{B8}, the factorization \m{\cutoff{S}{c}=\theta_{S,c}\omega_{S,c}} and the invariance of
\m{\omega_{S,c}} under \m{\Uoone^\infty} \inp{B1}. In the same way, because \m{u\,y} is a boxed point exactly if $y$ is
\inp{S3}, the left side of the claim is
\begin{equation}\label{e:avg2}
\int\mathbf 1_{\bxp{S}}(y)\;\omega_{S,c}(y)\,W(y)\Bigl(\int\theta_{S,c}(u\,y)\,F(u\,y)\,d\eta^0(u)\Bigr)d\mu^0(y).
\end{equation}
So it suffices to show, for every boxed point~$y$,
\begin{equation}\label{e:avg3}
\int\theta_{S,c}(u\,y)\,F(u\,y)\,d\eta^0(u)\ \le\ C_0\int\theta_{S,c}(u\,y)\,d\eta^0(u).
\end{equation}
Fix $y$ and \m{u_1\in\Uoone^\infty}. Since \m{\sum_{u'\in\clUo{S}{c}}\theta_{S,c}(\ri{u'}\,y')=1} for all~$y'$
\inp{S9}, and \m{\ri{u'}\in\Uoone^\infty} \inp{S3},
\begin{align*}
\int\theta_{S,c}(uy)\,F(u\,u_1y)\,d\eta^0(u)
&=\sum_{u'}\int\theta_{S,c}(uy)\;\theta_{S,c}(\ri{u'}^{-1}u\,u_1y)\,F(\ri{u'}^{-1}u\,u_1y)\,d\eta^0(u)\\
&=\sum_{u'}\int\theta_{S,c}(\ri{u'}\,uy)\;\theta_{S,c}(u\,u_1y)\,F(u\,u_1y)\,d\eta^0(u)\\
&=\int\theta_{S,c}(u\,u_1y)\,F(u\,u_1y)\,d\eta^0(u)=\int\theta_{S,c}(uy)\,F(uy)\,d\eta^0(u).
\end{align*}
In the first line we inserted the sum over \m{u'^{-1}} at the point \m{u\,u_1y} and used the invariance of $F$
under \m{\clUo{S}{c}}; in the second we substituted \m{u\mapsto\ri{u'}u}; in the third we summed over $u'$ at
the point \m{uy} and used the right invariance of~$\eta^0$. Apply this with \m{u_1=\phi(q)} and integrate over
\m{q\in Q_0}: as \m{\mu(Q_0)\ge1}, Tonelli's theorem and the hypothesis give
\[\int\theta_{S,c}(uy)\,F(uy)\,d\eta^0(u)\ \le\ \int\theta_{S,c}(uy)\Bigl(\int_{Q_0}F\bigl(u\,\phi(q)\,y\bigr)d\mu(q)\Bigr)d\eta^0(u)
\ \le\ C_0\int\theta_{S,c}(uy)\,d\eta^0(u),\]
which is \eqref{e:avg3}. Multiplying \eqref{e:avg3} by \m{\omega_{S,c}(y)W(y)}, integrating against $\mu^0$ over
the boxed points and comparing \eqref{e:avg2} with \eqref{e:avg1} proves the proposition.
\end{proof}
\leanline{\lean{lintegralboxCutoffleofaverage}}

\sect{Thin classes: the proof of Theorem A}{s:thin}
\begin{lemma}[the terms of one class]\label{l:class}
In Setting~\ref{set:S}, let \m{j\in\{0,\dots,\sro+1\}}, \m{c,c'\in\Cl{S}}, \m{s\in\GRes{S}}, and
\m{\gamma\in\GS{S}{E}} with \m{\gamma\in\terms{S}{c}{s}{c'}} and \m{\HasLev{S}{j}{\gamma}}. Let $\delta$ be an element of
the double coset \m{\clUo{S}{c}\,\gamma\,\clUo{S}{c'}}. Then \m{\HasLev{S}{j}{\delta}},
\m{\Zf{S}{\delta}=\Zf{S}{\gamma}}, \m{\delta\in\terms{S}{c}{s}{c'}}, and for every middle index the conjugate of the
corresponding row of \m{\fmid{\pos{S}{c}{c'}{\delta}}-\fmid{\pos{S}{c}{c'}{\gamma}}} lies in
\m{\gq\cdot\mathcal R(\pos{S}{c}{c'}{\gamma})}.
\end{lemma}
\begin{proof}
Write \m{\delta=u\gamma u'} with \m{u\in\clUo{S}{c}}, \m{u'\in\clUo{S}{c'}}. The level is the same by~\inp{S10}; the
frame Gram matrix is the same by \inp{L8}, since $u^{-1}$ and $u'$ lie in \m{\UoE}; and
\m{\delta\in\terms{S}{c}{1\cdot s\cdot1}{c'}} by \inp{S6}, since \m{u\in\terms{S}{c}{1}{c}} and
\m{u'\in\terms{S}{c'}{1}{c'}}. For the last statement put \m{\mathsf k=\pos{S}{c}{c'}{\gamma}} and
\m{\mathsf u=\pos{S}{c}{c}{u}\in\KgS{S}}. Since $u'$ fixes the frame and \m{\hc{S}{c'}} lies in \m{\HT{\Tl}}, the
matrix \m{\pos{S}{c}{c'}{\delta}\,\mathfrak f} is \m{\mathsf u\,\mathsf k\,\mathfrak f}. The element \m{\mathsf u} lies in
\m{\UT{\Tl}}, because \m{\hc{S}{c}\in\HT{\Tl}} normalizes this group; so for every vector $t$ the middle
coordinates of \m{\mathsf u\,t} are those of $t$ plus \m{\sum_{a<\sro}\mathsf u_{i,e^*_a}\,t_{e^*_a}} at the middle index~$i$.
For \m{t=\mathsf k\fr{a'}} the coordinate \m{t_{e^*_a}} is \m{h(e_a,\mathsf k\fr{a'})=\Zl{\mathsf k}_{a,a'}}. Hence the row $i$ of
\m{\fmid{\mathsf u\mathsf k}-\fmid{\mathsf k}} is \m{\sum_{a<\sro}\mathsf u_{i,e^*_a}} times the row $a$ of \m{\Zl{\mathsf k}},
and \m{\mathsf u_{i,e^*_a}\in\gq\,\hO} by~\inp{S7}. As \m{\hO} and $\gq$ are stable under conjugation, the claim follows.
\end{proof}
\leanline{\lean{termsofclassspec}}

\begin{proof}[Proof of Theorem A\,(a)]
Put \m{j_0=\min(\sro,j)=\#\IdxB{j}}, \m{j_1=j-j_0\in\{0,1\}}, \m{Z=\Zf{S}{\gamma}},
\m{\mathsf k=\pos{S}{c}{c'}{\gamma}}, and let $\xi$ be the matrix of the rows of $Z$ with index in~\m{\IdxB{j}}.

\emph{Step 1: averaging.} Let \m{Q=\prod_\tau\Mat_{\kappa\times\IdxB{j}}(\CC)} with Lebesgue measure~$\mu$, and let
$Q_0$ be the set of the \m{q\in Q} all of whose entries have absolute value less than~$1$; then
\m{\mu(Q_0)=\pi^{\fE\sk j_0}\ge1}. For each $\tau$ let $R_\tau$ be the diagonal matrix indexed by \m{\IdxB{j}} with
the entries \m{\pi\,\gq/(\bb\cdot\bscale{S}{\tau}{e_a})}, and for \m{q\in Q} let
\m{\phi(q)_\tau=u((q_\tau R_\tau)^*)\in\Uoone} in the notation of~\inp{L9}, the matrix \m{(q_\tau R_\tau)^*} being
completed by zero rows. The map $\phi$ is continuous. Let
\[F(y)=\sum_{\delta\in\mathcal D}\schw{S}{c'}{\delta}{y}.\]
By monotone convergence the left side of \textup{(a)} is the integral of
\m{\bcut{S}{c}(y)\prod_\tau\|\Wt{\cut}{S}{\tau}(y_\tau)\|^2F(y)}. The function $F$ is measurable, and it is invariant
under \m{\clUo{S}{c}}: for \m{u_0\in\clUo{S}{c}}, \m{\Delta(\ri{u_0}_\tau y_\tau)=\Delta(y_\tau)} by \inp{S3}
and~\inp{R1}, so \m{\schw{S}{c'}{\delta}{\ri{u_0}y}=\schw{S}{c'}{u_0^{-1}\delta}{y}}, and \m{\delta\mapsto u_0^{-1}\delta}
permutes~$\mathcal D$. By Proposition~\ref{p:avg} it suffices to prove, for every \m{u\in\Ginf} with all
\m{u_\tau\in\Uoone} and every \m{y\in\bxp{S}},
\begin{equation}\label{e:A1}
\int_{Q_0}F\bigl(u\,\phi(q)\,y\bigr)\,d\mu(q)\ \le\ \tl^{2}\cdot(\bb^{4}\gq)^{-j_0\sk\fE}\cdot
\Bigl(\frac{\xri{S}{j}{\mathsf k}}{\sqrt{\mnorm{S}{j}{\gamma}}}\Bigr)^{\sk+1}.
\end{equation}
Fix such $u$ and~$y$.

\emph{Step 2: colouring.} By Lemma~\ref{l:class}, for \m{\delta\in\mathcal D} the conjugates of the middle rows of
\m{\fmid{\pos{S}{c}{c'}{\delta}}-\fmid{\mathsf k}} lie in \m{\gq\,\mathcal R(\mathsf k)}. We apply \inp{B9}\,(ii)
to the set $\mathcal D$, with $I_1=\Mx{\fs}$, which has \m{\sk+1} elements, \m{K_1=\gq\,\mathcal R(\mathsf k)},
\m{K_0=\gq\,\mathcal R_j(\mathsf k)} and the weights \m{\int_{Q_0}\schw{S}{c'}{\delta}{u\phi(q)y}\,d\mu(q)}. Here
\m{[K_1:K_0]=\xri{S}{j}{\mathsf k}}, which is not zero by~\inp{P1}. So \eqref{e:A1} follows if, for every
non-empty \m{\mathcal B\subseteq\mathcal D} such that for all \m{\delta,\delta'\in\mathcal B} the conjugates of the
middle rows of \m{\fmid{\pos{S}{c}{c'}{\delta}}-\fmid{\pos{S}{c}{c'}{\delta'}}} lie in \m{\gq\,\mathcal R_j(\mathsf k)},
\begin{equation}\label{e:A2}
\sum_{\delta\in\mathcal B}\int_{Q_0}\schw{S}{c'}{\delta}{u\,\phi(q)\,y}\,d\mu(q)\ \le\
\tl^{2}\cdot(\bb^{4}\gq)^{-j_0\sk\fE}\cdot\mnorm{S}{j}{\gamma}^{-(\sk+1)/2}.
\end{equation}
Fix such a set $\mathcal B$ and an element \m{\delta_0\in\mathcal B}.

\emph{Step 3: the terms of $\mathcal B$ are lattice translates.} Let \m{\delta\in\mathcal B}. By
Lemma~\ref{l:class}, $\delta$ and $\delta_0$ have level~$j$ and the frame Gram matrix~$Z$, and
\m{\mathcal R_j(\mathsf k)=\hO\langle\fadm{\xi}^*\rangle} by~\inp{S6}. So \inp{R8}, with \m{\Sigma=\IdxB{j}}, gives a
matrix \m{\ell_\delta} over $E$ with rows in $\kappa$ and columns in \m{\IdxB{j}} such that
\m{\clCX{S}{c}\fadm{\ell_\delta}} is integral and
\begin{equation}\label{e:A3}
\delta\mathfrak f-\delta_0\mathfrak f\ \text{ has the rows }\ \gq\cdot\ell_\delta\,\xi\ \text{ on }\kappa
\text{ and vanishes on the row of }z\text{ and on the rows }e^*_a.
\end{equation}
Put \m{L_\delta=((\bb/\pi)\,(\ell_\delta)_\tau)_\tau\in Q}, the embeddings being applied to the entries.

\emph{Step 4: one kernel for all terms.} For each $\tau$ let \m{Z^\tau=\Zr{\ri{\gamma}_\tau}}. By \inp{S4} its leading
block \m{Z^\tau_j} of size $j$ has the determinant
\m{\prod_{a<j}\fscale{S}{\tau}{a}\cdot\lmin{j}{Z}_\tau}, which is not zero because $\gamma$ has level~$j$. Since $y$
is a boxed point, \m{\bsize{y_\tau}\le\bb^{6\sn+2}}. Let \m{\mathcal K_\tau} be the kernel of \inp{R2} for the point
\m{y_\tau}, the bound \m{\theta_0=\bb^{6\sn+2}}, the set \m{\Sigma=\IdxB{j}\subseteq\{0,\dots,j-1\}}, the matrix
\m{R_0=R_\tau}, the matrix \m{Z^\tau}, and $V_0$ the matrix of the $\kappa$-coordinates of the vectors
\m{u_\tau^{-1}\ri{\delta_0}_\tau\fr{a}}, \m{a<j}. We claim that for all \m{\delta\in\mathcal B}, \m{q\in Q} and~$\tau$
\begin{equation}\label{e:A4}
\Bigl(\frac{\Delta((u\phi(q)y)_\tau)}{\Delta((\ri{\delta}^{-1}u\phi(q)y)_\tau)}\Bigr)^{\sk+\sn}\ \le\
\mathcal K_\tau\bigl(q_\tau+(L_\delta)_\tau\bigr).
\end{equation}
Indeed, let \m{u_1=u_\tau\phi(q)_\tau\in\Uoone} and \m{\gamma_1=u_1^{-1}\ri{\delta}_\tau}. By \inp{R1} the left side
of \eqref{e:A4} is \m{(\Delta(y_\tau)/\Delta(\gamma_1^{-1}y_\tau))^{\sk+\sn}}. We check the three conditions
of~\inp{R2} for $\gamma_1$ and \m{A=q_\tau+(L_\delta)_\tau}. First, \m{\Zr{\gamma_1}=\Zr{\ri{\delta}_\tau}=Z^\tau}
by \inp{L8} and~\inp{S4}, since \m{\Zf{S}{\delta}=Z}. Second, \m{\IsNF{j}{\gamma_1}} by~\inp{S12}. Third, an element
of \m{\Uoone} leaves the $e^*$-coordinates of a vector unchanged, because it is an isometry that fixes the
vectors~$e_a$, and it changes the $\kappa$-coordinates only by a combination of the $e^*$-coordinates; for
\m{u(\alpha_1)^{-1}} the change is \m{\alpha_1^*} applied to the $e^*$-coordinates \inp{L9}. By \eqref{e:A3} and the definition of \m{\rinf{S}}, the vector
\m{\ri{\delta}_\tau\fr{a}-\ri{\delta_0}_\tau\fr{a}} has no $e^*$-coordinates and the $\kappa$-coordinates
\m{\gq\,((\ell_\delta\,\xi)_{\cdot,a})_\tau\,\bscale{S}{\tau}{\fr{a}}}; so the same holds after applying
\m{u_\tau^{-1}}. The $e^*_{a'}$-coordinate of \m{u_\tau^{-1}\ri{\delta}_\tau\fr{a}} is
\m{h(e_{a'},u_\tau^{-1}\ri{\delta}_\tau\fr{a})=Z^\tau_{a',a}}, because \m{u_\tau} fixes~$e_{a'}$. Applying
\m{\phi(q)_\tau^{-1}=u((q_\tau R_\tau)^*)^{-1}} we find that the $\kappa$-coordinates of \m{\gamma_1\fr{a}}, for
\m{a<j}, are the column $a$ of
\[V_0+\gq\,\bigl((\ell_\delta\,\xi)_{\cdot,a}\bigr)_\tau\,\bscale{S}{\tau}{\fr{a}}+q_\tau R_\tau Z^\tau_\Sigma.\]
Finally \m{Z^\tau_{a',a}=\bscale{S}{\tau}{e_{a'}}\,\bscale{S}{\tau}{\fr{a}}\,(Z_{a',a})_\tau} by \inp{S4}, so by the
choice of $R_\tau$ the middle term is the column $a$ of \m{(L_\delta)_\tau R_\tau Z^\tau_\Sigma}. This is the third
condition, and \eqref{e:A4} follows.

\emph{Step 5: translation.} For \m{\delta\in\mathcal B} and \m{q\in Q} put
\[\Phi_\delta(q)=\mathbf 1_{Q_0}(q-L_\delta)\cdot\schw{S}{c'}{\delta}{u\,\phi(q-L_\delta)\,y}.\]
By the invariance of Lebesgue measure, \m{\int_{Q_0}\schw{S}{c'}{\delta}{u\phi(q)y}\,d\mu(q)=\int_Q\Phi_\delta\,d\mu}, and
by \eqref{e:A4} at the point \m{q-L_\delta}
\begin{equation}\label{e:A5}
\Phi_\delta(q)\ \le\ \mathbf 1_{Q_0}(q-L_\delta)\cdot\bcut{S}{c'}\bigl(\ri{\delta}^{-1}u\,\phi(q-L_\delta)\,y\bigr)\cdot
\prod_\tau\mathcal K_\tau(q_\tau).
\end{equation}
Fix \m{q\in Q}. If \m{\Phi_\delta(q)\ne0} and \m{\Phi_{\delta'}(q)\ne0}, then \m{q-L_\delta} and \m{q-L_{\delta'}} lie
in~$Q_0$, so every entry of \m{(\ell_\delta-\ell_{\delta'})_\tau} has absolute value less than
\m{2\pi/\bb\le8/\bb}, and for every column the sum over $\kappa$ of the squares of these absolute values is less
than \m{64\,\sk/\bb^2\le64/\bb<1}, by \m{\sk\le\bb}. As \m{\clCX{S}{c}\fadm{\ell_\delta-\ell_{\delta'}}} is
integral, \inp{R6} gives \m{\ell_\delta=\ell_{\delta'}}. So there is one matrix \m{\ell} with
\m{\ell_\delta=\ell} for all $\delta$ with \m{\Phi_\delta(q)\ne0}; let $L$ be the corresponding point of $Q$ and
\m{y''=u\,\phi(q-L)\,y}.

\emph{Step 6: the frame count.} We apply \inp{B3} to the index set~$\mathcal B$, with the classes
\m{c_\delta=c'}, the elements \m{\gamma_\delta=\delta_\delta=\delta}, the weights \m{\Phi_\delta(q)}, the points
\m{y''} and \m{y''}, and \m{C_1=\prod_\tau\mathcal K_\tau(q_\tau)}. If \m{\Phi_\delta(q)\ne0}, then
\m{\delta\in\terms{S}{c}{s}{c'}} (Lemma~\ref{l:class}), \m{\ri{\delta}^{-1}y''} is a boxed point because the boxed
cutoff does not vanish there, and \m{\Phi_\delta(q)\le\cutoff{S}{c'}(\ri{\delta}^{-1}y'')\cdot C_1} by~\eqref{e:A5}.
For two such \m{\delta,\delta'}, by \eqref{e:A3} with \m{\ell_\delta=\ell_{\delta'}}, the matrices
\m{\delta\mathfrak f} and \m{\delta'\mathfrak f} have the same middle rows and the same $e^*$-rows; so by
\inp{B12} the differences of the vectors \m{(\ri{\delta}_\tau^{-1}y''_\tau)^{-1}\fr{a}} lie in the complex span of
the \m{\sro} vectors \m{(y''_\tau)^{-1}e_{a'}}. These spans form a real subspace $\mathcal V$ of dimension at most
\m{2\,\sro\,\sn\,\fE}. Hence, with~\inp{S1} and \m{12\,\sro\,\sn^3\le10^5\sn^9},
\[\sum_{\delta\in\mathcal B}\Phi_\delta(q)\ \le\ \bb^{12\sro\sn^3\fE}\cdot\prod_\tau\mathcal K_\tau(q_\tau)\ \le\
\tl\cdot\prod_\tau\mathcal K_\tau(q_\tau).\]
Integrating over $q$, the left side of \eqref{e:A2} is at most \m{\tl\cdot\prod_\tau\int\mathcal K_\tau}.

\emph{Step 7: the integral of the kernels.} Let \m{P_\tau=\prod_{a<j_0}|\kapS{S}{\tau}|^{2(\sro-a)}\ge1}
\inp{A2}. Then \m{|\det R_\tau|^2=(\pi\gq/\bb)^{2j_0}\,\Nn^{j_0}/P_\tau} and
\m{|\det Z^\tau_j|=P_\tau\,\Nn^{-j}\,|\lmin{j}{Z}_\tau|}. With \m{\theta_0=\bb^{6\sn+2}}, the bound of \inp{R2} is
\[\int\mathcal K_\tau\ \le\ \frac{(2\theta_0^2\,\Nn)^{j(\sk+\sn)}}{P_\tau^{\,\sn}\,|\lmin{j}{Z}_\tau|^{\sk+\sn}}\cdot
\frac{\bb^{2j_0\sk}}{(j_0!)^{\sk}\,\theta_0^{2j_0\sk}\,\pi^{j_0\sk}\,\gq^{2j_0\sk}\,\Nn^{j_0\sk}}.\]
The number \m{\lmin{j}{Z}} is a non-zero algebraic integer \inp{S6}; so
\m{\Pi_0=\prod_\tau|\lmin{j}{Z}_\tau|\ge1} by~\inp{L6}, and \m{\Pi_0^2=|\Nm(\lmin{j}{Z})|=\mnorm{S}{j}{\gamma}\,\Nn^{2j_1\fE}}.
Hence \m{\Pi_0^{-(\sk+\sn)}\le\Pi_0^{-(\sk+1)}=\mnorm{S}{j}{\gamma}^{-(\sk+1)/2}\,\Nn^{-j_1\fE(\sk+1)}}. Using
\m{P_\tau\ge1}, \m{j_0!\ge1}, \m{\pi\ge2} and \m{j=j_0+j_1}, the product over $\tau$ is at most
\begin{gather*}
\mnorm{S}{j}{\gamma}^{-(\sk+1)/2}\cdot(\bb^4\gq)^{-j_0\sk\fE}\cdot\beta_1^{\,\fE}\cdot\beta_2^{\,\fE},\\
\beta_1=(2\theta_0^2)^{j_1\sk}\,\bigl(\bb^{6}/\gq\bigr)^{j_0\sk},\qquad\beta_2=(2\theta_0^2\,\Nn)^{j\sn}.
\end{gather*}
The number $\beta_1$ is at most~$1$: for \m{j_1=0} because \m{\gq\ge\bb^{19}}; for \m{j_1=1}, where \m{j_0=\sro}, because
\m{2\theta_0^2\,(\bb^6/\gq)^{\sro}\le2\,\bb^{12\sn+4-13\sro}=2\,\bb^{16-\sro}\le1}. The number $\beta_2$ is at most
\m{(2\,\bb^{12\sn+4}\,\bb^{60\sn})^{\sn^2}\le\bb^{(72\sn+5)\sn^2}}, by \inp{A2}, and
\m{\bb^{(72\sn+5)\sn^2\fE}\le\tl} by~\inp{S1}. So \m{\prod_\tau\int\mathcal K_\tau\le\tl\cdot(\bb^4\gq)^{-j_0\sk\fE}\cdot\mnorm{S}{j}{\gamma}^{-(\sk+1)/2}},
and with Step~6 this is \eqref{e:A2}.
\end{proof}
\leanline{\lean{tsumschurIntegralle}}

\begin{proof}[Proof of Theorem A\,(b)]
Put \m{\mathsf k=\pos{S}{c}{c'}{\gamma}}, \m{j_0=\min(\sro,j)} and \m{D_2=\mnorm{S}{j}{\gamma}}, which is positive
by~\inp{B11}. By Lemma~\ref{l:cs} and \inp{L1} for the counting measure on~$\mathcal D$,
\[\Bigl(\sum_{\delta\in\mathcal D}\ampabs{\cut}{S}{c}{c'}{\delta}\Bigr)^{2}\ \le\
\sum_{\delta\in\mathcal D}\schi{S}{c}{c'}{\delta}\cdot\sum_{\delta\in\mathcal D}\schi{S}{c'}{c}{\delta^{-1}}.\]
The first factor is bounded by part~(a). In the second, \m{\delta^{-1}} runs through the double coset
\m{\clUo{S}{c'}\,\gamma^{-1}\,\clUo{S}{c}}; here \m{\gamma^{-1}\in\terms{S}{c'}{s^{-1}}{c}} \inp{S6}, \m{\gamma^{-1}} has
level $j$ \inp{B11}, \m{\pos{S}{c'}{c}{\gamma^{-1}}=\mathsf k^{-1}}, and \m{\mnorm{S}{j}{\gamma^{-1}}=D_2} because
\m{\Zf{S}{\gamma^{-1}}=\Zf{S}{\gamma}^*}. So part~(a) applies to it as well, and
\[\Bigl(\sum_{\delta\in\mathcal D}\ampabs{\cut}{S}{c}{c'}{\delta}\Bigr)^{2}\ \le\ \nuS^{2}\,(\#\Bc{S})^{2}\,\tl^{4}\,
(\bb^4\gq)^{-2j_0\sk\fE}\Bigl(\frac{\xri{S}{j}{\mathsf k}\,\xri{S}{j}{\mathsf k^{-1}}}{D_2}\Bigr)^{\sk+1}.\]
On the other side, by the hypothesis of~(b) the class gauge is
\m{\Nn^{\fE}D_2/(\xri{S}{j}{\mathsf k}\,\xri{S}{j}{\mathsf k^{-1}})}, so that by Definition~\ref{d:cs} and the formula for
\m{\cw{S}{0}}
\begin{align*}
\csize{S}{j}{\cgauge{S}{j}{c}{c'}{\gamma}}^{2}={}&\frac{\tl^{4}\,\mathsf D_j}{\#\Sbar\cdot\#\Cl{S}\cdot\Nn^{\fE}\cdot\gq^{2j_0\sk\fE}}\\
&\cdot\Bigl(\frac{\xri{S}{j}{\mathsf k}\,\xri{S}{j}{\mathsf k^{-1}}}{D_2}\Bigr)^{\sk+1}.
\end{align*}
Comparing, it remains to see that
\m{\Nn^{\fE}\cdot\dd\cdot\#\Bc{S}\le\mathsf D_j\cdot\bb^{8j_0\sk\fE}}. Now \m{\Nn^{\fE}\le\bb^{60\sn\fE}\le\tl} by \inp{A2}
and~\inp{S1}, \m{\#\Bc{S}\le\bb^{2\sk\fE}=\tl^{8}} by \inp{A4}, and \m{\bb^{8j_0\sk\fE}=\tl^{32j_0}}. If
\m{j\le\sro} then \m{j_0=j} and \m{\mathsf D_j\,\tl^{32j}=\tl^{13}\dd\ge\tl^{9}\dd}; if \m{j=\sro+1} then
\m{\mathsf D_j=\dm\ge\tl^{44}\dd} by~\inp{A3}.
\end{proof}
\leanline{\lean{tsumamplitudeAbsclassle}}

\sect{Counting: discriminants, shared vectors, pairs}{s:count}
In this section a power \m{x^{-l}} of a real number \m{x\ge0}, as a term of a sum in \m{[0,\infty]}, is read as $0$
when \m{x=0}, as in the formal text; this concerns the discriminants \m{\latd{S}{c}(U)} and the indices.

\begin{proposition}[inverse discriminants, summed over classes and subspaces]\label{p:latdisc}
In Setting~\ref{set:data}, for every \m{m\in\NN},
\[\sum_{c\in\Cl{S}}\ \sum_{\substack{U\subseteq E^\kappa\\ \dim_EU+m=\sk}}\latd{S}{c}(U)^{-(\sk+1)}\ \le\
\#\Cl{S}\cdot\Nn^{-\sk\fE}\cdot\tl^{\,m+1},\]
the inner sum over the $E$-subspaces $U$ of \m{E^\kappa} of codimension~$m$.
\end{proposition}
\begin{proof}
Put \m{B_0=\tl^{m+1}\Nn^{-\sk\fE}} and, for a class~$c$, let \m{\Phi(c)\in[0,\infty]} be the inner sum. Let
\m{K'=\Kf{E}{\kappa}\cap\incH{S}^{-1}(\KS{S})} and let $\mathfrak A$ be the set of the left cosets of $K'$ in
\m{\Kf{E}{\kappa}}; by \inp{A5} it is finite, with \m{\#\mathfrak A\ge(\Nn^{\sk}/\bb)^{\fE}}. Fix a representative
\m{\tilde a} of each \m{a\in\mathfrak A}, and write \m{[x]\in\Cl{S}} for the class of \m{x\in\Gf{E}{\kappa}}.

\emph{Claim.} For every \m{x\in\Gf{E}{\kappa}}: \m{\sum_{a\in\mathfrak A}\Phi([x\tilde a])\le\#\mathfrak A\cdot B_0}.

Let \m{a\in\mathfrak A}. The representative of the class \m{[x\tilde a]} is \m{\gamma_a\,x\,\tilde a\,\mathsf k_a} with
$\gamma_a$ rational and \m{\mathsf k_a\in\KcolS{S}}. Put \m{y_a=\tilde a\,\mathsf k_a}. Since \m{\KcolS{S}} is
contained in \m{\Kf{E}{\kappa}} and, by \inp{A1} and \m{\KgS{S}\subseteq\KS{S}}, in $K'$, the element $y_a$ lies in
\m{\Kf{E}{\kappa}} and represents~$a$. Let $r_a$ be the row $k_0$ of \m{y_a^{-1}}, an integral row vector. By
\inp{S8}, \m{\clCX{S}{[x\tilde a]}=\mbc{\fs}{\piS{S}}\,y_a^{-1}x^{-1}\gamma_a^{-1}}.

Let $U$ be a subspace of codimension~$m$, and \m{M=\fadm{x_0}\,w_0} the basis matrix of its lattice for the
class \m{[x\tilde a]}, with $x_0$ a basis matrix of $U$ over~$E$ and $w_0$ invertible, so that
\m{\latd{S}{[x\tilde a]}(U)^2=D_1} with \m{D_1=\gix{\det\fadm{x_0^*x_0}}\,\gix{\det w_0}^2}. If \m{D_1=0} the term of $U$
is~$0$. Otherwise, by \inp{S8} the integrality of \m{\clCX{S}{[x\tilde a]}M} means that
\m{y_a^{-1}x^{-1}\gamma_a^{-1}M} is integral and that \m{r_a\,x^{-1}\gamma_a^{-1}M} has entries in \m{\piS{S}\hO}; as
$y_a$ is integral, also \m{x^{-1}\gamma_a^{-1}M} is integral. We apply the second part of \inp{S15} to
\m{H=x^{-1}}, the subspace \m{U'=\gamma_a^{-1}U}, the matrix \m{\gamma_a^{-1}x_0}, which has the same Gram matrix
\m{x_0^*x_0} because $\gamma_a$ is unitary, the matrix $w_0$, \m{Y_0=r_a} and \m{K_0=\piS{S}\hO}. It gives
\begin{gather*}
\latd{S}{[x\tilde a]}(U)^{-(\sk+1)}\ \le\ \bigl(\sld{x^{-1}}{U'}\cdot i_a(U')\bigr)^{-(\sk+1)},\\
i_a(U')=\bigl[\hO\langle r_a\,x^{-1}M_{U'}\rangle:\piS{S}\hO\bigr]\ge1,
\end{gather*}
where \m{M_{U'}=\slb{x^{-1}}{U'}}. Since \m{U\mapsto\gamma_a^{-1}U} is injective and preserves the dimension, summing
over $U$ and then over~$a$,
\begin{equation}\label{e:ld1}
\sum_{a}\Phi([x\tilde a])\ \le\ \sum_{U'}\sld{x^{-1}}{U'}^{-(\sk+1)}\cdot\sum_{a\in\mathfrak A}i_a(U')^{-(\sk+1)},
\end{equation}
the outer sum over the subspaces $U'$ of codimension~$m$.
Fix $U'$. We apply the congruence count \inp{R5} to the set $\mathfrak A$ with \m{d_0=\piS{S}}, a set $I_1$ with one
element, \m{I_2=\kappa}, $I_3$ the set of the \m{\sk-m} columns of \m{M_{U'}}, \m{m_1=m}, the integral matrices
\m{Y_a=r_a}, the matrix \m{M_0=x^{-1}M_{U'}}, which is integral with an integral left inverse, \m{K_0=\piS{S}\hO},
\m{\eta_a=i_a(U')} and \m{l_1=\sk+1}. Here \m{C_1=1} is admissible: if all entries of \m{r_a-r_{a'}} lie in
\m{\piS{S}\hO}, then so do those of \m{(r_a-r_{a'})\,y_{a'}}, which is the row $k_0$ of \m{y_a^{-1}y_{a'}-1}; by
\inp{S8} this means \m{y_a^{-1}y_{a'}\in K'}, that is \m{a=a'}. Hence
\m{\sum_a i_a(U')^{-(\sk+1)}\le\gix{\piS{S}}^{1+m}}, and \m{\gix{\piS{S}}=\Nn^{\fE}} because
\m{\gix{\piS{S}}^2=\gix{\Nn}=\Nn^{2\fE}}. With the first part of \inp{S15}, \eqref{e:ld1} gives
\begin{align*}
\sum_{a}\Phi([x\tilde a])\ &\le\ \Nn^{\fE(m+1)}\cdot\bigl(200\,\fE\,\sk\,\rdE{E}\bigr)^{\fE(\sk+1)m}\\
&\le\ \bb^{\,60\sn\fE(m+1)+\fE(\sk+1)m/1000},
\end{align*}
by \inp{A2} and Setting~\ref{set:S}. On the other hand
\begin{gather*}
\#\mathfrak A\cdot B_0\ \ge\ \bb^{-\fE}\,\tl^{m+1}=\bb^{\,\sk\fE(m+1)/4-\fE}\qquad\text{and}\\
60\,\sn\,(m+1)+\frac{(\sk+1)\,m}{1000}+1\ \le\ \frac{\sk\,(m+1)}{4},
\end{gather*}
because \m{60\sn+1\le\sk/8} and \m{(\sk+1)/1000\le\sk/8}. This proves the claim.

\emph{Averaging.} Every class is \m{[x\tilde a]} for suitable $x$ and~$a$, so the claim shows that $\Phi$ has
finite values. Let \m{F_0(x)=B_0-\Phi([x])}, a real function on \m{\Gf{E}{\kappa}} that is invariant under the
rational points on the left and under \m{\KcolS{S}} on the right. By the claim,
\m{\sum_{a}F_0(x\tilde a)\ge0} for all~$x$. By the properties \inp{L2} of the class integral,
\[0\ \le\ \int_{\Gamma_0\backslash}\sum_{a\in\mathfrak A}F_0(\,\cdot\,\tilde a)=\#\mathfrak A\cdot\int_{\Gamma_0\backslash}F_0
=\#\mathfrak A\cdot\sum_{c\in\Cl{S}}\bigl(B_0-\Phi(c)\bigr),\]
which is the proposition.
\end{proof}
\leanline{\lean{sumlatticeDiscinvle}}

\begin{lemma}[shared vectors are short]\label{l:short}
In Setting~\ref{set:S}, let $S$ be admissible, \m{j\in\{0,\dots,\sro+1\}}, and \m{\gamma\in\GS{S}{E}} with
\m{\IsNr{S}{\gamma}} and \m{\HasLev{S}{j}{\gamma}}. Let \m{u\in\UoE}, let \m{a'\le\sro} with \m{j\le a'}, and
\m{Y=\padl{j}{\Zf{S}{\gamma}}}. Then for every infinite place $\tau$ and every \m{a<\sro}, the coordinate at $e_a$
of the vector \m{t=(u\gamma)\mathfrak f\cdot\adj(Y)_{\cdot,a'}}, where \m{\adj(Y)_{\cdot,a'}=\det Y\cdot(Y^{-1})_{\cdot,a'}} is the column
$a'$ of the adjugate, satisfies
\[2\,|(t_{e_a})_\tau|\ <\ \qS{S}{0}^{\,\sn^4}.\]
\end{lemma}
\begin{proof}
By \inp{B11} there is \m{v\in E^{\sn}} with \m{\mathfrak f\,v=\gamma\mathfrak f\,(Y^{-1})_{\cdot,a'}}. As $u$ fixes the
frame, \m{t=\det Y\cdot u\,\mathfrak f\,v=\det Y\cdot\mathfrak f\,v}, and \m{t_{e_a}=\det Y\cdot v_a}. Fix~$\tau$.

Since $\gamma$ is near, there is \m{y\in\bxp{S}} with \m{\ri{\gamma}^{-1}y\in\bxp{S}}. By \inp{R4} there is
\m{u_1\in\Uone} with \m{\|u_1y_\tau\|\le\bb^{6\sn+2}}. Let $G'$ and $G''$ be the complex square matrices of size
\m{\sn} with the entries
\[G'_{l,a''}=\frac{\bscale{S}{\tau}{\fr{l}}}{\bscale{S}{\tau}{\fr{a''}}}\,(u_1)_{\fr{l},\fr{a''}},\qquad
G''_{l,a''}=\frac{\bscale{S}{\tau}{\fr{l}}}{\bscale{S}{\tau}{\fr{a''}}}\,(u_1\,\ri{\gamma}_\tau)_{\fr{l},\fr{a''}}.\]
The block of $u_1$ on the frame is upper triangular with diagonal entries~$1$ \inp{L8}, so \m{\det G'=1}. All
entries of $G'$ and $G''$ have absolute value at most \m{\bb^{5\sn^2}}: the entry of \m{u_1} or \m{u_1\ri{\gamma}_\tau}
is a coordinate of \m{(u_1y_\tau)\,y'^{-1}\fr{a''}} for \m{y'=y_\tau} or \m{y'=(\ri{\gamma}^{-1}y)_\tau}, hence at most
\m{\bb^{6\sn+2}\cdot\bxs{S}=\bb^{6\sn+2+2\sn^2}} because both points are boxed, and the quotient of scales is a
power of \m{\kapS{S}{\tau}} with exponent between \m{-\sro} and~\m{\sro}, at most \m{\bb^{23\sro}} by~\inp{A2}.
Conjugating the relation \m{\mathfrak f\,v=\gamma\mathfrak f\,(Y^{-1})_{\cdot,a'}} at $\tau$ by the scales, applying
$u_1$ and reading the coordinates at the frame gives \m{G'\,v_\tau=G''\,((Y^{-1})_{\cdot,a'})_\tau}. Therefore the
square matrix of size \m{2\sn}
\[\mathfrak M=\begin{pmatrix}G'&-G''\\0&Y_\tau\end{pmatrix}\]
maps the vector with the blocks \m{v_\tau} and \m{((Y^{-1})_{\cdot,a'})_\tau} to the vector with the blocks $0$ and the
unit vector of index~$a'$, and \m{\det\mathfrak M=\det Y_\tau}. By Cramer's rule, \m{(\det Y\cdot v_a)_\tau} is
the determinant of the matrix obtained from $\mathfrak M$ by replacing its column $a$ with that vector. All
entries of this matrix have absolute value at most \m{\bb^{5\sn^2}}: those of $Y_\tau$ are $0$, $1$ or entries of
\m{\Zf{S}{\gamma}} at~$\tau$, bounded by~\inp{S5}. By \inp{L7}, with \m{2\le\bb} and \m{\sn\le\bb},
\begin{align*}
2\,|(\det Y\cdot v_a)_\tau|\ &\le\ 2\,\bigl((2\sn)^2\,\bb^{5\sn^2}\bigr)^{2\sn}\ \le\ \bb^{\,10\sn^3+8\sn+1}\\
&<\ \bb^{10\sn^4}\ \le\ \qS{S}{0}^{\,\sn^4}.\qedhere
\end{align*}
\end{proof}
\leanline{\lean{normsharedvectorlt}}

\begin{proposition}[a weighted count of labelled terms]\label{p:count}
In Setting~\ref{set:data}, let \m{c\in\Cl{S}}, \m{j\in\{0,\dots,\sro+1\}}, \m{\mathsf t\in\Sbar} and \m{m\in\NN}; let
$\mathcal P$ be a set of pairs \m{p\in\Cl{S}\times\GS{S}{E}} such that \m{\IsLab{S}{j}{c}{\mathsf t}{p}} for all
\m{p\in\mathcal P}; let \m{U\subseteq E^\kappa} be an $E$-subspace with \m{\sk\le\dim_EU+m}; and let
\m{\eta\colon\GFin{S}\to\NN} be a function such that \m{\fvi{S}{c}{\telem{S}{c}{p}}{\gid{S}}{U}} divides
\m{\eta(\telem{S}{c}{p})} for every \m{p\in\mathcal P}. Then
\[\sum_{p\in\mathcal P}\ipw{S}{c}{\eta}{p}\ \le\ \bb^{\,2000\,\sn^6\,(m+1)\,\fE}.\]
\end{proposition}
\begin{proof}
Write \m{m'=\sk-\dim_EU\le m}. For \m{p=(c',\gamma)\in\mathcal P} the element $\gamma$ is near, has level~$j$, and
lies in \m{\terms{S}{c}{s}{c'}} for some~$s$.

\emph{Step 1: few frame Gram matrices.} By \inp{S6} and \inp{S5} the entries of \m{\Zf{S}{\gamma}}, for
\m{(c',\gamma)\in\mathcal P}, are algebraic integers of absolute value at most \m{\bb^{5\sn^2}} at every place; so by
\inp{B9}\,(iii) the map \m{(c',\gamma)\mapsto\Zf{S}{\gamma}} takes at most \m{N_1=(1+2\bb^{5\sn^2})^{2\sn^2\fE}}
values. By \inp{B9}\,(i) it suffices to bound the sum over a non-empty subset \m{\mathcal P'} on which
\m{\Zf{S}{\gamma}} is a fixed matrix~$Z$. Fix \m{\mathcal P'} and~$Z$.

\emph{Step 2: the modulus.} Let \m{Y=\padl{j}{Z}}; its entries are algebraic integers, and
\m{\det Y=\lmin{j}{Z}\ne0}. Put \m{m_0=\qS{S}{0}^{\,\sn^4}} and \m{d_0=m_0\cdot\gq\cdot\det(YY^*)}, a non-zero algebraic
integer of~$E$ with \m{\cj{d_0}=d_0}. By \inp{S5},
\m{|(d_0)_\tau|\le\bb^{11\sn^4}\cdot\bb^{23}\cdot\bb^{11\sn^3}}, so
\begin{equation}\label{e:cnt0}
\gix{d_0}=\prod_\tau|(d_0)_\tau|^2\ \le\ \bb^{(22\sn^4+22\sn^3+46)\fE}.
\end{equation}

\emph{Step 3: the congruence count.} Let \m{\mathcal R=\hO\langle\fadm{\xrows{Z}}^*\rangle}; by \inp{S6} it is
\m{\mathcal R(\telem{S}{c}{p})} for every \m{p\in\mathcal P'}. We apply \inp{R5} to \m{\mathcal P'} with $d_0$; with
\m{I_1=\{0,\dots,\sro\}}, \m{I_2=\Mx{\fs}}, $I_3$ the set of the columns of the basis matrix \m{M_U} of the lattice
of $U$ and~$c$, so that \m{m_1=m'+1}; with \m{Y_p=(\fmid{\telem{S}{c}{p}})^*}, integral by~\inp{S6};
\m{M_0=\Tl\,\clCX{S}{c}\,M_U}, integral with an integral left inverse; \m{K_0=\gq\,\mathcal R+d_0\hO^{I_1}};
\m{\eta_p=\eta(\telem{S}{c}{p})}; and \m{l_1=\sk}. The divisibility holds: \m{\hO\langle Y_pM_0\rangle} is the lattice
\m{\Lambda_U(\telem{S}{c}{p})}, and its index relative to $K_0$ divides its index relative to the smaller module
\m{\gq\,\mathcal R}, which is the frame value index of the hypothesis. So
\[\sum_{p\in\mathcal P'}\ipw{S}{c}{\eta}{p}\ \le\ \gix{d_0}^{\,\sn(\sn+m'+1)}\cdot C_2,\]
if every non-empty \m{\mathcal B\subseteq\mathcal P'} such that, for \m{p,p'\in\mathcal B}, the conjugates of the middle
rows of \m{\fmid{\telem{S}{c}{p}}-\fmid{\telem{S}{c}{p'}}} lie in~$K_0$, has at most $C_2$ elements. We show this for
\m{C_2=\gix{d_0}^{\sn\sro}}. Fix such a $\mathcal B$ and an element \m{p_1=(c_1,\gamma_1)} of it.

\emph{Step 4: equal middle rows.} Let \m{p=(c',\gamma)\in\mathcal B}. The module \m{d_0\hO^{I_1}} consists of
vectors $v$ with \m{\gq^{-1}\fadm{Y^{-1}}^*v} integral, because \m{\gq^{-1}(Y^{-1})^*d_0=m_0\det Y\cdot\adj(Y)^*}. So
\inp{R8} applies to the pair \m{\gamma,\gamma_1}, with $\Sigma$ the set of all \m{a<\sro}: there is a matrix
\m{\ell_p} over $E$ with rows in $\kappa$ and \m{\sro} columns such that \m{\clCX{S}{c}\fadm{\ell_p}} is integral
and \m{\gamma_1\mathfrak f-\gamma\mathfrak f} has the rows \m{\gq\,\ell_p\,\xrows{Z}} on $\kappa$ and vanishes on the row
of $z$ and the rows~$e^*_a$. Apply \inp{B10} with \m{I_1=\{0,\dots,\sro-1\}}, $p_1$ the identity, \m{\Psi=0} and
$U_1$ the matrix with the rows \m{-\ell_p} on $\kappa$ and zero rows elsewhere: its columns are orthogonal to the
frame and \m{\clC{S}{c}\fadm{U_1}} is integral \inp{B12}. We get \m{u_p\in\clUo{S}{c}} with
\m{u_p\,t'=t'-\gq\sum_{a<\sro}(U_1)_{\cdot,a}\,h(e_a,t')+(\text{a combination of the }e_a)} for all vectors~$t'$.
For \m{t'=\gamma\fr{a'}} we have \m{h(e_a,t')=Z_{a,a'}}; so \m{(u_p\gamma)\mathfrak f} and \m{\gamma_1\mathfrak f} have the
same rows outside the rows~$e_a$. Put \m{\gamma_p=u_p\gamma}; it lies over the same triple as $\gamma$ \inp{S6} and
has the frame Gram matrix~$Z$ \inp{L8}.

\emph{Step 5: colouring by the first rows.} By \inp{S6} and \inp{B12} the rows $e_a$ of the integral matrix
\m{\pos{S}{c}{c'}{\gamma_p}\,\mathfrak f} are those of \m{\gamma_p\mathfrak f}. We colour $\mathcal B$ by these \m{\sro}
rows modulo \m{d_0\hO^{I_1}}: by \inp{B9}\,(ii) and~(iv), \m{\#\mathcal B\le\gix{d_0}^{\sn\sro}}, provided that two
elements \m{p=(c',\gamma)} and \m{p'=(c'',\gamma')} of $\mathcal B$ are equal whenever the matrix $\Upsilon$ of the
rows $e_a$ of \m{\gamma_{p'}\mathfrak f-\gamma_p\mathfrak f} has all its entries in \m{d_0\hO}. Let \m{p,p'} be such, and
\m{\Upsilon=d_0\,H}; the entries of $H$ are algebraic integers.

Both \m{\gamma_p\mathfrak f} and \m{\gamma_{p'}\mathfrak f} have the Gram matrix of the frame, and they differ by
$\Upsilon$ in the isotropic rows $e_a$ only. Expanding, and using \m{h(e_a,\gamma_p\fr{a'})=Z_{a,a'}}, we get
\m{\Upsilon^*\xrows{Z}+(\xrows{Z})^*\Upsilon=0}, that is
\begin{equation}\label{e:cnt1}
H^*\,\xrows{Z}=-(\xrows{Z})^*\,H.
\end{equation}
Let $\alpha$ be the matrix of the first \m{\sro} columns of \m{\adj(Y)}. For \m{a'\ge j} the column $a'$ of
\m{H\adj(Y)} vanishes: its entries, multiplied by~$d_0$, are differences of two coordinates as in
Lemma~\ref{l:short}, applied to \m{\gamma'} with \m{u_{p'}} and to $\gamma$ with~$u_p$, so their absolute values
at every place are less than~$m_0$; dividing by $m_0$ we see that \m{\gq\det(YY^*)\,(H\adj(Y))_{a,a'}} is an
algebraic integer with all absolute values less than~$1$, hence zero by~\inp{L6}. Consequently
\begin{equation}\label{e:cnt2}
H\alpha\,\xrows{Z}=\det Y\cdot H\ \ \text{if }j\le\sro,\qquad\xrows{Z}\,\alpha=\det Y\cdot1\ \ \text{if }j=\sro+1:
\end{equation}
in the first case, the columns of \m{H\adj(Y)} of index \m{a\ge j} vanish and the rows of index \m{a<j} of
\m{\xrows{Z}} are those of~$Y$, so \m{H\alpha\,\xrows{Z}=H\adj(Y)\,Y}; in the second, \m{Y=Z}. Put
\m{\Theta=H\alpha} and
\[\Psi=m_0\cdot\bigl(\cj{\det Y}\cdot\Theta-\det Y\cdot\Theta^*-(\xrows{Z}\alpha)^*\,\Theta\bigr).\]
Its entries are algebraic integers, and \m{\Psi+\Psi^*=-m_0\,\alpha^*\bigl((\xrows{Z})^*H+H^*\xrows{Z}\bigr)\alpha=0}
by~\eqref{e:cnt1}. Moreover \m{\Theta^*\xrows{Z}=-(\xrows{Z}\alpha)^*H} by \eqref{e:cnt1}, and with \eqref{e:cnt2}
one finds in both cases \m{\Psi\,\xrows{Z}=m_0\det(YY^*)\,H}, that is \m{\gq\,\Psi\,\xrows{Z}=\Upsilon}. Now apply
\inp{B10} with \m{U_1=0} and this~$\Psi$: the element \m{r=1+\gq\,x_1\Psi x_1^*J} lies in \m{\clUo{S}{c}}, and it adds
to a vector $t'$ the combination \m{\gq\sum_{a,a''}e_a\Psi_{a,a''}h(e_{a''},t')}. For \m{t'=\gamma_p\fr{a'}} this is the
column $a'$ of \m{\gq\,\Psi\,\xrows{Z}=\Upsilon}, placed in the rows~$e_a$. Hence
\m{(r\,u_p\,\gamma)\mathfrak f=(u_{p'}\gamma')\mathfrak f}, and \m{p=p'} by~\inp{R7}.

\emph{Step 6: the exponents.} Altogether the sum over $\mathcal P$ is at most
\m{N_1\cdot\gix{d_0}^{\sn(\sn+m'+1)+\sn\sro}=N_1\cdot\gix{d_0}^{\sn(2\sn+m')}}. Here
\m{N_1\le\bb^{(5\sn^2+2)\,2\sn^2\fE}\le\bb^{12\sn^4\fE}}, and by \eqref{e:cnt0}, \m{m'\le m} and
\m{2\sn+m\le2\sn(m+1)}, the second factor is at most \m{\bb^{23\sn^4\cdot2\sn^2(m+1)\fE}}. The sum of the exponents is
at most \m{2000\,\sn^6(m+1)\fE}.
\end{proof}
\leanline{\lean{tsuminvPowWeightleofcodimle}}

\begin{lemma}\label{l:gindex}
In Setting~\ref{set:S}, let $I_1$ be a finite set and \m{M_1\subseteq M_2} two \m{\hO}-submodules of
\m{\AEf^{I_1}} with \m{[M_2:M_1]\ge1}. Then \m{[M_2:\gq M_2]=[M_1:\gq M_1]}.
\end{lemma}
\begin{proof}
Indices are multiplicative in the towers \m{\gq M_1\subseteq M_1\subseteq M_2} and
\m{\gq M_1\subseteq\gq M_2\subseteq M_2}, and \m{[\gq M_2:\gq M_1]=[M_2:M_1]} because multiplication by the idele
\m{\gid{S}} is bijective. So \m{[M_2:M_1]\cdot[M_1:\gq M_1]=[M_2:\gq M_2]\cdot[M_2:M_1]}, and \m{[M_2:M_1]} can be
cancelled.
\end{proof}
\leanline{\lean{submoduleIndexgIdelesmul}}

\begin{proposition}[the pair weights are summable]\label{p:pairs}
In Setting~\ref{set:data}, for every \m{j\in\{0,\dots,\sro+1\}},
\[\sum_{c\in\Cl{S}}\ \sum_{c'\in\Cl{S}}\ \sum_{\gamma\in\GS{S}{E}}\ \sum_{\gamma'\in\GS{S}{E}}\pairw{S}{j}{c}{c'}{\gamma}{\gamma'}\ \le\ 2\,\tl^{2}.\]
\end{proposition}
\begin{proof}
\emph{Step 1: regrouping.} Let \m{\pairw{S}{j}{c}{c'}{\gamma}{\gamma'}\ne0}. Then $\gamma$ and $\gamma'$ lie in
\m{\terms{S}{c}{s}{c'}} for one $s$, are near representatives for \m{(c,c')} and have level~$j$. Let
\m{U=\clink{\gamma}{\gamma'}}, \m{i=\sk-\dim_EU}, and let \m{\mathsf t} be the component of $s$ in~\m{\Sbar}. Then
\m{q=(c',\gamma)} and \m{p=(c',\gamma')} satisfy \m{\IsLab{S}{j}{c}{\mathsf t}{q}} and \m{\IsLab{S}{j}{c}{\mathsf t}{p}},
and the pair weight is \m{\termw{S}{j}{i}{c}{q}{p}{U}}. For a labelled $q=(c',\gamma)$ and a subspace $U$ let
\m{\mathcal P(q,U)} be the set of the \m{p=(c'',\gamma')} with \m{\IsLab{S}{j}{c}{\mathsf t}{p}},
\m{\Cong{S}{c}{c''}{\gamma}{\gamma'}} and \m{\clink{\gamma}{\gamma'}=U}. The left side of the proposition is then at
most
\begin{align*}
\sum_{i\in\NN}\ \sum_{\mathsf t\in\Sbar}\ \sum_{c}\ \sum_{\dim_EU+i=\sk}\ &\frac{\latd{S}{c}(U)^{-(\sk+1)}}{\cw{S}{i}}\cdot\Sigma_1(i,\mathsf t,c,U),\\
\Sigma_1(i,\mathsf t,c,U)={}&\sum_{q}\Bigl(\ipw{S}{c}{\eta_{j,U}}{q}\cdot\sum_{p\in\mathcal P(q,U)}\ipw{S}{c}{\sati{S}{\cdot}}{p}\Bigr),
\end{align*}
the sum over $q$ running over the pairs with \m{\IsLab{S}{j}{c}{\mathsf t}{q}}.

\emph{Step 2: the two inner sums.} Fix $i$, \m{\mathsf t}, $c$, $U$ and \m{q=(c',\gamma)}. Let $U'$ be the subspace
of the \m{t_0\in U} that are orthogonal, for the standard form of \m{E^\kappa}, to the $\kappa$-parts of all
vectors \m{\gamma\fr{a}}, \m{a\le\sro}. Then \m{\dim_EU'\ge\dim_EU-\sn}, so \m{\sk\le\dim_EU'+i+\sro+1}. We show in
Step~3 that \m{\fvi{S}{c}{\telem{S}{c}{p}}{\gid{S}}{U'}} divides \m{\sati{S}{\telem{S}{c}{p}}} for all
\m{p\in\mathcal P(q,U)}, and in Step~4 that \m{\fvi{S}{c}{\mathsf k}{\gid{S}}{U}} divides \m{\eta_{j,U}(\mathsf k)} for
\m{\mathsf k=\telem{S}{c}{q}}. Proposition~\ref{p:count}, applied to \m{\mathcal P(q,U)} with $U'$ and \m{m=i+\sro+1},
and to the set of all labelled $q$ with $U$ and \m{m=i}, then gives
\[\Sigma_1(i,\mathsf t,c,U)\ \le\ \bb^{\,2000\,\sn^6(2i+\sro+3)\fE}\ \le\ \tl^{\,i+1},\]
the last by \inp{S1}, since \m{2i+\sro+3\le(i+1)(\sro+3)} and \m{2000\,\sn^6(\sro+3)\le10^5\sn^9}. By
Proposition~\ref{p:latdisc}, the sum over \m{\mathsf t}, $c$ and $U$ of \m{\latd{S}{c}(U)^{-(\sk+1)}} is at most
\m{\#\Sbar\cdot\#\Cl{S}\cdot\Nn^{-\sk\fE}\cdot\tl^{i+1}=\cw{S}{i}\,(2\tl^2)^{-i}\,\tl^{i+1}}. So the terms with a given $i$
contribute at most \m{(2\tl^2)^{-i}\,\tl^{2i+2}=2^{-i}\tl^2}, and the sum over $i$ is at most \m{2\,\tl^2}.

\emph{Step 3: the saturation index.} Let \m{p=(c'',\gamma')\in\mathcal P(q,U)}; so \m{\gamma,\gamma'\in\terms{S}{c}{s}{c''}} for
some~$s$, and \m{U=\clink{\gamma}{\gamma'}}. Put \m{\mathsf k=\pos{S}{c}{c''}{\gamma}} and
\m{\mathsf k'=\pos{S}{c}{c''}{\gamma'}=\telem{S}{c}{p}}; then \m{\mathsf k'\mathsf k^{-1}\in\KgS{S}}. Let
\m{M_{U'}=\fadm{p_0}\,a_0} be the basis matrix of the lattice of $U'$ and~$c$, with $p_0$ a basis matrix of $U'$
over~$E$, and let $P_0$ be the matrix over $E$ with the rows $p_0$ on $\kappa$ and zero rows elsewhere.

\emph{(i)} Every column of \m{\gamma'\gamma^{-1}P_0-P_0} is a combination of the vectors $e_a$ and \m{\gamma'e_a},
\m{a<\sro}. Indeed, let $t$ be a column of $P_0$, with $\kappa$-part \m{t_0\in U'}. As \m{t_0\in\clink{\gamma}{\gamma'}}, there is
\m{t_2\in E^\kappa} such that the pair \m{((t_0,0),(t_2,0))}, the zeros being the $z$-coordinates, lies in
\m{\link{\gamma}\cap\link{\gamma'}}. By \inp{S11} for \m{\gamma'}, $t$ is the sum of \m{\gamma'\tilde t_2}, where
\m{\tilde t_2} has the $\kappa$-part $t_2$ and no other coordinates, and of a combination of the $e_a$ and
the~\m{\gamma'e_a}. On the other hand $t$ is orthogonal to all \m{\gamma\fr{a}}, by the definition of~$U'$; so
\m{\gamma^{-1}t} is orthogonal to the frame, has no $e^*$-coordinates and no $z$-coordinate; let $t_3$ be its
$\kappa$-part. By \inp{S11} for $\gamma$ the pair \m{((t_0,0),(t_3,0))} lies in \m{\link{\gamma}}, hence so does
\m{((0,0),(t_3-t_2,0))}; as the link is Lagrangian, \m{(t_3-t_2)^*(t_3-t_2)=0}, and \m{t_3=t_2} by the last
statement of~\inp{S11}. So \m{\gamma^{-1}t} is \m{\tilde t_2} plus a combination of the~$e_a$, and
\m{\gamma'\gamma^{-1}t-t} is a combination of the $e_a$ and the~\m{\gamma'e_a}.

\emph{(ii)} We apply \inp{L11} to the map \m{P_1\colon t'\mapsto\gq\cdot(h(\mathsf k'\fr{a},t'))_{a\le\sro}} on
\m{\AEf^{\Ix{\fs}}}, with \m{h=h_{\Tl}}, the modules \m{M_2=\hO^{\Ix{\fs}}\cap\operatorname{span}_{\AEf}(e_a,\mathsf k'e_a)},
\m{D_2=\hO\langle e_a,\mathsf k'e_a\rangle}, \m{K_0=\gq\,\mathcal R(\mathsf k')} and \m{K_1=\Lambda_{U'}(\mathsf k')}. First,
\m{P_1(e_a)} is $\gq$ times the conjugate of the row $a$ of \m{\Zl{\mathsf k'}}, and
\m{P_1(\mathsf k'e_a)=\gq\,(h(\fr{a'},e_a))_{a'}=0}; so \m{D_2\subseteq P_1^{-1}(K_0)}. Second, a general element of
\m{\Lambda_{U'}(\mathsf k')} is \m{(h(\mathsf k'\fr{a},\tilde t))_a} with
\m{\tilde t=\clC{S}{c}\,\fadm{P_0}\,a_0\,r_0} for an integral vector~$r_0$; the vector \m{\tilde t} is integral. Put
\m{\mu_0=-\gq^{-1}(\mathsf k'\mathsf k^{-1}-1)\,\tilde t}. It is integral by~\inp{S7}. It lies in the span of the $e_a$ and
\m{\mathsf k'e_a}, by~(i) and~\inp{B12}, because \m{(\mathsf k'\mathsf k^{-1}-1)\,\clC{S}{c}\,\fadm{P_0}=\clC{S}{c}\,\fadm{\gamma'\gamma^{-1}P_0-P_0}}.
And, \m{\mathsf k'\mathsf k^{-1}} being an isometry,
\[P_1(\mu_0)=\bigl(h(\mathsf k'\fr{a},\tilde t)\bigr)_a-\bigl(h(\mathsf k\fr{a},\tilde t)\bigr)_a=\bigl(h(\mathsf k'\fr{a},\tilde t)\bigr)_a,\]
since \m{h(\mathsf k\fr{a},\tilde t)} is the image in \m{\AEf} of \m{(\gamma\fr{a})^*\JT{\TE}P_0}, which is zero, times
\m{a_0r_0} \inp{B12}. So \m{K_1\subseteq P_1(M_2)}, and \inp{L11} shows that
\m{[K_1:K_0]=\fvi{S}{c}{\mathsf k'}{\gid{S}}{U'}} divides \m{[M_2:D_2]=\sati{S}{\mathsf k'}}.

\emph{Step 4: the index for~$q$.} Let \m{\mathsf k=\telem{S}{c}{q}}, \m{\Lambda=\Lambda_U(\mathsf k)},
\m{\mathcal R=\mathcal R(\mathsf k)} and \m{\mathcal R_j=\mathcal R_j(\mathsf k)}. Then
\m{[\Lambda:\gq\mathcal R]=[\Lambda:\mathcal R]\cdot[\Lambda\cap\mathcal R:\gq\mathcal R]}, the second factor divides
\m{[\mathcal R:\gq\mathcal R]}, and \m{[\mathcal R:\gq\mathcal R]=[\mathcal R_j:\gq\mathcal R_j]} by Lemma~\ref{l:gindex},
whose hypothesis is \m{\xri{S}{j}{\mathsf k}\ge1} \inp{P1}. Finally \m{\mathcal R_j} is the image of \m{\hO^{\IdxB{j}}}
under a linear map, so \m{[\mathcal R_j:\gq\mathcal R_j]} divides
\m{[\hO^{\IdxB{j}}:\gq\hO^{\IdxB{j}}]=\gix{\gid{S}}^{\#\IdxB{j}}} \inp{B9}. Hence \m{\fvi{S}{c}{\mathsf k}{\gid{S}}{U}}
divides \m{\gix{\gid{S}}^{\#\IdxB{j}}\cdot\fvi{S}{c}{\mathsf k}{1}{U}=\eta_{j,U}(\mathsf k)}.
\end{proof}
\leanline{\lean{sumpairWeightle}}

\sect{Block bounds: the proof of Theorem B}{s:blocks}
\begin{proof}[Proof of Theorem B\,(a)]
Write \m{\omega=\omega_{S,c'}} for the box partition of the class~$c'$ \inp{B1}.

\emph{Step 1: a choice of representatives.} Call \m{\gamma_1,\gamma_2\in\GS{S}{E}} equivalent if
\m{\gamma_2=u\gamma_1u'} with \m{u\in\clU{S}{c}} and \m{u'\in\clUo{S}{c'}}. In every equivalence class choose a
representative, and choose it with \m{\IsNRep{S}{c}{c'}{\cdot}} whenever the class contains such an element. Let
\m{\varphi(\gamma,y)=\omega(\ri{\gamma}^{-1}y)} if $\gamma$ is the chosen representative of its class, and
\m{\varphi(\gamma,y)=0} otherwise.

\emph{Step 2: the weight is normalized.} Let \m{\gamma\in\Xi\cap\terms{S}{c}{s}{c'}} and \m{y\in\Ginf}. We show
\[\sum_{(u,u')\in\clU{S}{c}\times\clU{S}{c'}}\varphi(u\gamma u'^{-1},\ri{u}y)=1.\] The term of
\m{(u,u')} is \m{\omega(\ri{u'}\,\ri{\gamma}^{-1}y)} if \m{u\gamma u'^{-1}} is a chosen representative, and $0$
otherwise. So the sum is \m{\sum_{u'}\alpha(u')\,\omega(\ri{u'}\,\ri{\gamma}^{-1}y)}, where \m{\alpha(u')} is the number
of the \m{u\in\clU{S}{c}} for which \m{u\gamma u'^{-1}} is a chosen representative, and by \inp{B1} it suffices to
see that \m{\sum_{u_0\in\clUo{S}{c'}}\alpha(u_0u')=1} for every \m{u'\in\clU{S}{c'}}. Put
\m{\gamma_0=\gamma u'^{-1}}; it lies in $\Xi$ because $\Xi$ is stable, so \m{\tdisc{S}{\gamma_0}\ne0}. We have to show
that there is exactly one pair \m{(u,u_0)\in\clU{S}{c}\times\clUo{S}{c'}} such that \m{u\gamma_0u_0^{-1}} is the chosen
representative of the class of~$\gamma_0$. There is one, by the definition of the equivalence. If
\m{u\gamma_0u_0^{-1}} and \m{\tilde u\gamma_0\tilde u_0^{-1}} are two and equal, then the element
\m{x=\tilde u^{-1}u} of \m{\UE} satisfies \m{x\gamma_0=\gamma_0\,\tilde u_0^{-1}u_0}, so $x$ fixes the vectors
\m{\gamma_0\fr{a}}. By \inp{S12}, $\gamma_0$ has
level \m{\sro+1}, so \m{\Zf{S}{\gamma_0}} is invertible, and the \m{\sro} vectors
\m{\gamma_0\mathfrak f\cdot(\Zf{S}{\gamma_0}^{-1})_{\cdot,a}}, \m{a<\sro}, are fixed by $x$ and have the identity matrix
as $e^*$-coordinates. By \inp{L10}, \m{x=1}; hence \m{u=\tilde u} and \m{u_0=\tilde u_0}.

\emph{Step 3: change of weight.} By Step~2 and \inp{B5}, for both sides the integrals
\m{\int\varphi(\gamma,y)\,\besker{\sigma}{S}{\gamma}{y}\,dy} converge, and for \m{\sigma=\cut} their sum over
\m{\gamma\in\Xi\cap\terms{S}{c}{s}{c'}} is \m{\ampsum{\cut}{S}{\Xi}{c}{s}{c'}}. So it suffices to show, for every
such~$\gamma$,
\begin{equation}\label{e:B1}
\Bigl\|\int\varphi(\gamma,y)\,\besker{\cut}{S}{\gamma}{y}\,dy\Bigr\|\ \le\
\begin{cases}\nuS\cdot\csize{S}{\sro+1}{\tdnorm{S}{\gamma}}&\text{if }\IsNRep{S}{c}{c'}{\gamma},\\ 0&\text{otherwise.}\end{cases}
\end{equation}

\emph{Step 4: representatives that are not near.} Suppose the integral in \eqref{e:B1} is not zero. Then $\gamma$
is the chosen representative of its class and there is $y$ with \m{\omega(\ri{\gamma}^{-1}y)\ne0} and
\m{\besker{\cut}{S}{\gamma}{y}\ne0}. By Lemma~\ref{l:incut} all \m{y_\tau} and all \m{(\ri{\gamma}^{-1}y)_\tau} are
in the cut; so \m{\ri{\gamma}^{-1}y\in\bxp{S}} by~\inp{B1}, and by \inp{B2} there is \m{u\in\clU{S}{c}} with
\m{\ri{u}y\in\bxp{S}}. Hence \m{u\gamma} is near. Its double coset \m{\clUo{S}{c}\,u\gamma\,\clUo{S}{c'}} has a
fixed representative, which is near by \inp{S3} and therefore a near representative for \m{(c,c')}; it is
equivalent to~$\gamma$. By the choice in Step~1, $\gamma$ itself is such a representative:
\m{\IsNRep{S}{c}{c'}{\gamma}}. This proves the second case of~\eqref{e:B1}.

\emph{Step 5: the open cell.} Let now $\gamma$ be the chosen representative of its class with
\m{\IsNRep{S}{c}{c'}{\gamma}}, and \m{x=\tdnorm{S}{\gamma}}. The element \m{\tdisc{S}{\gamma}} is a non-zero algebraic
integer \inp{S6}, so \m{x\ge1} by~\inp{L6}, and Lemma~\ref{l:cell} applies. Let \m{y^0\in\bxp{S}} with
\m{\ri{\gamma}^{-1}y^0\in\bxp{S}}. For each $\tau$ we have \m{\gdr{\Zr{\ri{\gamma}_\tau}}>0} by~\inp{S4}, and
\m{\bxs{S}\ge\bb}; so \inp{R3} applies to \m{\ri{\gamma}_\tau}, the point \m{y^0_\tau} and \m{\rho=\bxs{S}}, and gives
a compact set \m{\mathcal C_\tau}, a finite measure \m{n_\tau}, a probability space \m{(Q_\tau,m_\tau)} and a map
\m{t_\tau}. By \inp{L12} the mean-square bounds multiply. Applied to the bounded function
\m{\omega\le1}, this gives a point \m{q=(q_\tau)}, outside any given null set, with
\begin{equation}\label{e:B2}
\begin{aligned}
\Bigl\|\int_{\prod_\tau\mathcal C_\tau}\omega(y)\;\besker{\full}{S}{\gamma}{\ri{\gamma}\,t\,y}\,dy\Bigr\|^{2}\ &\le\ \prod_\tau n_\tau(\Gone{\fs})\\
&\le\ \prod_\tau\cellm{\ri{\gamma}_\tau}{\degt{S}{\tau}}{\bxs{S}}{\#\It{S}{\tau}},
\end{aligned}
\end{equation}
where \m{t=(t_\tau(q_\tau))_\tau\in\Ginf}; here we used that the kernel \m{\besker{\full}{S}{\gamma}{\ri{\gamma}ty}} is the
tensor product over $\tau$ of \m{\Wt{\full}{S}{\tau}(\ri{\gamma}_\tau t_\tau y_\tau)} and the conjugate of
\m{\Wt{\full}{S}{\tau}(t_\tau y_\tau)}, and that a probability measure gives a point at which a non-negative function is
at most its integral. We choose $q$ such that moreover every \m{t_\tau} lies in the subgroup generated by
\m{\Uoone} and \m{\Hone}, and such that every \m{y'\in\frbox{\bxs{S}}} with \m{\bsize{\ri{\gamma}_\tau t_\tau y'}\le\bxs{S}} lies
in~\m{\mathcal C_\tau}. By Lemma~\ref{l:cell} the right side of \eqref{e:B2} is at most
\m{(\nuS\,\csize{S}{\sro+1}{x}/2)^2}.

\emph{Step 6: from the cut to the full kernel.} Substituting \m{y\mapsto\ri{\gamma}\,t\,y} and using the invariance of
$\omega$ under $t$ \inp{B1},
\[\int\varphi(\gamma,y)\,\besker{\cut}{S}{\gamma}{y}\,dy=\int\omega(y)\,\besker{\cut}{S}{\gamma}{\ri{\gamma}\,t\,y}\,dy.\]
The integrand vanishes outside \m{\prod_\tau\mathcal C_\tau}: if it is not zero at~$y$, then by Lemma~\ref{l:incut}
\m{\ri{\gamma}_\tau t_\tau y_\tau} and \m{t_\tau y_\tau} are in the cut; as \m{\omega(ty)=\omega(y)\ne0}, the point
\m{t_\tau y_\tau} lies in \m{\frbox{\bxs{S}}} \inp{B1}, hence so does \m{y_\tau}, because $t_\tau$ fixes the frame vectors;
and \m{\bsize{\ri{\gamma}_\tau t_\tau y_\tau}\le\bb^{6\sn+2}\le\bxs{S}}; so \m{y_\tau\in\mathcal C_\tau} by the choice
of~$q$. On \m{\prod_\tau\mathcal C_\tau} we write the cut kernel as the full kernel minus the difference. The integral
of the full kernel is bounded by \eqref{e:B2}, in norm by \m{\tfrac12\nuS\,\csize{S}{\sro+1}{x}}. The integral of the
norm of the difference is at most
\[\int\varphi(\gamma,y)\,\bigl\|\besker{\full}{S}{\gamma}{y}-\besker{\cut}{S}{\gamma}{y}\bigr\|\,dy\ \le\
\tfrac12\,\nuS\,\bb^{-\sk^3\fE}\ \le\ \tfrac12\,\nuS\,\csize{S}{\sro+1}{x},\]
by \inp{B5} for the side \m{\full} and by the first inequality of Lemma~\ref{l:cell}. Adding the two bounds gives
the first case of~\eqref{e:B1}.
\end{proof}
\leanline{\lean{enormamplitudeSumleclassSizetop}}

\begin{proof}[Proof of Theorem B\,(b)]
If $s$ does not agree, the block vanishes by~\inp{B6}. Let $s$ agree, and let $\Xi'$ be the set of the
\m{\gamma\in\GS{S}{E}} with \m{\Nn^{\fE}\,\mnorm{S}{j}{\gamma}<\tdnorm{S}{\gamma}}. By \inp{S10} the sets
\m{\Xi\cap\Xi'} and \m{\Xi\setminus\Xi'} are stable, and by Lemma~\ref{l:split} the block is the sum of the blocks of
these two sets.

\emph{The part in $\Xi'$.} If \m{\gamma\in\Xi\cap\Xi'} then \m{\tdnorm{S}{\gamma}>0}, so \m{\tdisc{S}{\gamma}\ne0} and, since
$\gamma$ has level~$j$, \m{j=\sro+1}. So either \m{\Xi\cap\Xi'} is empty or \m{j=\sro+1} and part~(a) applies to
\m{\Xi\cap\Xi'}. At the level \m{\sro+1} the lattices \m{\mathcal R_j} and \m{\mathcal R} coincide, the indices
\m{\xri{S}{\sro+1}{\cdot}} are~$1$, and for \m{\gamma\in\Xi'} the class gauge is
\m{\cgauge{S}{\sro+1}{c}{c'}{\gamma}=\tdnorm{S}{\gamma}}. So the bound of part~(a) is the part of the sum in (b) over
\m{\gamma\in\Xi'}.

\emph{The part outside $\Xi'$.} By \inp{S2} the norm of the block of \m{\Xi\setminus\Xi'} is at most the sum of
the absolute amplitudes \m{\ampabs{\cut}{S}{c}{c'}{\gamma}} over
\m{\gamma\in(\Xi\setminus\Xi')\cap\terms{S}{c}{s}{c'}}. Let $\gamma$ be such an element with non-zero absolute amplitude.
By Lemma~\ref{l:cs} the Schur integral \m{\schi{S}{c}{c'}{\gamma}} is not zero, so there is $y$ with
\m{\bcut{S}{c}(y)\ne0} and \m{\bcut{S}{c'}(\ri{\gamma}^{-1}y)\ne0}; both points are boxed, and $\gamma$ is near. Let
$\delta$ be the fixed representative of the double coset \m{\clUo{S}{c}\,\gamma\,\clUo{S}{c'}}. It is near by~\inp{S3},
so \m{\IsNRep{S}{c}{c'}{\delta}}; it lies in \m{\Xi\setminus\Xi'} by stability and in \m{\terms{S}{c}{s}{c'}} by
Lemma~\ref{l:class}; and it has level~$j$. Hence the sum of the absolute amplitudes is at most the sum, over the
near representatives \m{\delta\in(\Xi\setminus\Xi')\cap\terms{S}{c}{s}{c'}}, of the sums over their double cosets,
and each of these is at most \m{\nuS\,\csize{S}{j}{\cgauge{S}{j}{c}{c'}{\delta}}} by Theorem~A\,(b), whose hypothesis
is \m{\delta\notin\Xi'}. This is the part of the sum in (b) over \m{\gamma\notin\Xi'}.
\end{proof}
\leanline{\lean{enormamplitudeSumleclassSize}}

\sect{Agreeing characters, and the proofs of Theorems D and E}{s:crude}
\begin{proposition}[terms with disagreeing characters cancel]\label{p:incoh}
In Setting~\ref{set:data}, let $\sigma$ be a side and \m{\Xi\subseteq\GS{S}{E}} a stable set, \m{\IsStab{S}{\Xi}},
such that no \m{\gamma\in\Xi} satisfies \m{\ChA{S}{\gamma}}. Then \m{\ampsum{\sigma}{S}{\Xi}{c}{s}{c'}=0} for all
\m{c,c'\in\Cl{S}} and \m{s\in\GRes{S}}.
\end{proposition}
\begin{proof}
The group \m{\clU{S}{c}} acts by left multiplication on the set of the cosets \m{x=\gamma\,\clU{S}{c'}}. For a coset
$x$ let \m{\Gamma_x\subseteq\clU{S}{c}} be its stabilizer and \m{\theta_x} a function as in \inp{S9} for~\m{\Gamma_x}.
In every orbit fix one coset. Let
\[\varphi(\gamma,y)=\cutoff{S}{c'}\bigl(\ri{\gamma}^{-1}\bigr)\cdot\theta_{\gamma\clU{S}{c'}}(y)\]
if the coset of $\gamma$ is the one fixed in its orbit, and \m{\varphi(\gamma,y)=0} otherwise; the first factor is
the value of the cutoff at the point \m{\ri{\gamma}^{-1}} of~\m{\Ginf}.

\emph{The weight is normalized.} Let \m{\gamma\in\GS{S}{E}}, \m{x=\gamma\clU{S}{c'}} and \m{y\in\Ginf}. The coset of
\m{u\gamma u'^{-1}} is \m{ux}, and \m{\sum_{u'\in\clU{S}{c'}}\cutoff{S}{c'}(\ri{u'}\,\ri{\gamma}^{-1}\ri{u}^{-1})=1}
by~\inp{B1}. So \m{\sum_{(u,u')}\varphi(u\gamma u'^{-1},\ri{u}y)} is the sum of \m{\theta_{ux}(\ri{u}y)} over the
\m{u\in\clU{S}{c}} for which \m{ux} is the fixed coset~$x_0$ of the orbit. If \m{u_1x=x_0}, these $u$ are the
\m{v\,u_1} with \m{v\in\Gamma_{x_0}}, and \m{\sum_{v\in\Gamma_{x_0}}\theta_{x_0}(\ri{v}\,\ri{u_1}y)=1}.

\emph{Every integral vanishes.} By \inp{B5} the block is the sum of the integrals
\m{\int\varphi(\gamma,y)\,\besker{\sigma}{S}{\gamma}{y}\,dy} over \m{\gamma\in\Xi\cap\terms{S}{c}{s}{c'}}. Let $\gamma$ be
such an element whose coset $x$ is the fixed one; the integral is a multiple of
\m{\int\theta_x(y)\,F_0(y)\,dy} with \m{F_0(y)=\besker{\sigma}{S}{\gamma}{y}}. The function $F_0$ is continuous. For
\m{u\in\Gamma_x} we have \m{u'=\gamma^{-1}u\gamma\in\clU{S}{c'}}, and \m{F_0(\ri{u}y)} is the tensor product of
the Bessel functions at \m{\ri{u}y} and the conjugates of those at \m{\ri{u'}\,\ri{\gamma}^{-1}y}; by \inp{B7} it
equals \m{F_0(y)}. So $F_0$ is invariant under the image of \m{\Gamma_x}, which is discrete, and by \inp{S9} the
integral is the integral of $F_0$ over the quotient. Since $\gamma$ does not satisfy \m{\ChA{S}{\gamma}}, there are
\m{u_2,u_2'\in\UE} with \m{u_2\gamma=\gamma u_2'} and
\m{c_0=\uchar{S}{\rfi{u_2}}^{-1}\uchar{S}{\rfi{u_2'}}\ne1}; by \inp{B7}, \m{F_0(\ri{u_2}y)=c_0\,F_0(y)} for all~$y$,
and \m{|c_0|=1}. By \inp{L3} the integral over the quotient is~$0$.
\end{proof}
\leanline{\lean{amplitudeSumeqzeroofnotcharsAgree}}

\begin{lemma}[shared vectors climb the flag]\label{l:flag}
In Setting~\ref{set:S}, let \m{\gamma\in\GS{S}{E}} with \m{\ChA{S}{\gamma}}, let \m{a\le\sro} and \m{c_1,c_2\in E}, and
suppose \m{\HCFV{\gamma}{c_1}{c_2}{a}}. Then \m{c_1\cj{c_1}=c_2\cj{c_2}}, and \m{\HCFV{\gamma}{c_1}{c_2}{a'}} for every
$a'$ with \m{a\le a'\le\sro}.
\end{lemma}
\begin{proof}
By descending induction on~$a$. Let \m{v,v'\in E^{\sn}} with \m{\mathfrak fv=\gamma\mathfrak fv'}, \m{v_a=c_1},
\m{v'_a=c_2} and zero entries above~$a$. For \m{a=\sro} the first statement is \inp{S13} and the second is empty.

Let \m{a<\sro} and suppose the lemma for \m{a+1}. Let \m{w,w'\in E^{\sn}} be the shifted vectors: \m{w_0=w'_0=0},
\m{w_{a''+1}=v_{a''}}, \m{w'_{a''+1}=v'_{a''}}. We show below that
\begin{equation}\label{e:flag}
\lambda(\mathfrak fw)=\lambda(\gamma\mathfrak fw')\quad\text{for every linear form }\lambda
\text{ with }\lambda(\fr{a''})=\lambda(\gamma\fr{a''})=0\text{ for }a''\le a,
\end{equation}
the forms being $E$-linear on \m{E^{\Ix{\fs}}}. By \inp{L4} there are then \m{v^+,v'^+} with \m{\mathfrak fv^+=\gamma\mathfrak fv'^+} that agree with $w$, $w'$ at the
indices above~$a$. So \m{\HCFV{\gamma}{c_1}{c_2}{a+1}}, and the induction hypothesis gives both statements.

To prove \eqref{e:flag}, let $\lambda$ be such a form. By \inp{S14} there is a vector $t$ with
\m{\lambda(\gamma t')=h(t,t')} for all~$t'$, with \m{h=h_{\TE}}; then \m{\lambda=h(\gamma t,\cdot)}. Let \m{x\in E} be
arbitrary and \m{r_0=-\gq\,\vpS{S}\,x}. We apply \inp{S14} twice, with the index \m{a+1}: to the coefficient vector
$v$ and the vector \m{\gamma t}, and to $v'$ and~$t$. The hypotheses hold, because \m{h(\gamma t,\fr{a''})=\lambda(\fr{a''})=0}
and \m{h(t,\fr{a''})=\lambda(\gamma\fr{a''})=0} for \m{a''\le a}. We get \m{u,u'\in\UE} with
\[\sd{u}=r_0\sum_{a''<\sro}\lambda(\fr{a''+1})\,v_{a''}=r_0\,\lambda(\mathfrak fw),\qquad
\sd{u'}=r_0\sum_{a''<\sro}\lambda(\gamma\fr{a''+1})\,v'_{a''}=r_0\,\lambda(\gamma\mathfrak fw').\]
They are the Eichler maps of the pairs \m{(r_0\mathfrak fv,\gamma t)} and \m{(r_0\mathfrak fv',t)}; as $\gamma$ is an
isometry and \m{\gamma(r_0\mathfrak fv')=r_0\mathfrak fv}, we have \m{\gamma u'\gamma^{-1}=u}. Since the characters agree at~$\gamma$,
\m{\uchar{S}{\rfi{u}}=\uchar{S}{\rfi{u'}}}, that is
\[\psi_{E,f}\bigl(x\cdot\lambda(\mathfrak fw)\bigr)=\psi_{E,f}\bigl(x\cdot\lambda(\gamma\mathfrak fw')\bigr).\]
This holds for all \m{x\in E}, so \eqref{e:flag} follows from~\inp{L5}.
\end{proof}
\leanline{\lean{hasCommonFlagVectorofle}}

\begin{proposition}[agreeing characters force a level]\label{p:level}
In Setting~\ref{set:S}, let \m{\gamma\in\GS{S}{E}} with \m{\ChA{S}{\gamma}} and \m{\tdisc{S}{\gamma}=0}. Then there is
\m{j\in\{0,\dots,\sro+1\}} with \m{\HasLev{S}{j}{\gamma}}.
\end{proposition}
\begin{proof}
The condition \m{\IsNF{\sro+1}{\gamma}} is empty, so there is a least $j$ with \m{\IsNF{j}{\gamma}}. Put
\m{Z=\Zf{S}{\gamma}}.

\emph{Claim.} If \m{\mathfrak fv=\gamma\mathfrak fv'} and \m{v'_a=0} for all \m{a\ge j}, then \m{v=v'=0}. We show
\m{v_a=v'_a=0} by descending induction on~$a$. Suppose it for the indices above~$a$. If \m{v'_a\ne0}, then
\m{\HCFV{\gamma}{v_a/v'_a}{1}{a}}, and Lemma~\ref{l:flag} gives \m{(v_a/v'_a)\cj{(v_a/v'_a)}=1} and
\m{\HCFV{\gamma}{v_a/v'_a}{1}{a'}} for all \m{a'\ge a}; so \m{\IsNF{a}{\gamma}}, and \m{j\le a} by minimality, against
\m{v'_a=0} for \m{a\ge j}. Hence \m{v'_a=0}. If \m{v_a\ne0}, then \m{\HCFV{\gamma}{1}{0}{a}}, and Lemma~\ref{l:flag}
gives \m{1=0}. Hence \m{v_a=0}.

\emph{The minor \m{\lmin{j}{Z}} is not zero.} Otherwise \m{\padl{j}{Z}} has a kernel vector \m{v'\ne0}: so
\m{v'_a=0} for \m{a\ge j} and \m{(Zv')_a=0} for \m{a<j}. In fact \m{Zv'=0}: for \m{j=\sro+1} this is said; for
\m{j\le\sro}, let \m{v_0\in E^{\sn}}; by \inp{S12} for \m{\gamma^{-1}} there are \m{v_1,v_2} with
\m{\gamma\mathfrak fv_1=\mathfrak fv_2} and \m{(v_2)_a=(v_0)_a} for \m{a\ge j}; then
\m{h(\mathfrak fv_2,\gamma\mathfrak fv')=h(\mathfrak fv_1,\mathfrak fv')=-\Nn\,\cj{(v_1)_{\sro}}\,v'_{\sro}=0}, and
\m{h(\mathfrak f(v_2-v_0),\gamma\mathfrak fv')=0} because \m{v_2-v_0} is supported on the indices below~$j$; so
\m{\sum_a\cj{(v_0)_a}\,(Zv')_a=h(\mathfrak fv_0,\gamma\mathfrak fv')=0} for all~$v_0$. Now \inp{S13} with \m{D_0=0} applies
to \m{t=\gamma\mathfrak fv'}: the number \m{-(\Nn|(v'_{\sro})_\tau|)^2} is \m{\le0}, so \m{t=\mathfrak fv} for
some~$v$, and the claim gives \m{v'=0}, a contradiction.

\emph{The level is at most \m{\sro}.} Suppose \m{j=\sro+1}; then \m{\det Z=\lmin{j}{Z}\ne0}. Let $v'$ be the last
column of the adjugate of~$Z$; then \m{Zv'} has the entry \m{\det Z} at \m{\sro} and $0$ elsewhere, and
\m{v'_{\sro}=\lmin{\sro}{Z}}. By \inp{S13} the vector \m{t=\gamma\mathfrak fv'} satisfies, at every~$\tau$,
\m{\Nn\sum_i|(t_i)_\tau|^2=|\det Z_\tau|^2-\Nn^2|\lmin{\sro}{Z}_\tau|^2=\tdisc{S}{\gamma}_\tau=0}; so \m{t=\mathfrak fv} for
some~$v$, and the claim, whose hypothesis on $v'$ is empty for \m{j=\sro+1}, gives \m{v'=0}, hence \m{\det Z=0}, a
contradiction. So \m{j\le\sro}; as \m{\tdisc{S}{\gamma}=0}, both sides of the last condition in the definition of the
level are false, and \m{\HasLev{S}{j}{\gamma}}.
\end{proof}
\leanline{\lean{existshasLevel}}

\begin{proof}[Proof of Theorem D\,(a)]
Write \m{A(c,s,c')=\ampsum{\cut}{S}{\Xi}{c}{s}{c'}}, \m{j_0=\min(\sro,j)} and \m{j_1=j-j_0}. Since \m{\nuS>0}, we have to
show
\begin{equation}\label{e:D0}
\sum_{c,c'}\sum_{s\in\Omega}\|A(c,s,c')\|^2\ \le\ w\,\nuS^2+\frac{C\,\nuS^2}{\bb^{30}\,4^{\sro+1-j}}.
\end{equation}

\emph{Step 1: splitting by links.} For a subspace $\Lambda$ of \m{(E^{\Mx{\fs}})^2} let
\m{\Xi_\Lambda=\Xi\cap\{\gamma:\link{\gamma}=\Lambda\}}; it is stable by~\inp{S10}, and its elements have level~$j$.
Let \m{A_\Lambda(c,s,c')} be its block. By the absolute summability \inp{B4},
\m{A(c,s,c')=\sum_\Lambda A_\Lambda(c,s,c')}, so
\[\|A(c,s,c')\|^2\ \le\ \Bigl(\sum_\Lambda\|A_\Lambda(c,s,c')\|\Bigr)^{2}=\sum_{(\Lambda,\Lambda')}\|A_\Lambda(c,s,c')\|\,\|A_{\Lambda'}(c,s,c')\|.\]
The pairs \m{(\Lambda,\Lambda)} with \m{\Lambda\in R} contribute to \eqref{e:D0} exactly
\m{\linkE{\cut}{S}{\Xi}{\Omega}{R}\le w\,\nuS^2}. Call the other pairs \emph{off-diagonal}; it remains to bound
their contribution~$\mathfrak S$.

\emph{Step 2: block bounds.} By Theorem~B\,(b) for \m{\Xi_\Lambda}, the block \m{A_\Lambda(c,s,c')} vanishes unless
$s$ agrees, and then its norm is at most the sum of \m{\nuS\,\csize{S}{j}{\cgauge{S}{j}{c}{c'}{\gamma}}} over the
near representatives \m{\gamma\in\Xi_\Lambda\cap\terms{S}{c}{s}{c'}}. Hence \m{\mathfrak S} is at most \m{\nuS^2} times the
sum of the products
\m{\csize{S}{j}{\cgauge{S}{j}{c}{c'}{\gamma}}\cdot\csize{S}{j}{\cgauge{S}{j}{c}{c'}{\gamma'}}} over the tuples
\m{(c,c',s,\gamma,\gamma')} in which \m{s\in\Omega} agrees, \m{\gamma,\gamma'\in\Xi\cap\terms{S}{c}{s}{c'}} are near
representatives for \m{(c,c')}, and \m{(\link{\gamma},\link{\gamma'})} is off-diagonal.

\emph{Step 3: one pair.} Fix such a tuple, and write \m{G=\cgauge{S}{j}{c}{c'}{\gamma}},
\m{G'=\cgauge{S}{j}{c}{c'}{\gamma'}}. We claim
\begin{equation}\label{e:D1}
\csize{S}{j}{G}\cdot\csize{S}{j}{G'}\ \le\ \frac{C}{\tl^{3}}\cdot\bigl(\pairw{S}{j}{c}{c'}{\gamma}{\gamma'}+\pairw{S}{j}{c}{c'}{\gamma'}{\gamma}\bigr).
\end{equation}
Put \m{U=\clink{\gamma}{\gamma'}=\clink{\gamma'}{\gamma}}, \m{\mathsf k=\pos{S}{c}{c'}{\gamma}},
\m{\mathsf k'=\pos{S}{c}{c'}{\gamma'}}, \m{d_1=\latd{S}{c}(U)}, \m{o=\fvi{S}{c}{\mathsf k}{1}{U}},
\m{o'=\fvi{S}{c}{\mathsf k'}{1}{U}}, \m{\mathfrak d=\sati{S}{\mathsf k}} and \m{\mathfrak d'=\sati{S}{\mathsf k'}}. By \inp{P1}
and \inp{P2}, applied to \m{(\gamma,\gamma')} and to \m{(\gamma',\gamma)}, the numbers \m{d_1}, $G$, $G'$ are positive,
\m{o,o',\mathfrak d,\mathfrak d'\ge1}, and there is \m{i\in\NN} with \m{\sk\le\dim_EU+i+j_1}, with
\m{\link{\gamma}=\link{\gamma'}} if \m{i=0}, and with
\[d_1^{\,2}\cdot o\,\mathfrak d\cdot o'\mathfrak d'\cdot X'^{\,2}\ \le\ (2/\gq)^{2i\fE}\cdot G\,G',\]
where \m{X'=X} if \m{\sro<j} and \m{i=0}, and \m{X'=1} otherwise; \m{X'\ge1}. By Definition~\ref{d:cs},
\begin{align*}
\csize{S}{j}{G}\cdot\csize{S}{j}{G'}&=\frac{\tl^4\,\mathsf D_j}{\cw{S}{0}\,\gq^{2j_0\sk\fE}}\,(GG')^{-(\sk+1)/2}\\
&\le\ \frac{\tl^4\,\mathsf D_j}{\cw{S}{0}\,\gq^{2j_0\sk\fE}}\cdot\frac{(2/\gq)^{i\fE(\sk+1)}}{X'^{\,\sk+1}\,d_1^{\,\sk+1}}\cdot
\bigl(o\,\mathfrak d'\cdot o'\mathfrak d\bigr)^{-(\sk+1)/2}.
\end{align*}
As \m{o\,\mathfrak d'\ge1} and \m{o'\mathfrak d\ge1}, the last factor is at most
\m{(o\,\mathfrak d')^{-\sk/2}(o'\mathfrak d)^{-\sk/2}\le\tfrac12\bigl((o\,\mathfrak d')^{-\sk}+(o'\mathfrak d)^{-\sk}\bigr)}. On the other
side, with \m{i_0=\sk-\dim_EU},
\[\pairw{S}{j}{c}{c'}{\gamma}{\gamma'}+\pairw{S}{j}{c}{c'}{\gamma'}{\gamma}=
\frac{d_1^{-(\sk+1)}}{\cw{S}{i_0}\cdot\gq^{2j_0\sk\fE}}\cdot\bigl((o\,\mathfrak d')^{-\sk}+(o'\mathfrak d)^{-\sk}\bigr),\]
by the definition of the pair weights, with \m{\gix{\gid{S}}^{\#\IdxB{j}}=\gq^{2j_0\fE}}. Since
\m{\cw{S}{i_0}/\cw{S}{0}=(2\tl^2)^{i_0}\le(2\tl^2)^{i+j_1}} and
\m{(2/\gq)^{i\fE(\sk+1)}(2\tl^2)^{i}\le\codimF{S}^{\,i}}, \eqref{e:D1} follows from
\begin{equation}\label{e:D2}
\tl^4\cdot\mathsf D_j\cdot(2\tl^2)^{j_1}\cdot\codimF{S}^{\,i}\cdot X'^{-\sk}\ \le\ C\,\tl^{-3}.
\end{equation}
We prove \eqref{e:D2} by cases, using \m{\tl\ge2} and~\inp{A3}.
\emph{Case \m{j=\sro+1}, \m{i=0}.} Then \m{\link{\gamma}=\link{\gamma'}}; as the pair is off-diagonal, this link is
not in~$R$, so \m{\tl^{12}X^{-\sk}\dm\le C} by the hypothesis of the theorem; the left side of \eqref{e:D2} is
\m{2\,\tl^6\,\dm\,X^{-\sk}\le\tl^{-3}\cdot\tl^{12}\dm X^{-\sk}}.
\emph{Case \m{j=\sro+1}, \m{i\ge1}.} Then \m{X'=1} and \m{\codimF{S}^{\,i}\le\codimF{S}}; the left side is at most
\m{2\,\tl^6\,\codimF{S}\,\dm\le2\,\tl^{-12}\dd\le\tl^{-3}C}, by \m{\tl^{18}\codimF{S}\dm\le\dd\le C}.
\emph{Case \m{j\le\sro}.} Then \m{j_1=0}, \m{X'=1}, \m{\codimF{S}^{\,i}\le1}, and the left side is at most
\m{\tl^{17-32j}\,\dd}. For \m{j=0} the element $\gamma$ is in normal form of level~$0$ and lies in
\m{\Xi\cap\terms{S}{c}{s}{c'}} with \m{s\in\Omega}, so \m{\tl^{20}\dd\le C} by Setting~\ref{set:en}. For \m{j\ge1},
\m{\tl^{17-32j}\dd\le\tl^{-3}\dd\le\tl^{-3}C}.

\emph{Step 4: summation.} In a tuple of Step~2 the element $s$ is determined by \m{(c,c',\gamma)}. So by
\eqref{e:D1} and Proposition~\ref{p:pairs}, applied twice,
\m{\mathfrak S\le\nuS^2\cdot C\,\tl^{-3}\cdot4\,\tl^2=4\,C\,\nuS^2/\tl}. Finally
\m{4\cdot\bb^{30}\cdot4^{\sro+1-j}\le\bb^{\sro+32-j}\le\tl} by \m{\bb\ge4} and~\inp{S1}. This gives~\eqref{e:D0}.
\end{proof}
\leanline{\lean{energyleoflevelpairs}}

\begin{proof}[Proof of Theorem D\,(b)]
We prove, by induction on~$l$, the statement for all stable subsets \m{\Xi_1\subseteq\Xi} in place of~$\Xi$; the
hypotheses of Setting~\ref{set:en} hold for them with the same $C$, $\Omega$ and~$X$. For \m{l=0} no element of
\m{\Xi_1} satisfies \m{\ChA{S}{\cdot}}, all blocks vanish by Proposition~\ref{p:incoh}, and the energy is~$0$. Let the
statement hold for some \m{l\le\sro}, and let every \m{\gamma\in\Xi_1} with \m{\ChA{S}{\gamma}} have a level below
\m{l+1}. Let \m{\Xi_2} be the set of the elements of \m{\Xi_1} of level~$l$ and \m{\Xi_3=\Xi_1\setminus\Xi_2}; both are
stable by~\inp{S10}. By Lemma~\ref{l:split} and Lemma~\ref{l:add} with \m{\varepsilon=1}, the energy of the blocks of
\m{\Xi_1} is at most twice that of \m{\Xi_3} plus twice that of~\m{\Xi_2}. An element of \m{\Xi_3} with agreeing characters has a
level below \m{l+1} that is not~$l$, so the induction hypothesis applies to~\m{\Xi_3}. To \m{\Xi_2} we apply part~(a)
with \m{j=l\le\sro}, \m{w=0} and $R$ empty; the link energy is then~$0$. Both energies are at most
\m{C\,\bb^{-30}4^{-(\sro+1-l)}}, and four times this bound is \m{C\,\bb^{-30}4^{-(\sro+1-(l+1))}}.
\end{proof}
\leanline{\lean{energyleoflevellt}}

\begin{proof}[Proof of Theorem E]
Let \m{\Xi_2=\Xi\cap\{\gamma:\tdisc{S}{\gamma}\ne0\}} and \m{\Xi_3=\Xi\setminus\Xi_2}; both are stable by~\inp{S10}. The
elements of \m{\Xi_2} have level \m{\sro+1} \inp{S12}; so Theorem~D\,(a) with \m{j=\sro+1} gives the energy bound
\m{w+C\,\bb^{-30}} for the cut blocks of~\m{\Xi_2}. An element of \m{\Xi_3} with agreeing characters has a level
by Proposition~\ref{p:level}, and this level is not \m{\sro+1} because \m{\tdisc{S}{\gamma}=0}; so Theorem~D\,(b) with
\m{l=\sro+1} gives the bound \m{C\,\bb^{-30}} for the cut blocks of~\m{\Xi_3}. By Lemma~\ref{l:split} and
Lemma~\ref{l:add} with \m{\varepsilon=\bb^{4}}, the cut blocks of $\Xi$ have energy at most
\[B_1=(1+\bb^{4})\,C\,\bb^{-30}+(1+\bb^{-4})(w+C\,\bb^{-30}).\] By Theorem~C, \m{\dd\le C}, and Lemma~\ref{l:add} once more
with \m{\varepsilon=\bb^{4}}, applied to the difference of the full and the cut blocks and to the cut blocks,
\[\energy{S}{\ampsum{\full}{S}{\Xi}{\cdot}{\cdot}{\cdot}}{\Omega}\ \le\ (1+\bb^{4})\,\bb^{-30}\,C+(1+\bb^{-4})\,B_1.\]
In this bound the coefficient of $w$ is \m{(1+\bb^{-4})^2}, and the coefficient of $C$ is
\m{\bb^{-30}\,(2\,\bb^{4}+4+3\,\bb^{-4}+\bb^{-8})\le\bb^{-12}}.
\end{proof}
\leanline{\lean{energyleoflinkEnergyle}}

\sect{Notes on the formal text}{s:notes}
\subhead{Arguments that the text reading gets wrong.} The program that reads the arguments of the declarations
from the text lists, for several members of these files, arguments that Lean does not take; the statements of
this note follow the Lean text. \lean{energy} and \lean{linkEnergy} take the data $S$ implicitly, and the second
takes the side as an instance; the notations \m{\energy{S}{\mathcal A}{\Omega}} and
\m{\linkE{\sigma}{S}{\Xi}{\Omega}{R}} show both, as in the other notes of the unit. In the first file a second
declaration of the field and of the data hides the earlier instances, so that \lean{classSize}, \lean{boxCutoff},
\lean{schurWeight}, \lean{schurIntegral} and Lemma~\ref{l:incut} assume neither that $E$ is stably unramified nor
that $S$ is admissible, and Lemma~\ref{l:cell} assumes only admissibility; the same happens in the third file,
where the lists for Proposition~\ref{p:pairs} and Theorem~B name a class and an element that are bound variables.
Lemma~\ref{l:class}, Lemma~\ref{l:flag} and Proposition~\ref{p:level} need no standing hypothesis at all.

\subhead{The side.} The formal text selects the Bessel function by an instance. The default is the full
function; from the definition of the Schur integral on, and in the second, third and fourth files, a local
instance makes it the cut function. So Theorems A, B and D are statements about cut functions, Theorem~C compares
the default with an explicit cut, and Theorem~E names the full function explicitly while its hypothesis on the
link energy is for the cut. Proposition~\ref{p:incoh} holds for both sides.

\subhead{Conventions of the formal text.} Sums and integrals of non-negative terms are taken in \m{[0,\infty]}.
An inverse power of a real number that is $0$ is~$0$; this enters the weights
\m{\ipw{S}{c}{\eta}{p}}, the term weights and the sums of Propositions \ref{p:latdisc} and~\ref{p:count}. The class
size is defined for every real $x$ by the conventions for square roots and quotients; it is used only at positive
arguments, the class gauge being positive by~\inp{P1}. The amplitudes are Bochner integrals, and the amplitude
sums are unconditional sums; both are $0$ by convention when they do not converge, which under the hypotheses of
Setting~\ref{set:data} does not happen.

\subhead{Statements.} Theorem~B\,(b) is stated in Lean with one sum over the elements $\gamma$ subject to the
condition that $s$ agrees; this is the same as the two cases above. In Theorems D and E the exponent of $4$ is
written with the reversal of the index~$j$. The hypothesis of Proposition~\ref{p:avg} asks nothing of the set
$Q_0$ but \m{\mu(Q_0)\ge1}. In Lemma~\ref{l:short} the formal text writes the vector as \m{\det Y} times a row of
the transposed inverse; it is the column of the adjugate. Lemma~\ref{l:gindex} is stated for the idele of~\m{\gq}
and uses only that it is an integral unit of the adeles. Theorem~E is an assembly: its content lies in Theorems
C and~D and in Proposition~\ref{p:level}. Theorem~D\,(b) assembles Theorem~D\,(a) and
Proposition~\ref{p:incoh}.

\subhead{Exponents.} The formal proofs end, in nine places, with comparisons of exponents of \m{\bb} that are left to
decision procedures, in one case (Lemma~\ref{l:cell}) after taking logarithms and forming an explicit linear
combination of the hypotheses. In this note all of them are carried out by hand, with cruder intermediate bounds
than the formal text needs: Theorem~C; both inequalities of Lemma~\ref{l:cell}; Step~7 of Theorem~A\,(a) and the
end of Theorem~A\,(b); Propositions \ref{p:latdisc}, \ref{p:count} and~\ref{p:pairs}; Lemma~\ref{l:short}; and
\eqref{e:D2}. In Lemma~\ref{l:cell} the bookkeeping of the powers of \m{\prod_\tau|\kapS{S}{\tau}|} is ours and exact,
where the formal text first bounds some of these factors by a power of~\m{\bb}.

\subhead{What is taken for granted.} Everything in Section~\ref{s:inputs}. In particular the three analytic
inputs from the real place (the volume kernel \inp{R2}, the open-cell bound \inp{R3}, the congruence count
\inp{R5}), the inequality of pair invariants \inp{P2}, the frame count \inp{B3}, the change of weight \inp{B5}
and the cut error \inp{B4} carry much of the weight of Theorems A, B and D; they are proved in the notes named
there.

\subhead{How this note was written.} From the dossier of the unit: the texts of the members of the four files and
their structured proofs. I read the four Lean files completely, the proofs line by line, in particular those
that the readers had marked as hiding arguments in automation: the logarithmic linear combination of
Lemma~\ref{l:cell}, the averaging over classes in Proposition~\ref{p:latdisc}, the inclusion of lattices in Step~3
of Proposition~\ref{p:pairs}, the numerical ends of Theorem~A and of \eqref{e:D2}, and the open-cell step of
Theorem~B\,(a). Every statement of this note was compared with its Lean statement, argument by argument, and with
the list of arguments that the program prints. Of the Inputs I read the Lean statements, in the files of the
set-up, of the box partitions and amplitude sums, of the pairing of links, of the real place and of the
background library, and for the definitions of the weights and indices their Lean text; I did not read their
proofs. The proofs given here follow the formal ones in structure; the arithmetic of exponents is ours, the
proof of the vanishing of the columns in Step~5 of Proposition~\ref{p:count} and the two cases of
\eqref{e:cnt2} are written out where the formal text calls a simplifier, and Step~3 of
Proposition~\ref{p:pairs} is arranged by vectors instead of matrices. The references for the background library
are those of the catalogue of the write-up and of the second reader; I did not open the cited works. The
program check of the unit's dossier covered \UnitSteps{} leaf steps of \UnitProofs{} proofs.

\head{Appendix A. Concordance with the Lean artifact}
The formal names of the objects: \texttt{S.b}, \texttt{S.N}, \texttt{S.g} are \m{\bb}, \m{\Nn}, \m{\gq};
\texttt{numInfPlaces E} is \m{\fE}; \texttt{𝔰.ϱ}, \texttt{𝔰.k}, \texttt{𝔰.n} are \m{\sro}, \m{\sk}, \m{\sn};
\texttt{tailScale}, \texttt{unipCovolume}, \texttt{besselEnergy}, \texttt{mainEnergy}, \texttt{countWeight},
\texttt{codimFactor} are \m{\tl}, \m{\nuS}, \m{\dd}, \m{\dm}, \m{\cw{S}{i}}, \m{\codimF{S}}; \texttt{termsOver},
\texttt{termPosition}, \texttt{classUnip}, \texttt{classUnip₀}, \texttt{classCutoff} are \m{\terms{S}{c}{s}{c'}},
\m{\pos{S}{c}{c'}{\gamma}}, \m{\clU{S}{c}}, \m{\clUo{S}{c}}, \m{\cutoff{S}{c}}; \texttt{amplitudeAbs},
\texttt{amplitudeSum} are \m{\ampabs{\sigma}{S}{c}{c'}{\gamma}}, \m{\ampsum{\sigma}{S}{\Xi}{c}{s}{c'}};
\texttt{termDiscNorm}, \texttt{minorNorm}, \texttt{classGauge}, \texttt{xRowsIndex} are \m{\tdnorm{S}{\gamma}},
\m{\mnorm{S}{j}{\gamma}}, \m{\cgauge{S}{j}{c}{c'}{\gamma}}, \m{\xri{S}{j}{\mathsf k}}; \texttt{HasLevel},
\texttt{IsNearRep}, \texttt{IsStable}, \texttt{Agrees}, \texttt{CharsAgree}, \texttt{FixesNormAtLeast} are the
relations \m{\HasLev{S}{j}{\gamma}}, \m{\IsNRep{S}{c}{c'}{\gamma}}, \m{\IsStab{S}{\Xi}}, \m{\Agr{S}{s}},
\m{\ChA{S}{\gamma}}, \m{\FixN{s}{X}}. The Lean letters \texttt{T}, \texttt{U} for sets of rational elements are
$\Xi$, $\Xi'$; \texttt{v} for an infinite place is~$\tau$; \texttt{w} and \texttt{C} are $w$ and~$C$; the
natural numbers \texttt{s} and \texttt{n} of the counting statements are $m$ and~$l$; the subspace \texttt{T} is~$U$.
Three statements as they stand in the artifact:
\quoted{energy}
\quoted{tsumschurIntegralle}
\quoted{energyleoflinkEnergyle}

{\fontsize{8.6}{11.6}\selectfont\renewcommand{\leanfont}{\ttfamily\fontsize{7.7}{10}\selectfont\color{leancol}}
\setlength{\LTleft}{0pt}\setlength{\LTright}{\fill}\setlength{\tabcolsep}{0pt}
\begin{longtable}{@{}p{27mm}p{67mm}p{35mm}r@{}}
\textbf{In this note} & \textbf{Lean declaration} & \textbf{File} & \textbf{Lines}\\\midrule\endhead
\ConcordanceRows
\end{longtable}}

\head{Notation}
\begingroup\fontsize{8.6}{10.9}\selectfont\raggedright
\setlength{\columnsep}{9mm}\setlength{\parskip}{0pt}
\begin{multicols}{2}
\nt{$E$, $\tau$, $x_\tau$, $\fE$}{the CM field with its CM type, an infinite place, the image of \m{x\in E} at $\tau$, the number of infinite places}
\nt{$\AEf$, $\hO$, $\psi_{E,f}$}{finite adeles, integral adeles, the standard additive character}
\nt{$\Nm$, $\rdE{E}$, $\dE{E}$, $\volO{E}$}{norm, root discriminant, discriminant, volume of the integral adeles}
\nt{$\gix{a}$, $[M':M]$}{generalized index of an adele; index of \m{M\cap M'} in $M'$ ($0$ if infinite)}
\nt{$\hO\langle M\rangle$, $\fadm{M}$, $M^*$}{span of the columns; a matrix over $E$ read over the adeles; conjugate transpose}
\nt{$\fs=(\varrho,\kappa,k_0)$}{the shape; \m{\sro}, \m{\sk}, \m{\sn}, \m{\rr{\fs}} its numbers; \m{\Ix{\fs}}, \m{\Mx{\fs}} its index sets}
\nt{$e_a$, $e^*_a$, $z$, $\fr{a}$, $\mathfrak f$}{basis vectors; the frame vectors; the matrix of the frame}
\nt{$T$, $\JT{T}$, $\GT{T}$, $h_T$, $h$}{an $\fs$-form, its split Gram matrix, its unitary group, the hermitian form}
\nt{$\UT{T}$, $\UoT{T}$, $\HT{T}$, $\ST{T}$}{Bessel unipotent subgroup, its part fixing the frame, the group fixing $e_a,e^*_a,z$, the Bessel subgroup}
\nt{$\ZT{T}{\gamma}$, $\Zr{y}$, $\lmin{l}{Z}$, $\fdisc{N'}{Z}$, $\gdr{Z}$, $\sd{u}$}{frame Gram matrices, leading minor, discriminants, sum above the diagonal}
\nt{$\Gone{\fs}$, $\Ginf$, $\Tone$}{the real group, its product over the places, the form over \m{\CC}}
\nt{$\bsize{y}$, $\InCut{y}{b_0}$, $\frbox{\rho}$, $\cellm{y}{\deg_0}{\rho}{d_0}$, $\occ{\fs}$}{Bessel size, the cut, the frame box, the cell mass, a constant}
\nt{$S$, $\Fill{E}{\fs}$}{the filling data; \m{\bb}, \m{\Nn}, \m{\qS{S}{i}}, \m{\gq}, \m{\piS{S}}, \m{\vpS{S}}, \m{\ChS{S}{i}}, \m{\KcolS{S}}, \m{\BsysS{S}{\tau}} its components}
\nt{$\Trat{S}{A}$, $\Tlat{S}{A}$, $\GS{S}{A}$, $\GFin{S}$}{the forms and groups of the data over a ring $A$}
\nt{$\KS{S}$, $\KgS{S}$, $\GRes{S}$, $\redS{S}$, $\Sbar$}{integral subgroup, principal level, the finite group, reduction, reduced Bessel image}
\nt{$\rinf{S}$, $\rfin{S}$, $\incH{S}$, $\bscale{S}{\tau}{i}$, $\kapS{S}{\tau}$, $\fscale{S}{\tau}{a}$}{maps to the real and finite adelic groups; scales}
\nt{$\Cl{S}$, $c,c',c''$, $x_c$, $\hc{S}{c}$, $[x]$}{classes, representatives, the class of $x$}
\nt{$\pos{S}{c}{c'}{\gamma}$, $\terms{S}{c}{s}{c'}$, $\clU{S}{c}$, $\clUo{S}{c}$, $\clC{S}{c}$, $\clCX{S}{c}$}{positions, terms, arithmetic unipotent groups, class coordinates}
\nt{$\Zf{S}{\gamma}$, $\tdisc{S}{\gamma}$, $\tdnorm{S}{\gamma}$}{frame products, discriminant and its norm}
\nt{$\Wt{\sigma}{S}{\tau}$, $\It{S}{\tau}$, $\degt{S}{\tau}$, $\Bc{S}$, $\bdeg{S}$, $\chi$}{Bessel functions, index sets, degrees, a character}
\nt{$\sigma$, $\full$, $\cut$}{a side}
\nt{$\besker{\sigma}{S}{\gamma}{y}$, $\cutoff{S}{c}$, $\theta_{S,c}$, $\omega_{S,c}$, $\IsBP{S}{c}{\omega}$}{kernel, cutoff and its two factors, box partitions}
\nt{$\amp{\sigma}{S}{c}{c'}{\gamma}$, $\ampabs{\sigma}{S}{c}{c'}{\gamma}$, $\ampsum{\sigma}{S}{\Xi}{c}{s}{c'}$}{amplitude, absolute amplitude, amplitude sum}
\nt{$\bxs{S}$, $\bxp{S}$, $\IsNr{S}{\gamma}$, $\IsNRep{S}{c}{c'}{\gamma}$}{box size, boxed points, near elements, near representatives}
\nt{$\HCFV{\gamma}{c_1}{c_2}{a}$, $\IsNF{j}{\gamma}$, $\HasLev{S}{j}{\gamma}$}{shared flag vectors, normal forms, levels}
\nt{$\link{\gamma}$, $\clink{\gamma}{\gamma'}$, $\Lambda$, $R$}{link, common link, a link, a set of subspaces}
\nt{$\uchar{S}{u}$, $\ChA{S}{\gamma}$, $\Agr{S}{s}$, $\IsStab{S}{\Xi}$, $\FixN{s}{X}$}{the unipotent character and the conditions built on it; stable sets; fixed vectors}
\nt{$\tl$, $\codimF{S}$, $\cdim{S}$, $\dd$, $\dm$, $\cw{S}{i}$, $\nuS$, $\fsprod{S}$, $\siegC{S}$, $\packC{S}$, $\mass{E}{\rr{\fs}}$}{the constants of the data}
\nt{$\energy{S}{\mathcal A}{\Omega}$, $\linkE{\sigma}{S}{\Xi}{\Omega}{R}$}{energy, link energy (Definition~\ref{d:energy}); $\mathcal A,\mathcal A'$ families, $V$ a normed group, $\Omega$ a subset of \m{\GRes{S}}}
\nt{$\csize{S}{j}{x}$, $\mathsf D_j$}{class size (Definition~\ref{d:cs})}
\nt{$\bcut{S}{c}$, $\schw{S}{c}{\gamma}{y}$, $\schi{S}{c}{c'}{\gamma'}$, $\Delta$}{boxed cutoff, Schur weight, Schur integral (Definition~\ref{d:schur}), frame volume}
\nt{$\mathsf k,\mathsf k',\mathsf u$, $\Zl{\mathsf k}$, $\xrows{Z}$, $\IdxB{j}$, $\mathcal R(\mathsf k)$, $\mathcal R_j(\mathsf k)$}{elements of \m{\GFin{S}}; frame Gram matrix, its first rows, row indices below $j$, row lattices}
\nt{$\xri{S}{j}{\mathsf k}$, $\sati{S}{\mathsf k}$, $\mnorm{S}{j}{\gamma}$, $\cgauge{S}{j}{c}{c'}{\gamma}$}{index of rows, saturation index, minor norm, class gauge}
\nt{$\latd{S}{c}(U)$, $\sld{H}{U}$, $\slb{H}{U}$, $M_U$, $\Lambda_U(\mathsf k)$, $\fvi{S}{c}{\mathsf k}{a}{U}$}{lattice discriminants, basis matrices, frame value index; $U$, $U'$ subspaces of \m{E^\kappa}}
\nt{$\pinv{S}{c}{c'}{\gamma}{\gamma'}$, $\telem{S}{c}{p}$, $\ipw{S}{c}{\eta}{p}$, $\IsLab{S}{j}{c}{\mathsf t}{p}$, $\Cong{S}{c}{c'}{\gamma}{\gamma'}$}{pair invariant, term element, inverse power weight, labelled terms, congruent terms}
\nt{$\termw{S}{j}{i}{c}{p}{p'}{U}$, $\pairw{S}{j}{c}{c'}{\gamma}{\gamma'}$, $\eta_{j,U}$, $\gid{S}$, $\padl{j}{Z}$, $\fmid{\mathsf k}$, $\mbc{\fs}{\piS{S}}$}{term and pair weights, the idele of \m{\gq}, padded matrix, middle rows of the frame image, a matrix of Note~S}
\nt{$\gamma,\gamma',\delta$, $\gamma_0,\gamma_1,\gamma_p$}{rational elements of \m{\GS{S}{E}}; in \inp{R2}, \inp{R3} elements of \m{\Gone{\fs}}}
\nt{$u,u',u_0,u_1,u_2,u_p,r$}{unipotent elements, rational or real; $x$ also an element of \m{\UE} in Theorem~B}
\nt{$y,y',y'',y^0$}{points of \m{\Ginf} or of \m{\Gone{\fs}}}
\nt{$s,s'$, $\mathsf t$}{elements of \m{\GRes{S}}; a label in \m{\Sbar}}
\nt{$j$, $j_0$, $j_1$, $l$}{levels; \m{j_0=\min(\sro,j)}, \m{j_1=j-j_0}; the bound of Theorem~D\,(b), also an index of minors}
\nt{$a,a',a''$, $i$, $e$, $m$, $m'$, $n_1$, $l_1$, $m_1$}{indices of the frame and of rows; natural numbers: codimensions, exponents, sizes}
\nt{$t,t',t_0,t_1,t_2,t_3,\tilde t,\tilde t_2$, $v,v',v_0,w,w'$}{vectors of the hermitian spaces; coefficient vectors in \m{E^{\sn}}; $w$ also the real bound of Theorems D, E}
\nt{$C$, $X$, $X'$, $w$, $\varepsilon$, $x$}{the constants of Setting~\ref{set:en} and of Theorems D, E; real numbers}
\nt{$Q$, $Q_0$, $\mu$, $\phi$, $F$, $C_0$, $q$}{the measure space, subset, map, function and constant of Proposition~\ref{p:avg}; a point}
\nt{$\eta^0$, $\mu^0$, $\Uoone^\infty$, $W$}{measures of \inp{B8}; the product group; the product of squared norms}
\nt{$\mathcal D$, $\mathcal B$, $\mathcal P$, $\mathcal P'$, $\mathcal P(q,U)$, $p,p',p_1,q$}{a double coset, a colour class, index sets of terms and pairs; pairs}
\nt{$Z$, $Z^\tau$, $Z^\tau_j$, $Z_\Sigma$, $Y$, $\xi$, $\ell_\delta$, $\ell$, $L_\delta$, $L$, $R_\tau$, $P_\tau$, $\Pi_0$, $D_2$}{matrices and numbers in the proof of Theorem~A}
\nt{$\mathcal K$, $\mathcal K_\tau$, $\theta_0$, $\Sigma$, $R_0$, $V_0$, $A$, $\alpha_1$, $u(\alpha_1)$}{the volume kernel and its data \inp{R2}; the parametrization \inp{L9}}
\nt{$\Phi_\delta$, $\varphi$, $\varphi_p$, $\alpha$, $\theta$, $\theta_x$, $\Gamma$, $\Gamma_x$, $\Gamma_0$, $x_0$}{weights; cutoffs; discrete groups; a coset}
\nt{$\mathcal C$, $\mathcal C_\tau$, $n_\tau$, $m_\tau$, $Q_\tau$, $t_\tau$, $n_1$, $m_1$, $Q_1$, $t_1$, $\rho$}{the data of the open-cell bound \inp{R3}}
\nt{$\Phi$, $B_0$, $\mathfrak A$, $\tilde a$, $K'$, $y_a$, $r_a$, $\mathsf k_a$, $\gamma_a$, $i_a$, $F_0$, $D_1$, $x_0$, $w_0$, $H$}{Proposition~\ref{p:latdisc}; in \inp{S15} a matrix with a left inverse}
\nt{$G'$, $G''$, $\mathfrak M$}{the matrices of Lemma~\ref{l:short}}
\nt{$d_0$, $m_0$, $N_1$, $C_1$, $C_2$, $\Upsilon$, $H$, $\Theta$, $\Psi$, $\alpha$, $\ell_p$, $U_1$, $x_1$, $J$, $\vartheta_1$}{Proposition~\ref{p:count} and \inp{B10}; $J$ a split Gram matrix}
\nt{$P_0$, $p_0$, $a_0$, $r_0$, $\mu_0$, $P_1$, $M_1$, $M_2$, $D_2$, $K_0$, $K_1$, $I_1,I_2,I_3$, $Y_p$, $M_0$, $\eta_p$}{Proposition~\ref{p:pairs}, Lemma~\ref{l:gindex}, the inputs \inp{L11}, \inp{R5}; finite index sets}
\nt{$\lambda$, $c_0$, $c_1$, $c_2$, $\varkappa$, $\epsilon$, $\epsilon_1,\epsilon_2,\epsilon_3$, $L_0$, $R_0$, $G$, $G'$, $d_1$, $o$, $o'$, $\mathfrak d$, $\mathfrak d'$, $i_0$, $\mathfrak S$, $A_\Lambda$, $\Xi_\Lambda$, $\Xi_1,\Xi_2,\Xi_3$}{a linear form; numbers and sets in the proofs of Lemma~\ref{l:cell} and of Theorems D, E}
\nt{$\beta_1$, $\beta_2$, $\Sigma_1$, $B_1$}{auxiliary numbers in the proofs of Theorem~A, Proposition~\ref{p:pairs} and Theorem~E}
\nt{$\pi$}{the number \m{3.14\dots}; \m{\piS{S}} is a component of the data}
\nt{$F_1$, $F_2$, $b_0$, $f$, $g$, $k$, $n$, $N$, $N'$, $O$, $K$, $M$, $M'$, $I$, $B$}{functions in Lemma~\ref{l:cs}; a real number; letters that occur only inside the notations above: \m{\ninf{E}}, \m{\gS{S}}, \m{\kk{\fs}}, \m{\nn{\fs}}, \m{\NS{S}}, \m{\hO}, \m{\KS{S}}, matrices $M$, \m{\Ix{\fs}}, \m{\bS{S}}}
\end{multicols}
\endgroup

\end{document}
